% Production d’écrits utilitaires : le rapport (élaboration du plan détaillé) — Français 1re Tech — Cours signé OnBuch+
% Programme officiel MINESEC. Compiler avec : tectonic N14-production-ecrits-utilitaires-rapport-elaboration-plan.tex
\newcommand{\DOCMATIERE}{Français}
\newcommand{\DOCNIVEAU}{1re Tech}
\newcommand{\DOCMODULE}{Français – classe de Première technique}
\newcommand{\DOCLECON}{N14}
\newcommand{\DOCTITRE}{Production d’écrits utilitaires : le rapport (élaboration du plan détaillé)}
% OnBuch+ — préambule « cours signés » ENRICHI (compilé avec tectonic / XeLaTeX)
% Système visuel « Pop » d'OnBuch+ : fond crème (jamais blanc), bordures encre
% épaisses, ombres « dures » décalées, titres Archivo Black en capitales.
% Version MATHÉMATIQUES : graphes (pgfplots/tikz), boîtes pédagogiques
% (définition, propriété, méthode, attention, le savais-tu, exercice résolu,
% exercice, TP, bilan), page de garde et signature OnBuch+ ancrée en bas.
%
% Macros attendues dans main.tex AVANT \input{preamble} :
%   \DOCMATIERE  \DOCNIVEAU  \DOCMODULE  \DOCLECON (numéro)  \DOCTITRE
\documentclass[11pt]{article}

\usepackage{fontspec}
\usepackage[french]{babel}
\usepackage[a4paper,margin=2cm,top=3.4cm,bottom=2.8cm,headheight=1.4cm,footskip=1.3cm]{geometry}
\usepackage{amsmath,amssymb,amsfonts}
\usepackage{mathtools} % \xmapsto, \coloneqq... souvent utilisés par les modèles
\usepackage{cancel} % simplifications barrées (bilans de forces)
% \abs{z} (module, valeur absolue) et \norm{u} (norme), utilisés par les modèles
\providecommand{\abs}{}\let\abs\relax\DeclarePairedDelimiter{\abs}{\lvert}{\rvert}
\providecommand{\norm}{}\let\norm\relax\DeclarePairedDelimiter{\norm}{\lVert}{\rVert}
\usepackage[table]{xcolor} % [table] : fournit aussi \rowcolors (alternance de lignes)
\usepackage{pagecolor}
\usepackage{tikz}
\usetikzlibrary{arrows.meta,positioning,calc,shapes.geometric,shapes.misc,shapes.arrows,decorations.pathmorphing,decorations.pathreplacing,decorations.markings,patterns,shadows,mindmap,trees,fit,backgrounds,matrix,intersections,angles,quotes}
\usepackage{pgfplots}
\usepackage[version=4]{mhchem} % équations nucléaires \ce{^{235}_{92}U}
\usepackage[european,siunitx]{circuitikz} % schémas électriques
\pgfplotsset{compat=1.18}
\usepgfplotslibrary{fillbetween,groupplots} % aires entre courbes (calcul intégral)
\usepackage{enumitem}
\usepackage{fancyhdr}
% [most] charge toutes les bibliothèques tcolorbox (listings, minted, raster,
% documentation...) et gonfle inutilement la mémoire TeX sur un document long
% (~300 boîtes) : on ne charge que ce qui est réellement utilisé.
\usepackage[skins,breakable]{tcolorbox}
\usepackage{graphicx}
\usepackage{colortbl}
\usepackage{booktabs}
\usepackage{tabularx}
\usepackage{multirow}
\usepackage{diagbox} % case d'en-tête diagonale des tableaux à double entrée
% Colonnes extensibles centrée / gauche / droite (C, L, R), souvent utilisées par les modèles
\newcolumntype{C}{>{\centering\arraybackslash}X}
\newcolumntype{L}{>{\raggedright\arraybackslash}X}
\newcolumntype{R}{>{\raggedleft\arraybackslash}X}
\usepackage{array}
\usepackage{multicol}
\usepackage{siunitx}
% parse-numbers=false : accepte aussi \SI{2,5 \times 10^{-3}}{...}
\sisetup{locale=FR,output-decimal-marker={,},group-minimum-digits=4,parse-numbers=false}
\DeclareSIUnit\ohm{\ensuremath{\symup{\Omega}}} % U+2126 absent de la police
\DeclareSIUnit\division{div} % réglages d'oscilloscope (ms/div, V/div)
\DeclareSIUnit\tr{tr} % tours (vitesse de rotation, ex. \SI{1500}{\tr\per\minute})
\DeclareSIUnit\molar{M} % concentration molaire (ex. \SI{3}{\molar}, \SI{1}{\micro\molar})
\DeclareSIUnit\an{an} % année (le modèle confond parfois avec \year, un compteur TeX) — ex. \SI{5}{\kilo\meter\per\an}
\DeclareSIUnit\FCFA{FCFA} % franc CFA (ex. \SI{500}{\FCFA\per\kilo\gram})
\DeclareSIUnit\degreeBrix{\textdegree Brix} % teneur en sucre d'un sirop/jus (réfractométrie)
\DeclareSIUnit\month{mois} % le modèle utilise parfois \month, un compteur TeX, comme unité de durée
\DeclareSIUnit\microliter{\micro\liter} % le modèle écrit parfois \microliter en un seul token
\DeclareSIUnit\million{millions} % dénombrement (ex. \SI{15}{\million\per\mL} spermatozoïdes)
\usepackage{qrcode}
% \degree : commande fréquemment utilisée par les modèles pour un angle ou une
% température en dehors de \SI{}{} (non fournie par les paquets chargés).
\providecommand{\degree}{^{\circ}}
% \qed (fin de démonstration), utilisé par les modèles sans amsthm
\providecommand{\qed}{\hfill$\square$}
% \barre : double barre des valeurs interdites dans un tableau de variations (tkz-tab)
\providecommand{\barre}{\ensuremath{\|}}
% environnement proof (amsthm non chargé) utilisé par les modèles
\ifdefined\proof\else\newenvironment{proof}[1][Démonstration]{\par\noindent\textit{#1.}\ }{\qed\par}\fi
\usepackage{lastpage}
\usepackage{hyperref}
\hypersetup{hidelinks}

% ============================================================
% Palette « Pop » OnBuch+ — mêmes hex que la classe P (plus_ui.dart)
% ============================================================
\definecolor{poporange}{HTML}{F59321}
\definecolor{popdark}{HTML}{E07A0C}
\definecolor{popcream}{HTML}{FFF6E8}
\definecolor{popink}{HTML}{1C1714}
\definecolor{poppurple}{HTML}{7A5AE0}
\definecolor{popgreen}{HTML}{2E9E6A}
\definecolor{poppink}{HTML}{E0457B}
\definecolor{popblue}{HTML}{2F7DE1}
\definecolor{popgold}{HTML}{C98A12}
\definecolor{popmuted}{HTML}{8A7F76}
\definecolor{popline}{HTML}{EADFD2}
\definecolor{popgray}{HTML}{8A7F76}
\colorlet{popgrey}{popgray}
% Alias tolérants (le modèle nomme parfois les couleurs différemment)
\colorlet{popwhite}{white}
\colorlet{popblack}{popink}
\colorlet{popred}{poppink}
\colorlet{popyellow}{popgold}
% Teintes claires (fonds de tableaux/encadrés) — le modèle nomme parfois
% les couleurs avec un suffixe L (ex. popblueL) sans les avoir définies.
\colorlet{poporangeL}{poporange!20}
\colorlet{popdarkL}{popdark!20}
\colorlet{poppurpleL}{poppurple!20}
\colorlet{popgreenL}{popgreen!20}
\colorlet{poppinkL}{poppink!20}
\colorlet{popblueL}{popblue!20}
\colorlet{popgoldL}{popgold!20}
\colorlet{popredL}{popred!20}
\colorlet{popgrayL}{popgray!20}
% Teintes claires pour fonds de tableaux / zones
\colorlet{popblueL}{popblue!10!white}
\colorlet{poporangeL}{poporange!14!white}
\colorlet{popgreenL}{popgreen!12!white}
\colorlet{poppurpleL}{poppurple!12!white}
\colorlet{poppinkL}{poppink!12!white}
\colorlet{popgoldL}{popgold!14!white}

\pagecolor{popcream}

% ============================================================
% Typographie
% ============================================================
\defaultfontfeatures{Ligatures=TeX}
\setmainfont{PlusJakartaSans-Medium.ttf}[Path=../../../fonts/,BoldFont=PlusJakartaSans-Bold.ttf]
% Maths en sans-serif (Fira Math) avec chiffres et lettres droites de Plus
% Jakarta, pour que les formules se fondent dans le texte.
\usepackage[math-style=ISO,bold-style=ISO]{unicode-math}
\setmathfont{FiraMath-Regular.otf}[Path=../../../fonts/]
\setmathfont{PlusJakartaSans-Medium.ttf}[Path=../../../fonts/,range={up/num,up/{Latin,latin},\mathup}]
\usepackage{icomma} % virgule décimale sans espace en mode math
% Caractères mathématiques Unicode absents des polices de texte (Plus Jakarta,
% Archivo Black) : rendus par leur équivalent mathématique, en texte comme en
% formule — sinon ils s'affichent en carré vide (ex. « Arithmétique dans ℤ »).
\usepackage{newunicodechar}
\newunicodechar{ℝ}{\ensuremath{\mathbb{R}}}
\newunicodechar{ℤ}{\ensuremath{\mathbb{Z}}}
\newunicodechar{ℂ}{\ensuremath{\mathbb{C}}}
\newunicodechar{ℕ}{\ensuremath{\mathbb{N}}}
\newunicodechar{ℚ}{\ensuremath{\mathbb{Q}}}
\newunicodechar{𝔻}{\ensuremath{\mathbb{D}}}
\newunicodechar{α}{\ensuremath{\alpha}}
\newunicodechar{β}{\ensuremath{\beta}}
\newunicodechar{γ}{\ensuremath{\gamma}}
\newunicodechar{δ}{\ensuremath{\delta}}
\newunicodechar{ε}{\ensuremath{\varepsilon}}
\newunicodechar{θ}{\ensuremath{\theta}}
\newunicodechar{λ}{\ensuremath{\lambda}}
\newunicodechar{μ}{\ensuremath{\mu}}
\newunicodechar{σ}{\ensuremath{\sigma}}
\newunicodechar{τ}{\ensuremath{\tau}}
\newunicodechar{φ}{\ensuremath{\varphi}}
\newunicodechar{ψ}{\ensuremath{\psi}}
\newunicodechar{ω}{\ensuremath{\omega}}
\newunicodechar{Σ}{\ensuremath{\Sigma}}
% majuscules produites par \MakeUppercase dans les titres (α-aminés -> Α)
\newunicodechar{Α}{\ensuremath{\alpha}}
\newunicodechar{Β}{\ensuremath{\beta}}
\newunicodechar{Γ}{\ensuremath{\Gamma}}
\newunicodechar{Δ}{\ensuremath{\Delta}}
\newunicodechar{Ε}{\ensuremath{\varepsilon}}
\newunicodechar{Θ}{\ensuremath{\Theta}}
\newunicodechar{Λ}{\ensuremath{\Lambda}}
\newunicodechar{Μ}{\ensuremath{\mu}}
\newunicodechar{Τ}{\ensuremath{\tau}}
\newunicodechar{Φ}{\ensuremath{\Phi}}
\newunicodechar{Ψ}{\ensuremath{\Psi}}
\newunicodechar{Ω}{\ensuremath{\Omega}}
\newunicodechar{↦}{\ensuremath{\mapsto}}
\newunicodechar{∈}{\ensuremath{\in}}
\newunicodechar{∉}{\ensuremath{\notin}}
\newunicodechar{∧}{\ensuremath{\wedge}}
\newunicodechar{⇔}{\ensuremath{\Leftrightarrow}}
\newunicodechar{⟺}{\ensuremath{\Longleftrightarrow}}
\newunicodechar{⇒}{\ensuremath{\Rightarrow}}
\newunicodechar{⟹}{\ensuremath{\Longrightarrow}}
\newunicodechar{∘}{\ensuremath{\circ}}
\newunicodechar{∪}{\ensuremath{\cup}}
\newunicodechar{∩}{\ensuremath{\cap}}
\newunicodechar{≡}{\ensuremath{\equiv}}
\newunicodechar{⊂}{\ensuremath{\subset}}
\newunicodechar{⊥}{\ensuremath{\perp}}
\newunicodechar{∃}{\ensuremath{\exists}}
\newunicodechar{∀}{\ensuremath{\forall}}
\newunicodechar{∛}{\ensuremath{\sqrt[3]{\,}}}
\newunicodechar{∜}{\ensuremath{\sqrt[4]{\,}}}
\newunicodechar{⃗}{\ensuremath{{}^{\to}}}
\newunicodechar{ⁿ}{\ifmmode^{n}\else\textsuperscript{\ensuremath{n}}\fi}
\newunicodechar{ˣ}{\ifmmode^{x}\else\textsuperscript{\ensuremath{x}}\fi}
\newunicodechar{ᵇ}{\ifmmode^{b}\else\textsuperscript{\ensuremath{b}}\fi}
\newunicodechar{ʳ}{\ifmmode^{r}\else\textsuperscript{\ensuremath{r}}\fi}
\newunicodechar{ᵘ}{\ifmmode^{u}\else\textsuperscript{\ensuremath{u}}\fi}
\newunicodechar{ʸ}{\ifmmode^{y}\else\textsuperscript{\ensuremath{y}}\fi}
\newunicodechar{ᵃ}{\ifmmode^{a}\else\textsuperscript{\ensuremath{a}}\fi}
\newunicodechar{ᵉ}{\ifmmode^{e}\else\textsuperscript{\ensuremath{e}}\fi}
\newunicodechar{ᵏ}{\ifmmode^{k}\else\textsuperscript{\ensuremath{k}}\fi}
\newunicodechar{ᵖ}{\ifmmode^{p}\else\textsuperscript{\ensuremath{p}}\fi}
\newunicodechar{ᴺ}{\ifmmode^{N}\else\textsuperscript{\ensuremath{N}}\fi}
\newunicodechar{ᶜ}{\ifmmode^{c}\else\textsuperscript{\ensuremath{c}}\fi}
\newunicodechar{ʰ}{\ifmmode^{h}\else\textsuperscript{\ensuremath{h}}\fi}
\newunicodechar{ᵠ}{\ifmmode^{q}\else\textsuperscript{\ensuremath{q}}\fi}
\newunicodechar{ₙ}{\ifmmode_{n}\else\textsubscript{\ensuremath{n}}\fi}
\newunicodechar{ₐ}{\ifmmode_{a}\else\textsubscript{\ensuremath{a}}\fi}
\newunicodechar{ᵢ}{\ifmmode_{i}\else\textsubscript{\ensuremath{i}}\fi}
\newunicodechar{ᵣ}{\ifmmode_{r}\else\textsubscript{\ensuremath{r}}\fi}
\newunicodechar{ₖ}{\ifmmode_{k}\else\textsubscript{\ensuremath{k}}\fi}
\newunicodechar{ᵦ}{\ifmmode_{\beta}\else\textsubscript{\ensuremath{\beta}}\fi}
\newunicodechar{ₓ}{\ifmmode_{x}\else\textsubscript{\ensuremath{x}}\fi}
\newfontfamily\popdisplay{ArchivoBlack-Regular.ttf}[Path=../../../fonts/]
\newfontfamily\poplogo{SpaceGrotesk-Bold.ttf}[Path=../../../fonts/]
\newfontfamily\popsemi{PlusJakartaSans-SemiBold.ttf}[Path=../../../fonts/]
\newfontfamily\popxbold{PlusJakartaSans-ExtraBold.ttf}[Path=../../../fonts/]
\newfontfamily\popxboldls{PlusJakartaSans-ExtraBold.ttf}[Path=../../../fonts/,LetterSpace=3.0]

\setlist[enumerate]{leftmargin=1.7em,itemsep=5pt,topsep=4pt}
\setlist[itemize]{leftmargin=1.7em,itemsep=4pt,topsep=4pt,label={\color{poporange}\textbullet}}
\setlength{\parindent}{0pt}
\setlength{\parskip}{5pt}
\renewcommand{\arraystretch}{1.25}

% pgfplots : style Pop par défaut
\pgfplotsset{
  every axis/.append style={axis line style={popink,line width=1pt},tick style={popink},
    grid style={popline,line width=0.5pt},label style={font=\small},tick label style={font=\footnotesize}},
  popcurve/.style={poporange,line width=1.8pt,no markers,smooth},
}
% Même style disponible dans un \draw TikZ simple (hors pgfplots)
\tikzset{popcurve/.style={poporange,line width=1.8pt}}
\tikzset{popgrid/.style={popline,very thin}}
\colorlet{popcurve}{poporange} % le nom du style est parfois utilisé comme couleur

% ============================================================
% Pill / squiggle
% ============================================================
\newcommand{\poppill}[3]{%
  \tikz[baseline=-0.55ex]\node[draw=popink,line width=1.2pt,rounded corners=2.6pt,%
    fill=#2,inner xsep=6.5pt,inner ysep=3pt]{%
    \popxboldls\fontsize{7}{7}\selectfont\color{#3}\MakeUppercase{#1}};%
}
\newcommand{\popsquiggle}[1]{%
  \tikz[baseline=-0.4ex]{
    \draw[#1,line width=2.6pt,line cap=round]
      plot [smooth] coordinates {(0,0.09) (0.34,0.01) (0.68,0.21) (1.02,0.01) (1.36,0.10)};
  }%
}
% Mot-clé surligné (marqueur orange sous le texte)
\newcommand{\cle}[1]{{\popxbold\color{popdark}#1}}

% ============================================================
% En-tête / pied
% ============================================================
\pagestyle{fancy}
\fancyhf{}
\renewcommand{\headrulewidth}{0pt}
\renewcommand{\footrulewidth}{0pt}
\lhead{\raisebox{-0.15ex}{\poplogo\fontsize{13}{13}\selectfont\color{popink}OnBuch\color{poporange}+}}
\rhead{\poppill{\DOCMATIERE\ \textbullet\ \DOCNIVEAU\ \textbullet\ Leçon \DOCLECON}{white}{popink}}
\lfoot{\popxbold\fontsize{7.6}{7.6}\selectfont\color{popmuted}\MakeUppercase{Cours signé OnBuch+}\ \textbullet\ {\popsemi\DOCTITRE}}
\rfoot{\tikz[baseline=-0.6ex]\node[rounded corners=8.5pt,fill=popink,minimum height=17pt,minimum width=17pt,inner xsep=5pt,inner sep=0]{\popxbold\fontsize{8}{8}\selectfont\color{popcream}\thepage\,/\,\pageref*{LastPage}};}

% ============================================================
% Titres
% ============================================================
\newcommand{\chaptitre}[2]{%
  \begin{tcolorbox}[enhanced,colback=poporange,colframe=popink,boxrule=2.6pt,arc=11pt,
      drop shadow=popink,left=15pt,right=15pt,top=12pt,bottom=11pt,before skip=3pt,after skip=18pt]
    \poppill{#2}{popink}{popcream}\par\vspace{9pt}
    {\raggedright\hyphenpenalty=10000\exhyphenpenalty=10000\popdisplay\fontsize{22}{24}\selectfont\color{popcream}\MakeUppercase{#1}\par}
  \end{tcolorbox}
}
% Section numérotée : pastille orange avec le numéro + titre Archivo
\newcounter{popsec}
\newcommand{\coursec}[1]{%
  \stepcounter{popsec}\setcounter{popsub}{0}%
  \par\vspace{16pt}\noindent
  \begin{minipage}{\linewidth}
  \tikz[baseline=-0.7ex]\node[draw=popink,line width=1.6pt,fill=poporange,rounded corners=4pt,minimum size=22pt,inner sep=0]{\popdisplay\fontsize{12}{12}\selectfont\color{popcream}\thepopsec};%
  \hspace{8pt}{\popdisplay\fontsize{15}{17}\selectfont\color{popink}\MakeUppercase{#1}}\par
  \vspace{4pt}\noindent{\color{poporange}\rule{36pt}{3.6pt}}
  \end{minipage}\par\nopagebreak\vspace{8pt}
}
\newcounter{popsub}
\newcommand{\courssub}[1]{%
  \stepcounter{popsub}%
  \par\vspace{9pt}\noindent{\popxbold\fontsize{11.5}{13}\selectfont\color{popink}{\color{poporange}\thepopsec.\thepopsub}\hspace{6pt}#1}\par\nopagebreak\vspace{4pt}
}

% ============================================================
% Boîtes pédagogiques — même recette : fond blanc, bordure épaisse colorée,
% ombre dure encre, badge plein. Titre optionnel [#1].
% ============================================================
\tcbset{popbase/.style={breakable,enhanced,colback=white,boxrule=2.3pt,arc=9pt,
  shadow={3pt}{-3pt}{0pt}{popink},left=14pt,right=14pt,top=11pt,bottom=11pt,before skip=11pt,after skip=15pt,
  fontupper=\popsemi\fontsize{10.6}{14.5}\selectfont\color{popink},
  fontlower=\popsemi\fontsize{10.6}{14.5}\selectfont\color{popink}}}
% Coche dessinée (le glyphe ✓ n'existe pas dans les polices du document)
\newcommand{\popcheck}[1]{\tikz[baseline=-0.5ex]\draw[#1,line width=1.6pt,line cap=round,line join=round](-0.13,0.0)--(-0.04,-0.09)--(0.13,0.1);}
\providecommand{\coche}{\popcheck{popgreen}}
\providecommand{\resultat}[1]{\par\smallskip\noindent\fcolorbox{popgreen}{popgreenL}{\parbox{\dimexpr\linewidth-2\fboxsep-2\fboxrule\relax}{\textbf{Résultat.} #1}}\par\smallskip}
\newcommand{\popboxhead}[3]{\poppill{#1}{#2}{white}\ifx\relax#3\relax\else\hspace{6pt}{\popxbold\fontsize{10.6}{12}\selectfont\color{popink}#3}\fi\par\vspace{7pt}}

\newtcolorbox{aretenir}[1][]{popbase,colframe=popgreen,before upper={\popboxhead{À retenir}{popgreen}{#1}}}
\newtcolorbox{exemplebox}[1][]{popbase,colframe=poporange,before upper={\popboxhead{Exemple}{poporange}{#1}}}
\newtcolorbox{definition}[1][]{popbase,colframe=popblue,before upper={\popboxhead{Définition}{popblue}{#1}}}
\newtcolorbox{propriete}[1][]{popbase,colframe=poppurple,before upper={\popboxhead{Propriété}{poppurple}{#1}}}
\newtcolorbox{methode}[1][]{popbase,colframe=popgold,before upper={\popboxhead{Méthode}{popgold}{#1}}}
\newtcolorbox{attention}[1][]{popbase,colframe=poppink,before upper={\popboxhead{Attention}{poppink}{#1}}}
\newtcolorbox{experience}[1][]{popbase,colframe=popblue,colback=popblueL,before upper={\popboxhead{Expérience}{popblue}{#1}}}
% « Le savais-tu ? » : carte violette pleine (contexte, culture, Cameroun)
\newtcolorbox{savaistu}[1][]{popbase,colback=poppurple,colframe=popink,
  fontupper=\popsemi\fontsize{10.4}{14.2}\selectfont\color{white},
  before upper={\poppill{Le savais-tu\,?}{white}{poppurple}\ifx\relax#1\relax\else\hspace{6pt}{\popxbold\color{white}#1}\fi\par\vspace{7pt}}}
% Exercice résolu : énoncé en haut, \tcblower puis la solution sur fond crème
\newcounter{popexo}\newcounter{popexr}
\newtcolorbox{exoresolu}[1][]{popbase,colframe=poporange,colbacklower=popcream,
  segmentation style={popink,line width=1.2pt,dashed},
  before upper={\stepcounter{popexr}\popboxhead{Exercice résolu \thepopexr}{poporange}{#1}},
  before lower={\poppill{Solution}{popink}{popcream}\par\vspace{6pt}}}
% Exercice (non résolu en place) — la correction est dans la partie Corrigés
\newtcolorbox{exercice}[1][]{popbase,colframe=popink,
  before upper={\stepcounter{popexo}\popboxhead{Exercice \thepopexo}{popink}{#1}}}
% Corrigé
\newtcolorbox{corrige}[1][]{popbase,colframe=popgreen,colback=popgreenL,
  before upper={\popboxhead{Corrigé}{popgreen}{#1}}}
% Objectifs / prérequis / bilan
\newtcolorbox{objectifs}{popbase,colframe=popink,colback=poporangeL,
  before upper={\popboxhead{Objectifs de la leçon}{popink}{}}}
\newtcolorbox{prerequis}{popbase,colframe=popink,colback=popblueL,
  before upper={\popboxhead{Prérequis}{popblue}{}}}
\newtcolorbox{bilan}{popbase,colframe=popink,colback=white,boxrule=2.8pt,
  before upper={\popboxhead{Fiche bilan}{poporange}{L'essentiel à connaître pour l'examen}}}
% Figure encadrée avec légende
\newtcolorbox{popfigure}[1][]{enhanced,breakable=false,colback=white,colframe=popline,boxrule=1.4pt,arc=8pt,
  left=10pt,right=10pt,top=10pt,bottom=8pt,before skip=10pt,after skip=12pt,halign=center,
  lower separated=false,halign lower=center,
  fontlower=\popsemi\fontsize{9}{11.5}\selectfont\color{popmuted},
  #1}
\newcommand{\legende}[1]{\par\vspace{6pt}{\popsemi\fontsize{9}{11.5}\selectfont\color{popmuted}#1}}

% ============================================================
% Signature OnBuch+
% ============================================================
\newcommand{\poplogoinline}[1]{{\poplogo\fontsize{#1}{#1}\selectfont\color{popink}OnBuch\color{poporange}+}}
\newcommand{\signatureonbuch}{%
  \begin{tcolorbox}[enhanced,colback=popink,colframe=popink,boxrule=2.2pt,arc=10pt,
      drop shadow=poporange,left=14pt,right=14pt,top=10pt,bottom=10pt,before skip=0pt,after skip=0pt]
    \tikz[baseline=-0.7ex]\node[circle,fill=popgreen,draw=popcream,line width=1.3pt,minimum size=22pt,inner sep=0]{\popcheck{white}};%
    \hspace{9pt}\begin{minipage}[c]{0.62\linewidth}
      {\popxbold\fontsize{12}{13}\selectfont\color{popcream}Signé\ \textbullet\ {\poplogo OnBuch\color{poporange}+}}\par\vspace{2pt}
      {\raggedright\popsemi\fontsize{8}{10.5}\selectfont\color{popline}\MakeUppercase{Équipe pédagogique OnBuch+}\\\MakeUppercase{Conforme au programme officiel MINESEC}\par}
    \end{minipage}\hfill
    {\poplogo\fontsize{22}{22}\selectfont\color{popcream}OnBuch\color{poporange}+}
  \end{tcolorbox}%
}

% ============================================================
% Page de garde
% #1 titre · #2 matière · #3 niveau · #4 module · #5 n° leçon · #6 durée ·
% #7 description / objectifs (texte ou itemize)
% ============================================================
\newcommand{\pagegarde}[7]{%
  \thispagestyle{empty}
  \newgeometry{a4paper,margin=1.7cm,top=1.6cm,bottom=1.4cm}%
  \begin{tikzpicture}[remember picture,overlay]
    \fill[poporange!18] ([xshift=-3.2cm,yshift=-3.6cm]current page.north east) circle (4.6cm);
    \fill[poppurple!14] ([xshift=2.4cm,yshift=7.6cm]current page.south west) circle (3.8cm);
  \end{tikzpicture}%
  {\poplogo\fontsize{30}{30}\selectfont\color{popink}OnBuch\color{poporange}+}\hfill
  \raisebox{0.6ex}{\poppill{Cours signé \textbullet\ Premium}{popink}{popcream}}\par
  \vspace{1pt}\popsquiggle{poppurple}\par
  \vspace{8pt}
  \begin{tcolorbox}[enhanced,colback=poporange,colframe=popink,boxrule=2.8pt,arc=13pt,
      drop shadow=popink,left=18pt,right=18pt,top=12pt,bottom=13pt,after skip=10pt]
    \poppill{#2 \textbullet\ #3}{popink}{popcream}\hspace{5pt}\poppill{#4}{white}{popink}\par\vspace{8pt}
    {\popdisplay\fontsize{14}{14}\selectfont\color{popink}LEÇON #5}\par\vspace{4pt}
    {\raggedright\hyphenpenalty=10000\exhyphenpenalty=10000\popdisplay\fontsize{25}{27}\selectfont\color{popcream}\MakeUppercase{#1}\par}
    \vspace{8pt}
    {\popxbold\fontsize{9.5}{9.5}\selectfont\color{popink}\MakeUppercase{Durée conseillée : #6}}
  \end{tcolorbox}
  \begin{tcolorbox}[enhanced,colback=white,colframe=popink,boxrule=2.2pt,arc=10pt,
      drop shadow=popink,left=16pt,right=16pt,top=10pt,bottom=10pt,after skip=8pt]
    \poppill{Dans cette leçon}{poppurple}{white}\par\vspace{5pt}
    \popsemi\fontsize{9.3}{12.6}\selectfont\color{popink}#7
  \end{tcolorbox}
  \noindent\begin{minipage}[t]{0.487\linewidth}
    \begin{tcolorbox}[enhanced,colback=popgreen,colframe=popink,boxrule=2.2pt,arc=10pt,
        drop shadow=popink,left=11pt,right=11pt,top=10pt,bottom=10pt,height=3.1cm,valign=center]
      \tcbox[colback=white,colframe=popink,boxrule=1.3pt,arc=5pt,boxsep=0pt,nobeforeafter,
        left=3pt,right=3pt,top=3pt,bottom=3pt,valign=center]{\qrcode[height=1.9cm]{https://play.google.com/store/apps/details?id=com.luvvix.onbuch&hl=fr}}%
      \hspace{8pt}\begin{minipage}[c]{0.5\linewidth}
        {\popdisplay\fontsize{11}{12.5}\selectfont\color{white}\MakeUppercase{Télécharge OnBuch}}\par\vspace{4pt}
        {\raggedright\popsemi\fontsize{8.4}{11}\selectfont\color{white}Gratuit \textbullet\ Google Play\\Tuteur IA, annales, concours, résultats\par}
      \end{minipage}
    \end{tcolorbox}
  \end{minipage}\hfill
  \begin{minipage}[t]{0.487\linewidth}
    \begin{tcolorbox}[enhanced,colback=poppurple,colframe=popink,boxrule=2.2pt,arc=10pt,
        drop shadow=popink,left=11pt,right=11pt,top=10pt,bottom=10pt,height=3.1cm,valign=center]
      \tikz[baseline=-0.7ex]\node[circle,fill=white,draw=popink,line width=1.3pt,minimum size=1.6cm,inner sep=0]{\popdisplay\fontsize{24}{24}\selectfont\color{poppurple}+};%
      \hspace{8pt}\begin{minipage}[c]{0.6\linewidth}
        {\popdisplay\fontsize{11}{12.5}\selectfont\color{white}\MakeUppercase{Passe à OnBuch+}}\par\vspace{4pt}
        {\raggedright\popsemi\fontsize{8.4}{11}\selectfont\color{white}Tous les cours signés, Tuteur IA illimité, fiches PDF — onglet OnBuch+ de l'app\par}
      \end{minipage}
    \end{tcolorbox}
  \end{minipage}
  \vspace{8pt}
  \popcontenu
  \vfill
  \signatureonbuch
  \restoregeometry
  \newpage
}

% Rangée « Ce cours contient » (page de garde)
\newcommand{\poptile}[3]{% #1 couleur #2 gros texte #3 légende
  \begin{tcolorbox}[enhanced,colback=#1,colframe=popink,boxrule=2pt,arc=9pt,shadow={2.5pt}{-2.5pt}{0pt}{popink},
      left=6pt,right=6pt,top=9pt,bottom=9pt,halign=center,valign=center,height=1.75cm,width=\linewidth,nobeforeafter]
    {\popdisplay\fontsize{12.5}{13}\selectfont\color{white}\MakeUppercase{#2}}\par\vspace{3pt}
    {\centering\popsemi\fontsize{7.8}{9.5}\selectfont\color{white}#3\par}
  \end{tcolorbox}}
\newcommand{\popcontenu}{%
  \par\noindent{\popxboldls\fontsize{7.5}{7.5}\selectfont\color{popmuted}CE COURS SIGNÉ CONTIENT}\par\vspace{7pt}
  \noindent\begin{minipage}[t]{0.235\linewidth}\poptile{popblue}{Cours}{complet, illustré, pas à pas}\end{minipage}\hfill
  \begin{minipage}[t]{0.235\linewidth}\poptile{poporange}{Méthodes}{et exercices résolus}\end{minipage}\hfill
  \begin{minipage}[t]{0.235\linewidth}\poptile{poppink}{Exercices}{type examen + corrigés}\end{minipage}\hfill
  \begin{minipage}[t]{0.235\linewidth}\poptile{popgreen}{Fiche bilan}{l'essentiel à retenir}\end{minipage}\par}

% Sommaire
\newtcolorbox{sommaire}{enhanced,colback=white,colframe=popink,boxrule=2.2pt,arc=10pt,
  drop shadow=popink,left=18pt,right=18pt,top=13pt,bottom=13pt,before skip=6pt,after skip=20pt,
  before upper={\poppill{Sommaire}{poppurple}{white}\par\vspace{9pt}\popsemi\fontsize{10.6}{14.5}\selectfont\color{popink}}}

% Signature de fin de cours — poussée tout en bas de la dernière page
\newcommand{\findecours}{%
  \par\vfill\nopagebreak
  \begin{center}\popsquiggle{poporange}\end{center}\vspace{4pt}
  \signatureonbuch
}
\AtBeginDocument{\def\llbracket{\mathopen{[\mkern-3.5mu[}}\def\rrbracket{\mathclose{]\mkern-3.5mu]}}} % intervalles d'entiers (glyphe ⟦ absent de la police)
% Glyphes absents de Fira Math : redéfinis à partir de glyphes disponibles
\AtBeginDocument{%
  \def\setminus{\mathbin{\backslash}}\let\smallsetminus\setminus
  \def\vdots{\mathord{\vbox{\baselineskip3.2pt\lineskiplimit0pt\kern4pt\hbox{.}\hbox{.}\hbox{.}}}}%
  \def\ddots{\mathinner{\mkern1mu\raise6.5pt\hbox{.}\mkern2mu\raise3.6pt\hbox{.}\mkern2mu\raise0.7pt\hbox{.}\mkern1mu}}%
  \def\triangle{\mathord{\scalebox{0.9}{$\symup{\Delta}$}}}%
  \def\nabla{\mathord{\raisebox{\depth}{\scalebox{1}[-1]{$\symup{\Delta}$}}}}%
  \def\checkmark{\popcheck{popdark}}%
  \def\star{\mathbin{\ast}}\def\bigstar{\mathord{\ast}}%
  \def\diamond{\mathbin{\scalebox{0.6}{$\Diamond$}}}%
  \def\bigcup{\mathop{\mathchoice{\raisebox{-0.3em}{\scalebox{1.7}{$\cup$}}}{\scalebox{1.25}{$\cup$}}{\cup}{\cup}}\displaylimits}%
  \def\bigcap{\mathop{\mathchoice{\raisebox{-0.3em}{\scalebox{1.7}{$\cap$}}}{\scalebox{1.25}{$\cap$}}{\cap}{\cap}}\displaylimits}%
  \def\lesseqqgtr{\mathrel{\lessgtr}}\def\bigoplus{\mathop{\oplus}\displaylimits}%
}
\providecommand{\ssi}{si et seulement si} % abréviation fréquente
\AtBeginDocument{\providecommand{\wideparen}{\overparen}} % arc de cercle

%% FIGFIX-BEGIN v4 — correctifs globaux des figures (docs/onbuch-generation/tools/figfix)
\usepackage{adjustbox}\usepackage{etoolbox}
% toute figure TikZ/pgfplots plus large que la ligne est réduite à la largeur disponible
\BeforeBeginEnvironment{tikzpicture}{\begin{adjustbox}{max width=\linewidth}}
\AfterEndEnvironment{tikzpicture}{\end{adjustbox}}
% légendes pgfplots lisibles par-dessus les courbes
\pgfplotsset{every axis/.append style={legend style={fill=white,fill opacity=0.92,text opacity=1,draw=black!15,inner sep=3pt}}}
% étiquettes d'arêtes (syntaxe edge["..."]) avec fond blanc
\tikzset{every edge quotes/.append style={fill=white,inner sep=1.2pt}}
%% FIGFIX-END
\begin{document}

\pagegarde{Production d’écrits utilitaires : le rapport (élaboration du plan détaillé)}{Français}{1re Tech}{Français – classe de Première technique}{N14}{13 séances (fiche de progression harmonisée)}{Cette leçon guide l'élève dans la maîtrise de l'écriture utilitaire, spécifiquement l'élaboration du plan détaillé d'un rapport technique. Elle articule analyse de textes fonctionnels, outils linguistiques (phrase complexe, emprunts) et méthodologie de structuration de l'information.
\vspace{4pt}
\begin{multicols}{2}\small
\begin{itemize}
\item À la fin de cette leçon, je sais identifier les caractéristiques du texte explicatif et du rapport technique.
\item À la fin de cette leçon, je sais distinguer la structure externe (paratexte) et la structure interne (organisation des idées) d'un rapport.
\item À la fin de cette leçon, je sais analyser des textes fonctionnels et iconiques pour en extraire l'information pertinente.
\item À la fin de cette leçon, je sais construire des phrases complexes correctes (subordination, coordination) adaptées à l'écrit technique.
\item À la fin de cette leçon, je sais repérer et expliquer les emprunts lexicaux (francismes, anglicismes, camerounaisismes) dans un corpus professionnel.
\item À la fin de cette leçon, je sais élaborer un plan détaillé (thématique, chronologique ou analytique) pour un rapport de stage ou technique.
\end{itemize}\end{multicols}}

\begin{sommaire}
\begin{multicols}{2}
\begin{enumerate}
\item 1. Le rapport technique : genre, finalités et structures
\item 2. Outils linguistiques : La phrase complexe et l'emprunt lexical
\item 3. Stylistique du rapport : Textes explicatifs, tonalités et objectivité
\item 4. Élaboration du plan détaillé : Méthodologie pas à pas
\item 5. Amélioration du plan et du texte : Révision, cohérence et iconographie
\item 6. Situation d'intégration : Du terrain au rapport final
\item Activité d'intégration
\item Méthodes et exercices résolus
\item Exercices
\item Corrigés des exercices
\item Fiche bilan
\end{enumerate}
\end{multicols}
\end{sommaire}

\begin{objectifs}
\begin{itemize}
\item À la fin de cette leçon, je sais identifier les caractéristiques du texte explicatif et du rapport technique.
\item À la fin de cette leçon, je sais distinguer la structure externe (paratexte) et la structure interne (organisation des idées) d'un rapport.
\item À la fin de cette leçon, je sais analyser des textes fonctionnels et iconiques pour en extraire l'information pertinente.
\item À la fin de cette leçon, je sais construire des phrases complexes correctes (subordination, coordination) adaptées à l'écrit technique.
\item À la fin de cette leçon, je sais repérer et expliquer les emprunts lexicaux (francismes, anglicismes, camerounaisismes) dans un corpus professionnel.
\item À la fin de cette leçon, je sais élaborer un plan détaillé (thématique, chronologique ou analytique) pour un rapport de stage ou technique.
\item À la fin de cette leçon, je sais améliorer un plan et un texte par la révision (cohérence, cohésion, tonalité neutre/objective).
\item À la fin de cette leçon, je sais appliquer les codes de la communication professionnelle par internet (mail d'envoi, format PDF, objet clair).
\end{itemize}
\end{objectifs}

\begin{prerequis}
\begin{itemize}
\item Notion de texte argumentatif et narratif (classe de 4ème/3ème).
\item Maîtrise de la phrase simple (sujet, verbe, compléments) et connecteurs logiques de base.
\item Technique du résumé et du compte-rendu (programme de Seconde).
\item Connaissance basique des logiciels de traitement de texte (Word, Google Docs).
\item Culture générale sur le monde de l'entreprise et les stages au Cameroun.
\end{itemize}
\end{prerequis}

\newpage

\coursec{Le rapport technique : genre, finalités et structures}

\textbf{Situation d'accroche.} En stage de fin d'année dans une entreprise de BTP à Douala, Armand doit rendre un rapport d'activités à son tuteur et au chef de chantier. Il a toutes les notes brutes : dates des travaux, quantités de ciment utilisées, incidents survenus, observations sur le respect des délais. Mais ses feuilles volantes, ses messages WhatsApp et ses photos de chantier forment un désordre total. Il ne sait pas comment structurer tout cela pour que le document soit professionnel, clair et conforme aux attentes de sa hiérarchie. Comment l'aider à transformer ses données en un rapport technique irréprochable ?

C'est précisément l'objet de cette première section : comprendre ce qu'est un rapport, à quoi il sert, et comment il est bâti, de l'enveloppe extérieure jusqu'au squelette logique intérieur.

\courssub{Qu'est-ce qu'un rapport ?}

\textbf{Intuition d'abord.} Imagine que ton chef te demande : « Raconte-moi comment s'est passée la semaine sur le chantier. » Tu peux répondre de vive voix, mais dès demain, tes paroles seront oubliées. Si tu écris un document structuré, objectif et daté, il restera une trace utilisable : on pourra le consulter, le vérifier, prendre des décisions à partir de lui. Ce document, c'est un \cle{rapport}.

\begin{definition}[Le rapport]
Le \cle{rapport} est un texte \textbf{utilitaire} à visée \textbf{informative} ou \textbf{décisionnelle}, rédigé dans un style \textbf{objectif}, \textbf{impersonnel} et \textbf{structuré}. Il rend compte de faits, d'observations ou d'activités à un destinataire déterminé (chef, tuteur, administration, jury), souvent à la demande de celui-ci.
\end{definition}

Trois adjectifs résument l'essentiel :

\begin{itemize}
\item \textbf{Utilitaire} : le rapport sert à quelque chose de concret (informer, justifier, décider). Ce n'est ni un récit d'imagination ni un texte littéraire.
\item \textbf{Explicatif} : il présente les faits, les explique (causes, déroulement, conséquences) et propose parfois des solutions.
\item \textbf{Fonctionnel} : il obéit à des codes précis de présentation et de langue, reconnus par tous les professionnels.
\end{itemize}

\begin{savaistu}[D'où vient le mot « rapport » ?]
Le terme \emph{rapport} vient du latin \emph{reportare}, qui signifie « rapporter, porter en arrière ». Le mot implique donc un \textbf{retour vers un commanditaire} : celui qui rédige le rapport « porte en arrière » l'information vers celui qui l'a demandée. Le rapport n'existe jamais seul : il répond toujours à une demande, explicite ou implicite.
\end{savaistu}

\courssub{Les principaux types de rapports}

Tous les rapports ne se ressemblent pas. Leur contenu et leur organisation dépendent de leur \textbf{finalité}. Voici les quatre types qu'un élève de Première technique doit savoir distinguer.

\begin{center}
\begin{tabularx}{\linewidth}{|l|X|X|}
\hline
\rowcolor{popblueL} \textbf{Type de rapport} & \textbf{Finalité} & \textbf{Exemple de situation} \\
\hline
Rapport de stage & Rendre compte des activités, des compétences acquises et des observations faites pendant un stage. & Armand, stagiaire BTP à Douala, décrit ses trois semaines sur le chantier. \\
\hline
Rapport d'incident & Décrire un événement anormal (accident, panne, vol), ses causes et ses conséquences. & Une grue s'est renversée sur le chantier : le chef d'équipe rédige un rapport le jour même. \\
\hline
Rapport d'activité & Faire le bilan périodique (hebdomadaire, mensuel, annuel) du travail d'une personne ou d'un service. & Le chef de chantier envoie chaque vendredi l'état d'avancement des travaux à la direction. \\
\hline
Rapport d'enquête & Présenter les résultats d'une investigation (questionnaires, mesures, témoignages) et des recommandations. & Une enquête sur les causes des retards de livraison de matériaux dans la région du Littoral. \\
\hline
\end{tabularx}
\end{center}

\begin{attention}[Ne pas confondre les types]
Un même stage peut donner lieu à un rapport de stage (bilan global) ET à un rapport d'incident (si un événement particulier s'est produit). Ce qui change, c'est la \textbf{question à laquelle le rapport répond} : « Qu'ai-je fait et appris ? » (stage), « Que s'est-il passé et pourquoi ? » (incident), « Où en est le travail ? » (activité), « Que révèle l'investigation ? » (enquête).
\end{attention}

\courssub{La structure externe : l'enveloppe du rapport}

\textbf{Observation préalable.} Prends deux documents : d'un côté, une copie d'élève sans titre, sans date, avec des pages dans le désordre ; de l'autre, un document avec page de garde, sommaire, parties numérotées et annexes. Lequel t'inspire confiance avant même de le lire ? Le second, évidemment. La \cle{structure externe} est la présentation matérielle du rapport : c'est la première chose que voit le lecteur, et elle conditionne son jugement.

Un rapport technique complet comprend, dans l'ordre :

\begin{enumerate}
\item \textbf{La page de garde} : titre du rapport, nom et fonction de l'auteur, nom du destinataire, nom de l'organisme, lieu et date, éventuellement logo.
\item \textbf{Le sommaire} : le plan détaillé du rapport, avec la numérotation des parties et sous-parties et les numéros de pages.
\item \textbf{L'introduction} : elle présente le contexte (qui demande le rapport, pourquoi), l'objet du rapport, la méthode suivie et annonce le plan.
\item \textbf{Le développement} : le corps du rapport, divisé en parties et sous-parties numérotées (1., 1.1, 1.2, 2., etc.).
\item \textbf{La conclusion} : bilan, réponse à la question posée, recommandations éventuelles, perspectives.
\item \textbf{Les annexes} : documents complémentaires (tableaux de mesures, photos, formulaires, schémas) renvoyés depuis le développement (« voir annexe 2 »).
\item \textbf{La bibliographie} : liste des documents consultés (ouvrages, sites, normes), si le rapport s'appuie sur des sources.
\end{enumerate}

\begin{attention}[Sommaire ou table des matières ?]
Erreur fréquente : confondre le \textbf{sommaire} et la \textbf{table des matières}. Le sommaire est placé au \emph{début} du rapport : c'est le plan détaillé \textbf{avec numérotation} des parties et sous-parties, qui guide le lecteur. La table des matières se trouve souvent en \emph{fin} d'ouvrage et peut se limiter aux grandes rubriques, parfois sans numérotation détaillée. Dans un rapport de stage de Première technique, c'est le \textbf{sommaire} qui est exigé.
\end{attention}

\begin{exemplebox}[Une page de garde conforme]
\begin{center}
\fbox{\parbox{0.85\linewidth}{\centering
\textbf{RAPPORT DE STAGE}\\[0.3em]
\emph{Suivi du chantier de construction d'un bâtiment R+2 à Bonabéri}\\[0.6em]
Rédigé par : Armand ESSOMBA, élève de Première Génie Civil\\
Tuteur de stage : M. NGUEFACK, chef de chantier\\
Entreprise : SOCABAT SARL, Douala\\
Période : du 5 au 23 janvier 2026}}
\end{center}
\end{exemplebox}

\courssub{La structure interne : le squelette logique}

Si la structure externe est l'enveloppe, la \cle{structure interne} est le squelette : c'est l'ordre logique dans lequel les idées sont disposées dans le développement. Trois grands types de plans sont à connaître.

\textbf{Le plan thématique (par sujets).} On regroupe les informations par grands thèmes, indépendamment du temps. \emph{Exemple} : rapport de stage organisé en « I. Présentation de l'entreprise ; II. Les tâches effectuées ; III. Les compétences acquises ».

\textbf{Le plan chronologique (par étapes).} On suit l'ordre du temps : première semaine, deuxième semaine, troisième semaine ; ou : avant l'incident, pendant, après. Ce plan convient bien aux rapports d'activité et d'incident.

\textbf{Le plan analytique (problème / causes / solutions).} On part d'un constat ou d'un problème, on en analyse les causes, puis on propose des solutions ou des recommandations. C'est le plan roi du rapport d'enquête et du rapport d'incident approfondi.

\begin{methode}[Choisir et construire son plan]
\begin{enumerate}
\item \textbf{Lister les idées} : noter toutes les informations brutes recueillies (notes, mesures, observations), sans les trier.
\item \textbf{Les grouper par affinité} : réunir ce qui porte sur le même sujet ou la même étape.
\item \textbf{Ordonner les groupes} : choisir une logique chronologique, thématique ou déductive (analytique), en fonction de la finalité du rapport.
\item \textbf{Intituler les parties et sous-parties} : donner à chaque groupe un titre clair et numéroté (I, 1.1, 1.2...).
\end{enumerate}
\end{methode}

\begin{exemplebox}[Pourquoi la structure est décisive]
L'expérience des corrections au Probatoire et au Baccalauréat technique le montre chaque année : la majorité des rapports de stage mal notés le sont pour un \textbf{défaut de structure} (plan inexistant, parties déséquilibrées, enchaînement illogique), et non pour la pauvreté du contenu. Autrement dit : un bon contenu mal organisé vaut moins qu'un contenu moyen bien structuré. Travailler le plan, c'est travailler la note.
\end{exemplebox}

\begin{popfigure}
\begin{center}
\begin{tikzpicture}[
  grow=down,
  level 1/.style={sibling distance=32mm, level distance=16mm},
  level 2/.style={sibling distance=22mm, level distance=14mm},
  every node/.style={draw=popblue, fill=popblueL, rounded corners=2pt, font=\small, align=center, inner sep=3pt},
  edge from parent/.style={draw=popdark, -{Stealth[length=2mm]}}
]
\node[fill=poporangeL, draw=poporange, font=\small\bfseries, align=center]{TITRE DU RAPPORT\\(page de garde)}
  child { node[align=center]{Introduction\\contexte, objet, plan} }
  child { node {Partie 1}
    child { node {1.1} }
    child { node {1.2} } }
  child { node {Partie 2}
    child { node {2.1} }
    child { node {2.2} } }
  child { node {Partie 3}
    child { node {3.1} }
    child { node {3.2} } }
  child { node[align=center]{Conclusion\\bilan, recommandations} }
  child { node[align=center]{Annexes\\documents, photos} };
\end{tikzpicture}
\end{center}
\legende{Arbre de l'architecture type d'un rapport technique : de la racine (le titre) aux branches (introduction, parties subdivisées en sous-parties, conclusion, annexes).}
\end{popfigure}

\begin{popfigure}
\begin{center}
\begin{tikzpicture}[font=\small]
% Structure externe
\node[draw=popblue, fill=popblueL, rounded corners=3pt, minimum width=6.2cm, minimum height=0.9cm, align=center] (ext) at (0,4.6) {\textbf{STRUCTURE EXTERNE}\\l'enveloppe physique (ce qu'on voit)};
\node[draw=popblue, fill=white, rounded corners=2pt, align=left, text width=5.8cm] at (0,2.3) {Page de garde $\rightarrow$ Sommaire $\rightarrow$ Introduction $\rightarrow$ Développement $\rightarrow$ Conclusion $\rightarrow$ Annexes $\rightarrow$ Bibliographie};
% Structure interne
\node[draw=poporange, fill=poporangeL, rounded corners=3pt, minimum width=6.2cm, minimum height=0.9cm, align=center] (int) at (8.6,4.6) {\textbf{STRUCTURE INTERNE}\\le squelette logique (ce qu'on comprend)};
\node[draw=poporange, fill=white, rounded corners=2pt, align=left, text width=5.8cm] at (8.6,2.3) {Plan thématique (par sujets)\\ Plan chronologique (par étapes)\\ Plan analytique (problème / causes / solutions)};
% Lien
\draw[-{Stealth[length=3mm]}, very thick, popgreen] (3.4,3.6) -- (5.2,3.6) node[midway, above, popgreen, align=center, font=\footnotesize] {s'emboîtent};
\node[align=center, font=\footnotesize\itshape, popmuted] at (4.3,0.8) {Un bon rapport = une enveloppe soignée\\ + un squelette logique cohérent};
\end{tikzpicture}
\end{center}
\legende{Schéma comparatif : la structure externe (présentation matérielle du document) et la structure interne (organisation logique des idées) sont complémentaires et indispensables.}
\end{popfigure}

\courssub{Étude d'un texte fonctionnel : extrait de rapport de stage au MINTP}

Lisons maintenant un texte fonctionnel type (Texte N°1) : l'introduction d'un rapport de stage rédigé par un élève de Première technique accueilli au MINTP (Ministère des Travaux Publics), délégation régionale du Littoral.

\begin{exemplebox}[Texte N°1 : extrait de rapport de stage au MINTP]
\emph{« Dans le cadre de la formation pratique prévue au programme de la classe de Première Génie Civil, j'ai effectué un stage d'observation et de participation du 5 au 23 janvier 2026 à la délégation régionale du MINTP du Littoral, service Études et Suivi des chantiers.\\
Ce stage avait pour objectifs de découvrir l'organisation d'un service technique de l'État, de participer au suivi d'un chantier de réhabilitation routière et d'appliquer les notions de topographie vues en classe.\\
Le présent rapport rend compte des activités menées. Il s'organise en trois parties : la première présente la structure d'accueil, la deuxième décrit les tâches effectuées, la troisième fait le bilan des compétences acquises et propose des recommandations. »}
\end{exemplebox}

\textbf{Analyse guidée.} Ce court extrait remplit parfaitement les trois fonctions d'une introduction de rapport :

\begin{enumerate}
\item \textbf{Le contexte} : « Dans le cadre de la formation pratique... » — on sait qui écrit, où, quand et pourquoi.
\item \textbf{L'objet} : « Ce stage avait pour objectifs de... » — la finalité est énoncée clairement.
\item \textbf{L'annonce du plan} : « Il s'organise en trois parties... » — le lecteur connaît déjà la structure interne, ici un plan thématique.
\end{enumerate}

On remarque aussi le style : phrases déclaratives, vocabulaire technique précis (\emph{réhabilitation routière}, \emph{topographie}), ton objectif, absence de jugements personnels ou d'émotions. C'est le registre explicatif attendu dans tout rapport.

\begin{exoresolu}[Identifier les éléments de structure dans le Texte N°1]
\textbf{Énoncé.} À partir de l'extrait du rapport de stage au MINTP ci-dessus : 1) Relevez les informations qui appartiennent normalement à la page de garde. 2) Identifiez le type de plan interne annoncé et justifiez. 3) Montrez en quoi le style est « objectif et impersonnel ».
\tcblower
\textbf{Solution détaillée.}
\begin{enumerate}
\item \textbf{Éléments de page de garde} : le nom de la structure d'accueil (délégation régionale du MINTP du Littoral), la période du stage (du 5 au 23 janvier 2026), la classe et la spécialité de l'auteur (Première Génie Civil). Sur la page de garde réelle s'ajouteraient le titre du rapport, le nom de l'auteur et celui du tuteur.
\item \textbf{Type de plan} : il s'agit d'un \textbf{plan thématique} (par sujets) : structure d'accueil / tâches effectuées / bilan et recommandations. Les parties ne suivent pas l'ordre du temps mais regroupent les informations par thèmes. La troisième partie (bilan, recommandations) introduit une légère dimension analytique.
\item \textbf{Style objectif et impersonnel} : aucune émotion, aucune opinion personnelle (« j'ai adoré », « c'était génial » seraient hors de propos) ; le « je » n'apparaît qu'une fois, pour situer le contexte, et le rapport lui-même devient le sujet de la phrase (« le présent rapport rend compte ») ; le vocabulaire est technique et précis ; les verbes sont à l'indicatif, mode du constat.
\end{enumerate}
\end{exoresolu}

\courssub{Le rapport et la communication professionnelle par internet}

Aujourd'hui, un rapport se transmet le plus souvent par courrier électronique. Cette transmission obéit à des \cle{codes} précis que tout technicien doit maîtriser.

\begin{itemize}
\item \textbf{L'objet du mail} : court et explicite, il fonctionne comme un mini-titre. \emph{Exemple} : « Rapport de stage — A. ESSOMBA — SOCABAT Douala (janvier 2026) ». Un objet vide ou vague (« bonjour », « mon document ») est une faute professionnelle.
\item \textbf{La pièce jointe} : le rapport est envoyé en \textbf{PDF}, format qui garantit que la mise en page ne sera pas modifiée. Le nom du fichier doit être clair : \emph{Rapport\_stage\_Essomba\_2026.pdf} et non \emph{document final2.pdf}.
\item \textbf{Le corps du message} : bref, il annonce la pièce jointe et sa finalité.
\item \textbf{La formule de politesse} : adaptée au destinataire et à la hiérarchie : « Veuillez agréer, Monsieur le Chef de chantier, l'expression de mes salutations distinguées » pour un supérieur ; « Cordialement » pour un échange entre collègues.
\end{itemize}

\begin{exemplebox}[Mail d'accompagnement du rapport d'Armand]
\emph{Objet : Rapport d'activités — chantier Bonabéri (semaine du 12 au 16 janvier 2026)}\\
« Monsieur,\\
Conformément à votre demande, vous trouverez ci-joint le rapport d'activités de la semaine écoulée sur le chantier de Bonabéri, au format PDF.\\
Je reste à votre disposition pour tout complément d'information.\\
Veuillez agréer, Monsieur, l'expression de mes salutations distinguées.\\
Armand ESSOMBA, stagiaire Génie Civil »
\end{exemplebox}

\begin{attention}[Les pièges de la transmission numérique]
Trois erreurs discréditent un bon rapport : oublier la pièce jointe (relire avant d'envoyer !), envoyer un fichier modifiable au lieu d'un PDF (la mise en page peut être déformée), et utiliser un ton trop familier dans le mail (« Salut chef, voici le truc »). Le mail est un document professionnel : il fait partie du rapport au sens large.
\end{attention}

\begin{aretenir}[L'essentiel de la section]
\begin{itemize}
\item Le rapport est un texte \textbf{utilitaire, explicatif et fonctionnel}, à visée informative ou décisionnelle, rédigé dans un style objectif et impersonnel.
\item On distingue quatre types : rapport de \textbf{stage}, d'\textbf{incident}, d'\textbf{activité}, d'\textbf{enquête}.
\item La \textbf{structure externe} (page de garde, sommaire, introduction, développement, conclusion, annexes, bibliographie) est l'enveloppe ; la \textbf{structure interne} (plan thématique, chronologique ou analytique) est le squelette logique.
\item Le sommaire, placé au début, est le plan détaillé numéroté ; il ne faut pas le confondre avec la table des matières.
\item La transmission par internet obéit à des codes : objet explicite, pièce jointe PDF bien nommée, formule de politesse adaptée.
\end{itemize}
\end{aretenir}

Armand dispose désormais de tous les repères nécessaires : il sait ce qu'est un rapport, quel type rédiger, comment l'habiller extérieurement et comment choisir le squelette logique de son développement. La section suivante lui fournira les outils linguistiques — la phrase complexe et le vocabulaire emprunté — pour rédiger chaque partie avec correction et précision.

\coursec{Outils linguistiques : La phrase complexe et l'emprunt lexical}

Après avoir identifié le genre du rapport technique, ses finalités (informer, décider, tracer) et ses structures canoniques (structure externe : page de garde, sommaire, corps, annexes ; structure interne : introduction, développement, conclusion), nous devons désormais outiller notre écriture. Un rapport technique ne se lit pas comme un roman : il se \emph{scanne}. La clarté y est une contrainte contractuelle. Deux leviers linguistiques majeurs permettent d'atteindre cette exigence : la \textbf{syntaxe complexe} pour articuler la logique technique, et la \textbf{lexicologie spécialisée} (notamment l'emprunt) pour nommer la réalité du terrain avec précision. C'est l'objet de cette section.

\courssub{La phrase complexe : l'architecture de la pensée technique}

\courssub{Du style télégraphique à la fluidité logique}

Dans les notes de terrain, sur le chantier ou en réunion, la communication est souvent \emph{télégraphique} : phrases nominales, juxtapositions, énumérations brutes. Exemple : \og{}Pluie forte. Arrêt travaux. Retard livraison béton. Pénalité contrat.\fg{} Ce style de note est indispensable pour la prise de notes rapide, mais inacceptable dans un rapport final. Le lecteur (chef de projet, client, bailleur de fonds) doit saisir les \textbf{relations logiques} entre les faits : cause, conséquence, but, condition, opposition.

\begin{propriete}[La subordination, colonne vertébrale de l'explication technique]
La subordination permet d'établir des liens logiques \textbf{explicites} (cause, but, condition, concession, temps) indispensables à l'explication technique. Elle hiérarchise l'information : l'idée principale (proposition principale) porte l'énoncé central ; les idées secondaires (propositions subordonnées) la précisent, la justifient ou la conditionnent. Repère de rédaction pour un rapport technique : viser en moyenne \textbf{2 à 3 propositions par phrase}, sans jamais dépasser 4 propositions dans une même phrase.
\end{propriete}

\courssub{Rappel structuré : proposition et types de liaisons}

Rappelons les trois modes d'enchaînement des propositions (une proposition est un groupe de mots construit autour d'un verbe conjugué) :

\begin{center}
\begin{tabularx}{\linewidth}{|l|X|X|X|}
\hline
\rowcolor{popblueL}
\textbf{Type de liaison} & \textbf{Moyen de liaison} & \textbf{Relation logique principale} & \textbf{Exemple technique} \\
\hline
\textbf{Juxtaposition} & Ponctuation forte (., ;, :) & Simple succession (parataxe) & Le béton a séché. Le coffrage est retiré. \\
\hline
\textbf{Coordination} & Conjonctions de coordination (\textit{mais, ou, et, donc, or, ni, car}) & Égalité syntaxique : addition, opposition, alternative, conséquence & Le délai est court, \textbf{donc} nous devons optimiser les équipes. \\
\hline
\textbf{Subordination} & Mots subordonnants (conjonctions, pronoms relatifs) & Hiérarchie : \textit{Principale} $\rightarrow$ \textit{Subordonnée} & \textbf{Parce que} le délai est court, nous optimisons les équipes. \\
\hline
\end{tabularx}
\end{center}

\courssub{Focus sur les subordonnées : les \og{}moteurs logiques\fg{} du rapport}

Trois types de subordonnées sont massivement utilisés dans le rapport technique. Les maîtriser, c'est maîtriser l'argumentation.

\begin{definition}[Subordonnée relative (déterminative ou explicative)]
Introduite par un pronom relatif (\textit{qui, que, dont, où, lequel...}) qui a un \textbf{antécédent} dans la principale. Elle précise un nom (technique, matériel, délai, personne).
\begin{itemize}
    \item \textit{Déterminative (sans virgule) :} elle identifie l'antécédent, elle est indispensable au sens. \og{}Le \textbf{chef de chantier} \textbf{qui a signé le PV} est responsable.\fg{}
    \item \textit{Explicative (entre virgules) :} elle ajoute une information sur un antécédent déjà identifié. \og{}Le \textbf{chef de chantier}, \textbf{qui a quinze ans d'expérience}, a validé le planning.\fg{}
\end{itemize}
\end{definition}

\begin{definition}[Subordonnée conjonctive circonstancielle]
Introduite par une conjonction de subordination. Elle fonctionne comme un \textbf{complément circonstanciel} du verbe principal. Les quatre types essentiels pour le rapport :
\begin{center}
\begin{tabularx}{\linewidth}{|l|X|X|}
\hline
\rowcolor{popblueL}
\textbf{Type} & \textbf{Subordonnants fréquents} & \textbf{Exemple en contexte BTP} \\
\hline
\textbf{Cause} & \textit{parce que, comme, puisque, du fait que} & \textit{Du fait que} le sol est argileux, nous avons prévu des fondations profondes. \\
\hline
\textbf{Conséquence} & \textit{si bien que, de sorte que, au point que} & La pluie a été continue, \textit{de sorte que} le terrassement a pris trois jours de retard. \\
\hline
\textbf{But} & \textit{afin que, pour que, de peur que} (+ subjonctif) & Le chef a ordonné le nettoyage \textit{afin que} la réception se passe sans réserve. \\
\hline
\textbf{Condition} & \textit{si, pourvu que, à condition que} & \textit{Si} le béton atteint la résistance requise à 28 jours, \textit{alors} le décoffrage est autorisé. \\
\hline
\end{tabularx}
\end{center}
\end{definition}

\begin{definition}[Subordonnée complétive (sujet ou objet)]
Introduite par \textit{que} (ou \textit{si} pour l'interrogation indirecte). Elle fonctionne comme un \textbf{COD} ou un \textbf{sujet} du verbe principal (souvent un verbe de parole, de pensée, de volonté ou de perception).
\begin{itemize}
    \item \textit{Complétive objet :} Le directeur \textbf{exige que} le rapport \textbf{soit remis} avant vendredi. (Le verbe \textit{exiger} impose le subjonctif.)
    \item \textit{Complétive sujet :} \textbf{Que} le chantier respecte les normes \textbf{est} une obligation légale.
\end{itemize}
\end{definition}

\courssub{Illustration : analyse en constituants (arbre syntaxique)}

L'arbre ci-dessous montre l'imbrication typique d'une phrase de rapport : une principale (\textit{exiger}) gouvernant une complétive objet (\textit{que le rapport soit remis...}), elle-même étendue par une subordonnée de but (\textit{pour que le paiement soit déclenché}).

\begin{popfigure}
\begin{tikzpicture}[
    level distance=1.5cm,
    level 1/.style={sibling distance=5.5cm},
    level 2/.style={sibling distance=3.2cm},
    level 3/.style={sibling distance=2.6cm},
    edge from parent/.style={draw=popdark, thick},
    every node/.style={font=\footnotesize, align=center},
    boite/.style={draw=popblue, fill=popblueL, rounded corners=3pt, inner sep=3pt},
    feuille/.style={draw=popgreen, fill=popgreenL, rounded corners=3pt, inner sep=2pt},
    outil/.style={draw=poporange, fill=poporangeL, rounded corners=3pt, inner sep=2pt}
]
\node [boite] {Phrase}
    child { node [feuille, align=center]{SN sujet\\Le chef de chantier} }
    child { node [boite] {GV : a exigé}
        child { node [outil, align=center]{Subordonnant\\que} }
        child { node [boite] {Complétive objet}
            child { node [feuille, align=center]{SN sujet\\le rapport} }
            child { node [feuille, align=center]{GV passif\\soit remis} }
            child { node [feuille, align=center]{CC temps\\avant vendredi} }
            child { node [boite] {Subordonnée de but}
                child { node [outil] {pour que} }
                child { node [feuille, align=center]{le paiement\\soit déclenché} }
            }
        }
    };
\end{tikzpicture}
\legende{Arbre syntaxique simplifié de la phrase : \og{}Le chef de chantier a exigé que le rapport soit remis avant vendredi pour que le paiement soit déclenché.\fg{} La structure hiérarchique (principale $\rightarrow$ complétive $\rightarrow$ but) rend explicite la chaîne contractuelle : exigence $\rightarrow$ livrable $\rightarrow$ déclencheur financier.}
\end{popfigure}

\courssub{Manipulation guidée : du style télégraphique au style technique}

\begin{methode}[Complexifier une suite de phrases simples]
\textbf{Objectif :} transformer un \og{}style note\fg{} en phrase(s) complexe(s) fluide(s) et logique(s).
\begin{enumerate}
    \item \textbf{Segmenter et identifier :} repérer les \textbf{propositions} (verbes conjugués) et les \textbf{entités techniques} (sujets, objets).
    \item \textbf{Repérer le lien logique :} cause ? but ? condition ? opposition ? chronologie ? (Se demander : \textit{pourquoi ? pour quoi faire ? dans quel cas ?})
    \item \textbf{Choisir le subordonnant adapté :} \textit{parce que / comme} (cause) ; \textit{afin que / pour que} + subjonctif (but) ; \textit{si / à condition que} (condition) ; \textit{bien que / alors que} (concession, opposition) ; \textit{lorsque / dès que} (temps).
    \item \textbf{Construire la phrase principale :} choisir le verbe \textbf{porteur de l'information nouvelle} (souvent le résultat, la décision, l'observation clé).
    \item \textbf{Placer les subordonnées :} en début, au milieu ou à la fin selon ce qu'on veut mettre en avant (la \textbf{thématisation}).
    \item \textbf{Vérifier les temps et les modes :} indicatif (fait réel) ou subjonctif (volonté, nécessité, but) ; concordance des temps (\textit{si} + imparfait $\rightarrow$ conditionnel présent).
\end{enumerate}
\end{methode}

\begin{exoresolu}[Transformation : style note $\rightarrow$ style rapport]
\textbf{Énoncé (notes de terrain) :}
\og{}Réunion chantier 14/03. Présence : chef de projet, bureau de contrôle, entreprise. Problème : fissures mur Est. Cause : tassement différentiel. Décision : injections de résine. But : stabiliser structure. Délai : 15 jours. Condition : validation budget par maître d'ouvrage.\fg{}

Transformez ces notes en une ou deux phrases de rapport.
\tcblower
\textbf{Solution détaillée :}

\textbf{Étape 1 — Analyse des liens logiques :}
\begin{itemize}
    \item \textit{Fissures} $\leftarrow$ \textbf{Cause} : tassement différentiel.
    \item \textit{Injections} $\rightarrow$ \textbf{But} : stabiliser la structure.
    \item \textit{Injections} $\rightarrow$ \textbf{Condition} : validation du budget.
    \item \textit{Délai de 15 jours} $\rightarrow$ \textbf{Circonstance de temps}.
\end{itemize}

\textbf{Étape 2 — Proposition de rédaction :}
\begin{quote}
\og{}Lors de la réunion de chantier du 14 mars, en présence du chef de projet, du bureau de contrôle et de l'entreprise, il a été constaté que des fissures étaient apparues sur le mur Est \textit{en raison d'un tassement différentiel}. Il a donc été décidé de réaliser des injections de résine \textit{afin de stabiliser la structure}, dans un délai de quinze jours, \textit{à condition que le budget soit validé par le maître d'ouvrage}.\fg{}
\end{quote}

\textbf{Étape 3 — Justification syntaxique :}
\begin{itemize}
    \item \textit{Lors de la réunion...} : complément circonstanciel de temps, placé en tête pour situer le contexte.
    \item \textit{que des fissures étaient apparues...} : subordonnée complétive, COD de \textit{a été constaté} (tournure passive, adaptée au style objectif du rapport).
    \item \textit{en raison d'un tassement différentiel} : complément de cause.
    \item \textit{afin de stabiliser la structure} : infinitive de but (équivalent de \textit{afin que la structure soit stabilisée}).
    \item \textit{à condition que le budget soit validé...} : subordonnée de condition, avec le subjonctif présent \textit{soit validé} après \textit{à condition que}.
\end{itemize}
La rédaction tient en deux phrases de 3 et 4 propositions : l'information est dense mais reste lisible.
\end{exoresolu}

\begin{exercice}[Application rapide]
Transformez ces notes en une phrase complexe unique (ou deux si l'ensemble devient trop dense) :
\og{}Essais béton à 28 jours. Résultats conformes à la résistance exigée. Décoffrage autorisé. Surveillance de la fissuration maintenue.\fg{}

\textit{Indice : verbe principal possible = \textbf{autoriser} (au passif). Liens à exprimer : cause (résultats conformes), conséquence (décoffrage), but (surveillance maintenue).}
\end{exercice}

\begin{corrige}[Exercice : application rapide]
Proposition de rédaction : \og{}Les essais de compression réalisés sur le béton à 28 jours ayant donné des résultats conformes à la résistance exigée, le décoffrage a été autorisé, \textit{la surveillance de la fissuration étant toutefois maintenue}.\fg{} Variante avec subordonnée de cause explicite : \og{}Parce que les essais à 28 jours ont donné des résultats conformes, le décoffrage a été autorisé ; la surveillance de la fissuration reste cependant maintenue afin de garantir la durabilité de l'ouvrage.\fg{}
\end{corrige}

\begin{attention}[Pièges fréquents : la \og{}phrase valise\fg{} et le subjonctif oublié]
\begin{itemize}
    \item \textbf{La phrase valise :} ne rédigez pas une phrase de quinze lignes avec six subordonnées imbriquées. Au-delà de 3 ou 4 propositions, \textbf{coupez} la phrase.
    \item \textbf{Subjonctif obligatoire :} après \textit{afin que, pour que, à condition que, il faut que, il est nécessaire que} $\rightarrow$ \textbf{subjonctif}.
    \item \textbf{Erreur type :} \og{}Il faut que le béton \textbf{est} sec.\fg{} $\rightarrow$ \og{}Il faut que le béton \textbf{soit} sec.\fg{}
    \item \textbf{Concordance des temps :} \textit{Si le béton \textbf{était} sec (imparfait), on \textbf{pourrait} décoffrer (conditionnel).} Ne pas écrire : \textit{Si le béton \textbf{serait} sec...} (le conditionnel est interdit dans la subordonnée introduite par \textit{si}).
\end{itemize}
\end{attention}

\courssub{L'emprunt lexical : nommer la réalité technique au Cameroun}

\courssub{Définition et enjeu terminologique}

Le rapport technique manipule des \textbf{concepts spécialisés} (matériaux, procédés, normes, outils numériques, procédures administratives). La langue française, langue de scolarisation et d'administration au Cameroun, ne possède pas toujours de mot propre pour désigner ces réalités récentes ou importées. Le recours à l'\textbf{emprunt lexical} est donc une nécessité, non un défaut de style.

\begin{definition}[L'emprunt lexical]
L'emprunt est l'\textbf{intégration d'un mot (forme et sens)} d'une \textbf{langue source} dans une \textbf{langue d'accueil}, ici le français technique camerounais. Il répond à un \textbf{besoin de désignation} : la langue d'accueil ne possède pas de mot pour cette réalité.
\begin{itemize}
    \item \textbf{Emprunt non intégré (xénisme) :} il conserve la graphie et la prononciation de la langue source. Il s'écrit en \textbf{italique}. Exemples : \textit{software, hardware, benchmark, workflow}.
    \item \textbf{Emprunt intégré (francisé) :} il subit une adaptation phonétique, graphique et \textbf{morphologique} (marque du pluriel en \textit{-s}, genre grammatical fixé). Il s'écrit en \textbf{romain} (caractères droits). Exemples : \textbf{un planning, des plannings ; un reporting, des reportings}.
\end{itemize}
\end{definition}

\courssub{Typologie des emprunts dans le rapport technique camerounais}

Le vocabulaire technique utilisé au Cameroun (BTP, génie civil, topographie, informatique industrielle, administration) révèle trois strates principales.

\begin{center}
\begin{tabularx}{\linewidth}{|l|X|X|X|}
\hline
\rowcolor{popblueL}
\textbf{Strate} & \textbf{Langue source} & \textbf{Statut graphique} & \textbf{Exemples représentatifs} \\
\hline
\textbf{Anglicismes techniques} & Anglais (langue de l'ingénierie et du numérique) & Souvent \textit{non intégrés} (italique) pour les termes récents ; \textbf{intégrés} (romain) pour les termes anciens & \textit{software, hardware, cloud, dashboard, deadline} ; \textbf{planning, reporting, brainstorming, manager} \\
\hline
\textbf{Francismes administratifs} & Français de l'administration et des marchés publics & \textbf{Intégrés} (romain) & \textbf{procès-verbal (PV), appel d'offres, cahier des charges, ordre de service, réception provisoire, retenue de garantie, maître d'ouvrage, maître d'œuvre, bureau de contrôle} \\
\hline
\textbf{Camerounaisismes de terrain} & Français populaire camerounais, langues nationales, pidgin & \textbf{Intégrés} (romain) dans l'usage professionnel local & \textbf{tôle bac} (tôle ondulée), \textbf{parpaing} (bloc de béton), \textbf{fer à béton} (armature), \textbf{gravier latéritique}, \textbf{tout-venant}, \textbf{gâchée}, \textbf{banche}, \textbf{étai}, \textbf{taloche} \\
\hline
\end{tabularx}
\end{center}

\courssub{Illustration quantitative : répartition dans un corpus de rapports de stage}

Le diagramme ci-dessous synthétise l'analyse d'un corpus fictif mais représentatif de 10 rapports de stage en BTP (environ 50 000 mots), rédigés par des élèves de Première technique. Il illustre la dominance des anglicismes techniques (normes, logiciels, matériel) et la place significative des camerounaisismes (réalité du terrain, matériaux locaux, organisation du chantier).

\begin{popfigure}
\begin{tikzpicture}
\begin{axis}[
    ybar=5pt,
    bar width=22pt,
    width=0.9\linewidth,
    height=8cm,
    enlarge x limits=0.35,
    symbolic x coords={Anglicismes,Francismes,Camerounaisismes},
    xtick=data,
    x tick label style={font=\small, text=popdark},
    ylabel={Part des emprunts (en \%)},
    ylabel style={font=\small, text=popdark},
    ymin=0, ymax=55,
    ymajorgrids=true,
    grid style={dashed, popline},
    nodes near coords={\pgfmathprintnumber\pgfplotspointmeta\,\%},
    nodes near coords style={font=\footnotesize\bfseries, text=popdark},
]
\addplot[fill=popblue, draw=popblue!80!black] coordinates {(Anglicismes,45)};
\addplot[fill=poporange, draw=poporange!80!black] coordinates {(Francismes,30)};
\addplot[fill=popgreen, draw=popgreen!80!black] coordinates {(Camerounaisismes,25)};
\end{axis}
\end{tikzpicture}
\legende{Diagramme en barres : répartition typologique des emprunts lexicaux dans un corpus de 10 rapports de stage BTP (données indicatives). Les anglicismes dominent (45\,\%) pour le vocabulaire technique et du management ; les francismes administratifs (30\,\%) structurent la partie contractuelle ; les camerounaisismes (25\,\%) ancrent le rapport dans la réalité matérielle du terrain. Total : 45 + 30 + 25 = 100\,\%.}
\end{popfigure}

\courssub{Règles graphiques et morphologiques : italique ou romain ?}

C'est un critère de correction formelle au Probatoire. La règle est simple : \textbf{l'usage dicte la graphie}.

\begin{definition}[Règle de l'italique et du romain pour l'emprunt]
\begin{itemize}
    \item \textbf{Italique (xénisme) :} emprunt \textbf{non francisé}, encore perçu comme étranger. Il ne prend pas la marque du pluriel français (on n'écrit pas \textit{softwares}). Exemple : \og{}Le \textit{software} de conception ; des \textit{hardware} performants.\fg{}
    \item \textbf{Romain (emprunt intégré) :} emprunt \textbf{francisé}, entré dans les dictionnaires. Il suit la morphologie française : \textbf{pluriel en -s}, genre fixé. Exemples : \og{}un \textbf{planning} $\rightarrow$ des \textbf{plannings}\fg{} ; \og{}un \textbf{reporting} $\rightarrow$ des \textbf{reportings}\fg{}.
    \item \textbf{Zone de transition :} certains termes oscillent (\textit{email / e-mail / courriel}). \textbf{Recommandation :} privilégier le terme français officiel (\textbf{courriel}) dans la partie administrative du rapport ; tolérer \textit{email} (italique) dans une citation ou une annexe technique.
\end{itemize}
\end{definition}

\begin{exemplebox}[Cas pratiques : choisir la bonne graphie]
\begin{enumerate}
    \item \og{}Le \textit{software} de gestion de projet a planté.\fg{} $\rightarrow$ Mieux : \og{}Le \textbf{logiciel} de gestion de projet a planté.\fg{} (terme français officiel, en romain).
    \item \og{}Nous avons établi trois \textbf{plannings} prévisionnels.\fg{} $\rightarrow$ Correct : \textit{planning} est intégré, il s'écrit en romain et prend le \textit{-s} du pluriel.
    \item \og{}Le \textit{feedback} du client est attendu.\fg{} $\rightarrow$ Mieux en rédaction formelle : \og{}Le \textbf{retour} du client est attendu.\fg{}
    \item \og{}Il faut commander des \textbf{tôles bac}.\fg{} $\rightarrow$ Correct : camerounaisisme intégré, en romain ; le pluriel porte sur \textit{tôle}.
\end{enumerate}
\end{exemplebox}

\begin{attention}[Erreur fréquente : l'abus d'anglicismes \og{}de confort\fg{}]
L'abus d'anglicismes non adaptés (\og{}faire un \textit{reporting}\fg{} au lieu de \og{}rédiger un compte rendu\fg{} ; \og{}le \textit{process}\fg{} au lieu de \og{}la procédure / le processus\fg{}) \textbf{alourdit le style et trahit le registre soutenu} attendu dans un rapport officiel.
\begin{itemize}
    \item \textbf{Test 1 :} le mot figure-t-il dans un dictionnaire français courant ? $\rightarrow$ Alors il s'écrit en romain.
    \item \textbf{Test 2 :} existe-t-il un équivalent français précis et court, utilisé par les professionnels francophones ? $\rightarrow$ Alors \textbf{privilégier le français}.
    \item \textbf{Exception :} les sigles techniques universels (\textbf{BIM, CAO, DAO, GPS, SIG, ISO}) s'écrivent en \textbf{capitales romaines}, sans italique.
\end{itemize}
\end{attention}

\courssub{Exercice résolu : repérage et correction graphique}

\begin{exoresolu}[Correction d'extraits de rapports d'élèves]
\textbf{Énoncé :} Repérez les emprunts dans les phrases suivantes. Corrigez la graphie (italique ou romain) et proposez un équivalent français lorsque l'emprunt est abusif.
\begin{enumerate}
    \item \og{}Le \textit{chef de projet} a validé le \textit{planning} définitif.\fg{}
    \item \og{}Il faut faire un \textit{reporting} hebdomadaire au \textit{top management}.\fg{}
    \item \og{}Les \textit{fer à béton} sont en \textit{attente} sur le \textit{chantier}.\fg{}
    \item \og{}Le \textit{software} de modélisation BIM a détecté un \textit{clash}.\fg{}
\end{enumerate}
\tcblower
\textbf{Solution détaillée :}
\begin{enumerate}
    \item \textbf{Chef de projet} : terme français intégré $\rightarrow$ romain. \textbf{Planning} : anglicisme intégré $\rightarrow$ romain. Correction : \og{}Le \textbf{chef de projet} a validé le \textbf{planning} définitif.\fg{}
    \item \textbf{Reporting} : anglicisme intégré $\rightarrow$ romain (ou mieux : \textbf{compte rendu}). \textbf{Top management} : anglicisme non intégré et abusif en rédaction formelle $\rightarrow$ remplacer par \textbf{la direction générale}. Correction : \og{}Il faut établir un \textbf{compte rendu} hebdomadaire pour la \textbf{direction générale}.\fg{}
    \item \textbf{Fer à béton} : camerounaisisme technique intégré (équivalent : \textbf{armature}) $\rightarrow$ romain, avec le pluriel \textbf{fers à béton}. \textbf{Chantier} : mot français $\rightarrow$ romain. Correction : \og{}Les \textbf{fers à béton} sont en \textbf{stock} sur le \textbf{chantier}.\fg{}
    \item \textbf{Software} : non intégré $\rightarrow$ remplacer par \textbf{logiciel}. \textbf{BIM} : sigle universel $\rightarrow$ capitales romaines. \textbf{Clash} : anglicisme technique récent $\rightarrow$ remplacer par \textbf{interférence} ou \textbf{conflit géométrique}. Correction : \og{}Le \textbf{logiciel} de modélisation \textbf{BIM} a détecté une \textbf{interférence}.\fg{}
\end{enumerate}
\textbf{Bilan :} sur les emprunts repérés, trois étaient abusifs (remplaçables par un terme français précis) et deux étaient mal graphiés. La vigilance terminologique est une compétence évaluée au Probatoire.
\end{exoresolu}

\begin{aretenir}[Synthèse : outils linguistiques pour le rapport technique]
\begin{itemize}
    \item \textbf{Phrase complexe :} outil de la \textbf{logique}. Subordination (cause, but, condition) $>$ coordination $>$ juxtaposition. Viser 2 à 3 propositions par phrase.
    \item \textbf{Subjonctif :} obligatoire après les verbes de volonté et d'exigence et après les subordonnants de but et de condition (\textit{exiger que, il faut que, afin que, pour que, à condition que}).
    \item \textbf{Emprunt lexical :} une nécessité technique. Distinction stricte entre \textbf{italique (xénisme non intégré)} et \textbf{romain (emprunt intégré, pluriel en -s)}.
    \item \textbf{Trois strates au Cameroun :} anglicismes (technique et management), francismes (administration et marchés publics), camerounaisismes (terrain et matériaux locaux).
    \item \textbf{Éthique du terme :} privilégier le français officiel ; tolérer l'emprunt intégré ; réserver l'italique au xénisme technique sans équivalent.
\end{itemize}
\end{aretenir}

\begin{savaistu}[Le \og{}français de spécialité\fg{} au Cameroun : une norme en construction]
Le français technique camerounais n'est pas un \og{}mauvais français\fg{}. C'est une \textbf{variété légitime} du français, façonnée par les réalités professionnelles locales et étudiée par la recherche (travaux du CERDOTOLA, de l'Agence universitaire de la Francophonie, lexiques du français d'Afrique). Des termes comme \og{}tôle bac\fg{}, \og{}fer à béton\fg{} ou \og{}gravier latéritique\fg{} désignent des réalités sans équivalent exact dans le français hexagonal : les employer en romain dans un rapport adressé à une administration camerounaise (MINTP, MINDUH) est \textbf{correct et professionnel}. En revanche, écrire \og{}faire un reporting\fg{} dans un procès-verbal de réception officiel relève de l'\textbf{anglicisme de confort} : c'est une faute de registre, sanctionnée à l'examen.
\end{savaistu}

\coursec{Stylistique du rapport : Textes explicatifs, tonalités et objectivité}

Dans la section précédente, nous avons appris à construire des phrases complexes correctes et à enrichir notre vocabulaire grâce aux emprunts lexicaux. Mais une phrase grammaticalement parfaite peut encore être un mauvais texte de rapport. Pourquoi ? Parce que le rapport n'obéit pas seulement à des règles de langue : il obéit à une \cle{stylistique}, c'est-à-dire à une manière particulière de dire les choses. Un rapport qui raconte, qui émeut ou qui plaisante manque sa cible. Cette section vous apprend à adopter la \cle{tonalité neutre} et objective qui fait la crédibilité d'un document technique.

\courssub{Le texte explicatif : le modèle stylistique du rapport}

\textbf{Intuition : expliquer n'est ni raconter, ni convaincre}\par\medskip

Imaginez trois façons de parler du même événement, la mise en service d'une pompe à eau dans un village :

\begin{itemize}
\item \emph{Récit} : « Ce matin-là, nous sommes arrivés très tôt. Soudain, l'eau a jailli et tout le monde a applaudi. » (on raconte une histoire, dans l'ordre du temps) ;
\item \emph{Texte argumentatif} : « Il faut absolument installer des pompes partout, car l'eau est un droit. » (on cherche à convaincre) ;
\item \emph{Texte explicatif} : « La pompe fonctionne par aspiration : le mouvement du piston crée une dépression qui fait monter l'eau de la nappe phréatique vers la canalisation. » (on fait comprendre un mécanisme).
\end{itemize}

Le rapport technique appartient à la troisième famille.

\begin{definition}[Le texte explicatif]
Le \cle{texte explicatif} vise à faire comprendre un phénomène, un processus ou une situation, sans chercher à persuader ni à émouvoir. Sa visée est \emph{informative} : il répond aux questions « comment ? » et « pourquoi ? » par des faits, des mécanismes et des causes.
\end{definition}

\textbf{Les quatre caractéristiques du texte explicatif}\par\medskip

\begin{propriete}[Caractéristiques du texte explicatif]
Un texte explicatif se reconnaît à quatre traits :
\begin{enumerate}
\item une \cle{visée informative} : transmettre un savoir vérifiable ;
\item une \cle{organisation logique} : du général au particulier, de la cause à la conséquence, du problème à la solution ;
\item des \cle{connecteurs logiques} explicites : \emph{premièrement, ensuite, en outre, par conséquent, en effet, c'est pourquoi} ;
\item un \cle{vocabulaire précis et dénotatif} : chaque mot désigne une réalité objective (sa \emph{dénotation}), sans connotation affective.
\end{enumerate}
\end{propriete}

\begin{exemplebox}[Connecteurs logiques dans une phrase de rapport]
« Les fondations ont été coulées \emph{premièrement} ; \emph{ensuite}, le coffrage a été retiré après vingt-huit jours de cure ; \emph{par conséquent}, la structure peut désormais supporter la charge prévue. »

Chaque connecteur indique au lecteur la relation logique entre les étapes : ordre chronologique puis relation de conséquence.
\end{exemplebox}

\begin{attention}[Dénotation contre connotation]
Le mot \emph{bicoque} et le mot \emph{habitation précaire} désignent la même réalité, mais le premier transporte un jugement méprisant (connotation), le second reste descriptif (dénotation). Dans un rapport, on écrit « habitation précaire », jamais « bicoque ». De même, on écrit « précipitations abondantes : \SI{180}{mm} en 24 heures » et non « une pluie terrible ».
\end{attention}

\courssub{L'analyse des tonalités : ce que le rapport doit exclure}

\textbf{Observer avant de théoriser}\par\medskip

Lisons deux descriptions du même incident sur un chantier :

\begin{itemize}
\item Version A : « Drame évité de justesse ! Un pauvre ouvrier a failli être écrasé par une pelleteuse folle. Heureusement, le chauffeur, plus malin que la machine, a freiné à temps. »
\item Version B : « Le 12 mars à 10 h 15, un engin de chantier s'est déplacé en dehors de sa zone de circulation. L'opérateur a immobilisé l'engin. Aucun blessé n'est à déplorer. »
\end{itemize}

La version A mélange deux \cle{tonalités} interdites dans un rapport : le \emph{dramatique} (« drame », « pauvre ouvrier », « écrasé ») et le \emph{comique} (« pelleteuse folle », « plus malin que la machine »). La version B, elle, adopte la \cle{tonalité neutre} : faits datés, localisés, mesurables.

\begin{propriete}[La tonalité neutre]
La \cle{tonalité neutre} (ou objective, ou scientifique) s'obtient par :
\begin{enumerate}
\item la suppression des \cle{marqueurs énonciatifs} (\emph{je, tu, nous}) dans le corps du texte ;
\item l'usage du \cle{présent de vérité générale} (« l'eau bout à 100 °C ») ou du passé composé / passé simple factuel pour les constats datés ;
\item un \cle{vocabulaire spécialisé dénotatif} (termes techniques du métier) ;
\item l'absence d'adjectifs subjectifs, d'exclamations, d'ironie et de familiarité.
\end{enumerate}
\end{propriete}

\textbf{Pourquoi exclure le dramatique et le comique ?}\par\medskip

\begin{itemize}
\item Le \cle{dramatique} (pathos, urgence, exclamation) fait appel aux émotions du lecteur. Or un décideur qui lit un rapport doit juger des faits, non partager une émotion. « Il est urgent d'agir ! » devient « une intervention est nécessaire dans un délai de quinze jours ».
\item Le \cle{comique} (ironie, familiarité, plaisanterie) installe une complicité entre l'auteur et le lecteur. Cette complicité détruit la distance professionnelle et peut être mal interprétée par un destinataire officiel. « Le serveur a encore planté » devient « le serveur a connu une interruption de service de deux heures ».
\end{itemize}

\begin{aretenir}[À retenir]
Dans un rapport technique, on ne cherche ni à faire pleurer, ni à faire rire, ni à convaincre par l'émotion : on cherche à \emph{informer de manière vérifiable}. Toute phrase doit pouvoir être précédée mentalement de la formule « il est constaté que... ».
\end{aretenir}

\courssub{Les marqueurs linguistiques de l'objectivité}

L'objectivité n'est pas seulement une intention : elle se fabrique avec des outils grammaticaux précis.

\textbf{Les quatre outils de l'impersonnel}\par\medskip

\begin{enumerate}
\item \textbf{Le passif réflexif (ou pronominal à valeur passive)} : « \emph{on observe que} les sondages révèlent une couche d'argile », « \emph{il se constate} une usure des roulements ». Le sujet humain disparaît derrière l'observation.
\item \textbf{Les tournures impersonnelles} : « \emph{il est nécessaire de} renforcer la poutre », « \emph{il convient de} vérifier l'étanchéité », « \emph{il apparaît que} les coûts dépassent le devis ».
\item \textbf{Les noms d'action} (noms déverbaux) : au lieu de « nous avons analysé les échantillons », on écrit « \emph{l'analyse} des échantillons a été effectuée ». Autres exemples : \emph{la réalisation, la mise en place, le contrôle, l'évaluation, le terrassement}. Le nom d'action gomme l'agent et met en avant l'opération elle-même.
\item \textbf{L'absence de « je »} : le pronom \emph{je} est banni du corps du rapport. Il n'est toléré que dans la préface, les remerciements ou la lettre de transmission, parties qui ne relèvent pas du contenu technique.
\end{enumerate}

\begin{exemplebox}[Du subjectif à l'objectif]
\begin{center}
\begin{tabularx}{\linewidth}{|X|X|}
\hline
\rowcolor{popblueL} \textbf{Phrase subjective (à éviter)} & \textbf{Phrase objective (attendue)} \\
\hline
J'ai vu que le mur était fissuré. & Il est constaté une fissure verticale de \SI{2}{mm} sur le mur pignon. \\
\hline
On a bien avancé, c'était super. & L'avancement des travaux atteint 75 \% du planning prévisionnel. \\
\hline
Le ciment qu'ils ont livré, c'est de la camelote. & Le ciment livré ne répond pas à la norme spécifiée (résistance insuffisante). \\
\hline
\end{tabularx}
\end{center}
\end{exemplebox}

\begin{attention}[Erreur fréquente : l'adjectif subjectif à la place de la mesure]
N'écrivez jamais « un \emph{grand} réservoir », « une \emph{belle} route », « une \emph{terrible} inondation ». Remplacez chaque adjectif qualificatif par une \cle{donnée chiffrée} : « un réservoir de \SI{500}{m^3} », « une route bitumée de \SI{12}{km} », « une inondation ayant submergé \SI{300}{ha} ». La mesure est le meilleur garant de l'objectivité.
\end{attention}

\begin{savaistu}[Le « Nous » de modestie dans l'administration camerounaise]
Dans l'administration camerounaise, l'usage du « \emph{Nous} » de modestie (« Nous avons l'honneur de... », « Nous soussigné... ») reste la norme dans les lettres de transmission et les en-têtes officiels. Attention cependant : ce « Nous » protocolaire appartient à la correspondance, pas au corps technique du rapport, qui exige l'impersonnel. Savoir passer d'un registre à l'autre est une compétence professionnelle recherchée.
\end{savaistu}

\begin{popfigure}
\begin{center}
\begin{tikzpicture}[font=\small]
% Entonnoir
\draw[thick, popblue, fill=popblueL] (-4,3) -- (4,3) -- (1.2,0.4) -- (0.6,-1.2) -- (-0.6,-1.2) -- (-1.2,0.4) -- cycle;
% Entrées
\node[align=left, anchor=west] at (-8.2,3.4) {\textbf{Entrées (subjectif)}};
\node[align=left, anchor=west] at (-8.2,2.6) {Impressions personnelles};
\node[align=left, anchor=west] at (-8.2,2.0) {Émotions, adjectifs subjectifs};
\node[align=left, anchor=west] at (-8.2,1.4) {Jargon, familiarité, « je »};
\draw[-{Stealth}, poporange, thick] (-4.2,2.4) -- (-3.2,2.4);
% Filtre
\node[align=center, popdark] at (0,1.7) {\textbf{Filtre de l'objectivité}};
\node[align=center] at (0,1.1) {verbes d'état, noms d'action,};
\node[align=center] at (0,0.7) {passif, tournures impersonnelles,};
\node[align=center] at (0,0.3) {connecteurs logiques, mesures};
% Sortie
\draw[-{Stealth}, popgreen, thick] (0,-1.4) -- (0,-2.2);
\node[align=center, popgreen!60!black] at (0,-2.7) {\textbf{Sortie : texte explicatif standardisé}};
\node[align=center] at (0,-3.2) {constats datés, chiffrés, vérifiables};
\end{tikzpicture}
\end{center}
\legende{L'entonnoir de l'objectivité : les matériaux bruts de l'observation (impressions, émotions, jargon) traversent un filtre linguistique (passif, noms d'action, impersonnel, connecteurs) pour produire un texte explicatif standardisé, conforme aux exigences du rapport technique.}
\end{popfigure}

\courssub{Du texte iconique au texte explicatif : la photographie de chantier}

\textbf{Observer une image comme un technicien}\par\medskip

Un rapport s'accompagne souvent de documents iconiques : photographies, plans, schémas. Mais une image ne parle pas toute seule : le rédacteur doit la \emph{traduire} en texte explicatif. Prenons l'exemple d'une photographie du chantier du \cle{barrage de Lom Pangar} (région de l'Est, Cameroun), ouvrage de régulation hydraulique sur la rivière Sanaga.

Face à cette photo, l'œil non entraîné voit « des grues et des camions ». L'œil du technicien, lui, découpe l'image en \emph{zones fonctionnelles} et nomme chaque élément avec précision.

\begin{methode}[Décrire une image technique en quatre étapes]
\begin{enumerate}
\item \textbf{Identifier le sujet global} : de quoi s'agit-il ? (« Photographie aérienne du chantier du barrage de Lom Pangar, prise en juin 2015. »)
\item \textbf{Découper en zones spatiales et fonctionnelles} : premier plan, arrière-plan, centre ; zone de fondations, zone de déversoir, zone de stockage.
\item \textbf{Nommer précisément} chaque élément avec le vocabulaire technique : \emph{fondations, déversoir, coffrage, grue à tour, centrale à béton, piste de chantier}.
\item \textbf{Articuler par des connecteurs spatiaux} : \emph{au premier plan, à l'arrière-plan, au centre, à droite, en contrebas, au-dessus de}.
\end{enumerate}
\end{methode}

\begin{popfigure}
\begin{center}
\begin{tikzpicture}[font=\small]
% Calque 1 : photo brute
\draw[thick, popdark, fill=popcream] (0,0) rectangle (3.4,2.4);
\draw[popmuted] (0,0.8) -- (3.4,0.8);
\fill[popmuted!60] (0.4,0.1) rectangle (1.2,0.7);
\fill[popmuted!60] (2.2,0.2) rectangle (3.0,0.7);
\draw[popmuted] (1.7,0.8) -- (1.7,2.1);
\node[below, align=center] at (1.7,-0.15) {\textbf{1. Photo brute}\\ (perception globale)};
% Calque 2 : zones identifiées
\draw[thick, popdark, fill=popcream] (4.6,0) rectangle (8.0,2.4);
\draw[popblue, thick, dashed] (4.8,0.2) rectangle (6.0,0.9);
\draw[poporange, thick, dashed] (6.8,0.2) rectangle (7.8,0.9);
\draw[popgreen, thick, dashed] (6.2,1.0) rectangle (6.9,2.2);
\node[popblue] at (5.4,0.55) {\tiny fondations};
\node[poporange] at (7.3,0.55) {\tiny déversoir};
\node[popgreen!60!black] at (6.55,1.6) {\tiny grue};
\node[below, align=center] at (6.3,-0.15) {\textbf{2. Zones techniques}\\ identifiées};
% Calque 3 : légende structurée
\draw[thick, popdark, fill=popblueL] (9.2,0) rectangle (13.2,2.4);
\node[align=left, anchor=north west, font=\scriptsize] at (9.35,2.3) {\textbf{3. Légende technique :}\\ Au premier plan, les fondations\\ en béton armé ; à droite, le\\ déversoir ; au centre, une grue\\ à tour assure le levage.};
\node[below, align=center] at (11.2,-0.15) {\textbf{3. Texte explicatif}\\ structuré};
% Flèches
\draw[-{Stealth}, very thick, poporange] (3.5,1.2) -- (4.5,1.2);
\draw[-{Stealth}, very thick, poporange] (8.1,1.2) -- (9.1,1.2);
\end{tikzpicture}
\end{center}
\legende{Les trois calques du traitement d'un texte iconique (photo du chantier de Lom Pangar) : 1. perception brute ; 2. découpage en zones fonctionnelles (fondations, déversoir, engins) ; 3. rédaction d'une légende technique articulée par des connecteurs spatiaux.}
\end{popfigure}

\begin{exemplebox}[Légende technique rédigée]
« \emph{Figure 1 : Chantier du barrage de Lom Pangar (Sanaga, région de l'Est). Au premier plan, on observe les fondations en béton armé de l'ouvrage principal ; au centre, une grue à tour assure le levage des coffrages ; à l'arrière-plan, la centrale à béton alimente le chantier par bennes roulantes. L'ensemble des travaux de terrassement est achevé ; la phase de bétonnage est en cours.} »

Remarquez : aucun « je », aucun adjectif subjectif, des connecteurs spatiaux, des noms d'action (\emph{levage, terrassement, bétonnage}) et des verbes d'état ou passifs.
\end{exemplebox}

\courssub{Exercice de réécriture : du récit subjectif au rapport objectif}

\begin{exoresolu}[Transformer un récit subjectif en constat de rapport]
\textbf{Énoncé.} Un stagiaire écrit dans son carnet : « J'ai vu les ouvriers travailler dur sous le soleil, c'était vraiment impressionnant, ils n'ont pas traîné ! » Réécrivez cette phrase en style de rapport objectif.
\tcblower
\textbf{Solution détaillée.}

Étape 1 : repérer les marques de subjectivité à supprimer : « j'ai vu » (marqueur énonciatif), « travailler dur » (jugement), « impressionnant » (émotion), « ils n'ont pas traîné » (familiarité).

Étape 2 : remplacer chaque jugement par un fait observable ou mesurable : quelle tâche ? quelle durée ? quelles conditions ?

Étape 3 : reformuler à l'impersonnel avec un nom d'action.

\textbf{Version corrigée} : « Les ouvriers ont effectué les travaux de terrassement dans des conditions climatiques difficiles (température supérieure à 35 °C). Le chantier a progressé conformément au planning prévisionnel. »

On peut aussi écrire : « Il est constaté que le terrassement a été réalisé dans les délais prévus, malgré des températures élevées. »
\end{exoresolu}

\begin{exercice}[Réécritures]
Réécrivez les phrases suivantes en style de rapport objectif :
\begin{enumerate}
\item « Franchement, le fournisseur nous a livré du matériel pourri. »
\item « Hier, j'étais sur le site et boom ! une conduite a explosé, c'était la panique. »
\item « La nouvelle salle informatique est magnifique, les élèves l'adorent. »
\end{enumerate}
\end{exercice}

\begin{corrige}[Exercice : Réécritures]
\begin{enumerate}
\item « Le matériel livré par le fournisseur présente des non-conformités par rapport aux spécifications du bon de commande (liste en annexe). » — On supprime « franchement » (familier) et « pourri » (subjectif), remplacés par un constat vérifiable.
\item « Le 14 mai, une conduite d'eau a cédé au niveau du secteur B, provoquant une interruption de l'alimentation. Les équipes ont procédé à l'évacuation de la zone, puis à la réparation. » — Faits datés, localisés, sans pathos.
\item « La salle informatique rénovée comprend 25 postes équipés ; le taux de fréquentation des élèves a augmenté de 40 \% depuis sa mise en service. » — L'adjectif « magnifique » est remplacé par des données mesurables.
\end{enumerate}
\end{corrige}

\courssub{Remédiation : corriger les écarts de registre et les ambiguïtés}

\textbf{Les écarts de registre}\par\medskip

Le rapport exige un \cle{registre soutenu} (ou courant soutenu). Les mots familiers doivent être systématiquement remplacés :

\begin{center}
\begin{tabularx}{\linewidth}{|l|X|}
\hline
\rowcolor{popblueL} \textbf{Registre familier (à bannir)} & \textbf{Registre soutenu (attendu)} \\
\hline
le boulot, le taf & le travail, les travaux, la tâche \\
\hline
le bagnole, la caisse & le véhicule, la voiture de service \\
\hline
bosser & travailler, exécuter \\
\hline
le fric, l'oseille & les fonds, le budget, les ressources financières \\
\hline
ça a planté & une panne est survenue, le système a cessé de fonctionner \\
\hline
les gars du chantier & le personnel de chantier, les ouvriers \\
\hline
\end{tabularx}
\end{center}

\textbf{Les ambiguïtés syntaxiques}\par\medskip

Une phrase mal construite peut avoir deux sens : c'est une \cle{ambiguïté}, faute grave dans un rapport où chaque phrase doit être interprétée d'une seule façon.

\begin{exemplebox}[Ambiguïtés à corriger]
\begin{itemize}
\item « Le chef de chantier a dit au grutier qu'il devait partir. » — Qui doit partir ? Correction : « Le chef de chantier a ordonné au grutier de quitter le poste. »
\item « J'ai observé la fissure avec des jumelles. » — La fissure a-t-elle des jumelles ? Correction : « À l'aide de jumelles, il a été observé une fissure sur le parement amont. »
\item « L'engin a heurté la canalisation qui était en mauvais état. » — L'engin ou la canalisation ? Correction : « L'engin a heurté la canalisation, laquelle était déjà dégradée. »
\end{itemize}
\end{exemplebox}

\begin{attention}[Pièges de rédaction à vérifier avant de rendre le rapport]
\begin{enumerate}
\item Aucun \emph{je} dans le corps du texte (réservé à la préface et aux remerciements).
\item Aucun adjectif subjectif sans mesure chiffrée à la place.
\item Aucun mot familier, aucune exclamation, aucune ironie.
\item Chaque pronom (\emph{il, qui, laquelle}) renvoie à un seul nom possible.
\item Chaque étape est introduite par un connecteur logique ou spatial explicite.
\end{enumerate}
\end{attention}

\begin{aretenir}[L'essentiel de la section]
Le rapport technique adopte le modèle du \cle{texte explicatif} : visée informative, organisation logique, connecteurs explicites, vocabulaire dénotatif. Sa \cle{tonalité neutre} exclut le dramatique et le comique. L'objectivité se construit grammaticalement : passif réflexif, tournures impersonnelles, noms d'action, suppression du « je », mesures chiffrées. Enfin, un document iconique (photographie de chantier) se traduit en texte par un découpage en zones, un vocabulaire technique et des connecteurs spatiaux. Ces compétences stylistiques vous permettront, dans la section suivante, d'élaborer le plan détaillé de votre propre rapport.
\end{aretenir}

\coursec{Élaboration du plan détaillé : Méthodologie pas à pas}

Dans la section précédente, nous avons appris à donner à notre rapport le bon \emph{registre} : un style explicatif, objectif, dépourvu de tonalités dramatiques ou comiques. Mais un beau style ne sauve pas un rapport mal construit. Avant d'écrire la moindre phrase rédigée, le technicien — comme l'élève de Première technique — doit d'abord \cle{bâtir la charpente} de son texte : c'est le \cle{plan détaillé}. Cette section vous donne une méthode en six phases, applicable à tout rapport de stage, de TP ou d'activité.

\courssub{Qu'est-ce qu'un plan détaillé ?}

\begin{definition}[Plan détaillé]
Le \cle{plan détaillé} est la représentation écrite et hiérarchisée de l'organisation future d'un rapport. Il présente, sous forme de \cle{titres numérotés} (numérotation décimale : 1. / 1.1. / 1.1.1.), l'ensemble des parties, sous-parties et paragraphes du document, dans l'ordre exact où ils seront rédigés.
\end{definition}

L'intuition est simple : construire un rapport sans plan, c'est comme monter une machine sans schéma d'assemblage. On peut y arriver, mais on risque les pièces oubliées, les doublons et les incohérences. Le plan détaillé est le \emph{schéma d'assemblage} de votre texte.

\begin{propriete}[Propriété fondamentale du plan détaillé]
Un plan détaillé valide permet de rédiger le rapport en \og remplissant les cases \fg{} : chaque sous-partie numérotée (par exemple 2.3.1.) correspond à \textbf{un paragraphe unique développé}. Si un item du plan ne peut pas donner lieu à au moins un paragraphe complet, il est trop fin ; si un item exige plusieurs pages sans subdivision, il est trop large et doit être éclaté.
\end{propriete}

\begin{exemplebox}[Le bon dosage]
Un plan détaillé de rapport de stage (rédigé sur une période de 13 séances de préparation) doit comporter environ \textbf{10 à 15 items numérotés} (parties + sous-parties confondues) pour une rédaction finale de 10 à 15 pages. En dessous de 10 items, le plan est trop pauvre ; au-dessus de 15, il devient illisible.
\end{exemplebox}

\courssub{Phase 1 : Exploitation du sujet et de la commande}

Tout rapport naît d'une \cle{commande}. Avant de chercher des idées, il faut interroger la situation de production. Quatre questions obligatoires :

\begin{enumerate}
\item \textbf{Quels sont les mots-clés du sujet ?} Soulignez-les. Dans \og Maintenance préventive sur groupe électrogène CATERPILLAR C15 \fg{}, les mots-clés sont : \emph{maintenance préventive} (et non corrective !), \emph{groupe électrogène}, \emph{CATERPILLAR C15} (modèle précis, donc fiche technique à exploiter).
\item \textbf{Qui est le commanditaire ?} Le chef de service, l'entreprise d'accueil, l'enseignant. C'est lui qui fixe les attentes.
\item \textbf{Qui est le destinataire ?} Le jury du Probatoire ou du Baccalauréat technique, le directeur technique, un client. Le destinataire détermine le niveau de technicité.
\item \textbf{Quel est le délai et le volume attendu ?} Un rapport de 5 pages et un rapport de 15 pages n'ont pas le même plan.
\end{enumerate}

\begin{attention}[Piège de la phase 1]
Ne confondez jamais le \emph{sujet} (ce qu'on vous demande de traiter) avec le \emph{titre} (ce que vous écrirez en couverture). Le sujet impose des limites : si la commande porte sur la maintenance \textbf{préventive}, un long développement sur les pannes accidentelles (maintenance corrective) est hors sujet, sauf en sous-partie clairement justifiée.
\end{attention}

\courssub{Phase 2 : Collecte et tri de l'information}

Le plan ne s'invente pas : il \emph{émerge} des documents. Rassemblez toute la matière brute :

\begin{itemize}
\item les \cle{notes de terrain} (carnet de bord du stagiaire, dates, heures, observations) ;
\item les \cle{documents fournis} (fiches techniques du constructeur, consignes de sécurité, organigramme de l'entreprise) ;
\item les \cle{données chiffrées} (heures de fonctionnement, tensions mesurées, consommation de gasoil, fréquence des interventions) ;
\item les \cle{photos et textes iconiques} (état avant/après intervention, schémas, plaques signalétiques).
\end{itemize}

Puis effectuez le \cle{tri} : chaque document est classé en trois catégories — \textbf{essentiel} (ira dans le corps du rapport), \textbf{annexe} (document utile mais trop volumineux : fiche technique complète, photo), \textbf{rebut} (hors sujet ou redondant). Ce tri prépare déjà la structure : ce qui est essentiel nourrit les parties, ce qui est annexe nourrit la dernière section du rapport.

\courssub{Phase 3 : Regroupement thématique — la carte heuristique}

Les informations triées doivent maintenant être regroupées par \emph{affinité}. L'outil idéal est la \cle{carte heuristique} (ou mind-mapping) : on place le sujet au centre, on fait jaillir des branches principales (les grandes parties), puis des sous-branches (les sous-parties).

\begin{methode}[Construire sa carte heuristique]
\begin{enumerate}
\item Écrire le sujet au centre d'une feuille et l'entourer.
\item Noter sur des branches toutes les idées et données issues de la phase 2, sans les juger.
\item Regrouper les branches proches et leur donner un nom commun : ce sont les \textbf{grandes parties}.
\item Limiter à \textbf{2 ou 3 grandes parties} pour le corps du rapport (hors introduction et conclusion).
\item Vérifier que chaque branche principale contient au moins deux sous-branches.
\end{enumerate}
\end{methode}

\begin{popfigure}
\begin{center}
\begin{tikzpicture}[mindmap, grow cyclic, every node/.style={concept}, concept color=popblueL,
root concept/.append style={concept color=popblue, font=\small\bfseries, text=white},
level 1/.append style={level distance=3.4cm, sibling angle=90, font=\small\bfseries},
level 2/.append style={level distance=2.4cm, sibling angle=45, font=\scriptsize},
level 1 concept/.append style={concept color=poporangeL},
text=popdark]
\node[root concept]{Rapport de stage : maintenance du groupe CATERPILLAR C15}
  child[concept color=popgreenL]{ node{Contexte}
    child{ node{Entreprise d'accueil} }
    child{ node{Fiche technique du C15} }
  }
  child[concept color=poporangeL]{ node{Op\'erations}
    child{ node{Vidange et filtres} }
    child{ node{Contr\^ole \'electrique} }
    child{ node{Essais de charge} }
  }
  child[concept color=poppurpleL]{ node{R\'esultats}
    child{ node{Donn\'ees chiffr\'ees} }
    child{ node{Photos avant/apr\`es} }
  }
  child[concept color=poppinkL]{ node{Recommandations}
    child{ node{Planning pr\'eventif} }
    child{ node{Pi\`eces \`a stocker} }
  };
\end{tikzpicture}
\end{center}
\legende{Carte heuristique du rapport de stage : du sujet central jaillissent quatre branches principales (Contexte, Opérations, Résultats, Recommandations), chacune subdivisée en sous-branches qui deviendront les sous-parties du plan.}
\end{popfigure}

\begin{attention}[L'erreur du plan « en entonnoir »]
Erreur très fréquente au Probatoire : faire un plan \emph{en entonnoir} du type \og 1. Généralités — 2. Mon stage — 3. Conclusion \fg{}. Ce plan ne dit rien et ne guide aucune rédaction. Il faut un plan \textbf{analytique}, où chaque titre annonce un contenu précis : \og 1. Contexte de l'intervention — 2. Problème technique rencontré — 3. Solutions apportées et résultats \fg{}. Le lecteur doit comprendre la logique du rapport en lisant seulement le sommaire.
\end{attention}

\courssub{Phase 4 : Hiérarchisation et numérotation décimale}

La carte heuristique donne des \emph{groupes} ; le plan exige un \emph{ordre}. On transforme les branches en titres numérotés selon la \cle{numérotation décimale} :

\begin{itemize}
\item \textbf{Niveau 1} (parties) : 1. / 2. / 3.
\item \textbf{Niveau 2} (sous-parties) : 1.1. / 1.2. / 2.1.
\item \textbf{Niveau 3} (paragraphes) : 1.1.1. / 1.1.2.
\end{itemize}

L'ordonnancement suit une logique : \emph{chronologique} (avant / pendant / après l'intervention), \emph{analytique} (constat / causes / solutions) ou \emph{thématique} (du contexte vers les résultats). Pour un rapport technique, la logique analytique est la plus valorisée.

\begin{savaistu}[Une norme internationale]
La numérotation décimale (1.1.1.) est la norme \textbf{ISO 2145}, utilisée dans le monde entier pour les documents techniques. Elle est adoptée par l'Université de Yaoundé et les grandes écoles camerounaises (ENSP, ENSET) pour la présentation des mémoires. La maîtriser dès la Première technique, c'est acquérir un réflexe professionnel durable.
\end{savaistu}

\courssub{Phase 5 : Rédaction des titres de parties et sous-parties}

Un bon titre est \cle{nominal} (sans verbe conjugué : \og Contrôle du circuit électrique \fg{} et non \og Nous avons contrôlé le circuit électrique \fg{}). Il obéit à la règle des \og 3 P \fg{}.

\begin{methode}[La règle des 3 P pour les titres]
\begin{enumerate}
\item \textbf{Parallélisme} : tous les titres d'un même niveau ont la même structure grammaticale. Si 2.1. est \og Vidange et remplacement des filtres \fg{} (nom + complément), 2.2. doit suivre le même moule : \og Contrôle du circuit électrique \fg{}, et non \og Comment contrôler le circuit électrique ? \fg{}.
\item \textbf{Précision} : employer le vocabulaire technique exact (vidange, couple de serrage, essai de charge) et non des termes vagues (\og les travaux \fg{}, \og divers \fg{}).
\item \textbf{Proportionnalité} : les parties doivent avoir un volume comparable ; une partie de 6 pages à côté d'une partie d'une demi-page signale un déséquilibre du plan.
\end{enumerate}
\end{methode}

\courssub{Phase 6 : Validation du plan avant rédaction}

Avant de rédiger, on soumet le plan à un triple contrôle :

\begin{center}
\begin{tabularx}{\linewidth}{|l|X|X|}
\hline
\rowcolor{popblueL} \textbf{Critère} & \textbf{Question de contrôle} & \textbf{Si le test échoue} \\
\hline
Cohérence & Chaque sous-partie relève-t-elle bien de sa partie ? & Déplacer l'item mal placé. \\
\hline
Exhaustivité & Tous les mots-clés du sujet sont-ils traités ? & Ajouter la partie manquante. \\
\hline
Proportionnalité & Les parties ont-elles un poids comparable ? & Fusionner ou subdiviser. \\
\hline
\end{tabularx}
\end{center}

Un plan qui passe ces trois tests est \emph{validé} : la rédaction peut commencer, case par case.

\courssub{Application guidée : rapport de stage « Maintenance préventive sur groupe électrogène CATERPILLAR C15 »}

\textbf{Situation.} Vous avez effectué un stage dans une entreprise de Bafoussam. Vous disposez de : carnets de bord (dates, heures de fonctionnement), fiches techniques du constructeur, photos avant/après intervention. Appliquons les six phases.

\textbf{Phase 1 — Commande.} Mots-clés : maintenance \emph{préventive}, groupe électrogène, C15. Commanditaire : l'entreprise d'accueil. Destinataire : jury du Baccalauréat technique. Volume : 10-15 pages.

\textbf{Phases 2 et 3 — Tri et regroupement.} Les carnets de bord et fiches techniques alimentent le \emph{contexte} et les \emph{opérations} ; les mesures et photos alimentent les \emph{résultats} ; les remarques du chef de service alimentent les \emph{recommandations} (voir la carte heuristique ci-dessus). La fiche technique complète et les photos grand format partent en \emph{annexes}.

\textbf{Phases 4 et 5 — Plan numéroté.}

\begin{popfigure}
\begin{center}
\begin{tikzpicture}[font=\small, text=popdark]
\node[draw=popblue, fill=popblueL, rounded corners, text width=0.9\linewidth, align=left] (p1) at (0,0) {\textbf{1. Introduction} : objet du rapport, méthode, annonce du plan};
\node[draw=popgreen, fill=popgreenL, rounded corners, text width=0.9\linewidth, align=left, below=0.15cm of p1] (p2) {\textbf{2. Pr\'esentation de la structure d'accueil} \quad 2.1. Activit\'es de l'entreprise \quad 2.2. Service maintenance};
\node[draw=poporange, fill=poporangeL, rounded corners, text width=0.9\linewidth, align=left, below=0.15cm of p2] (p3) {\textbf{3. Activit\'es men\'ees} \quad 3.1. Maintenance pr\'eventive (vidange, filtres, contr\^oles) \quad 3.2. Maintenance corrective (incident constat\'e)};
\node[draw=poppurple, fill=poppurpleL, rounded corners, text width=0.9\linewidth, align=left, below=0.15cm of p3] (p4) {\textbf{4. Analyse des performances du groupe C15} \quad 4.1. Donn\'ees avant intervention \quad 4.2. Donn\'ees apr\`es intervention};
\node[draw=poppink, fill=poppinkL, rounded corners, text width=0.9\linewidth, align=left, below=0.15cm of p4] (p5) {\textbf{5. Conclusion} : bilan, comp\'etences acquises, recommandations};
\node[draw=popmuted, fill=popcream, rounded corners, text width=0.9\linewidth, align=left, below=0.15cm of p5] (p6) {\textbf{6. Annexes} : fiche technique CATERPILLAR C15, photos, carnet de bord};
\end{tikzpicture}
\end{center}
\legende{Modèle de plan détaillé numéroté : 6 grandes sections, 13 items numérotés au total, conforme à la règle des 10 à 15 items pour un rapport de 10 à 15 pages.}
\end{popfigure}

\begin{exoresolu}[Vérifier la validité du plan obtenu]
Énoncé. Le plan ci-dessus comporte-t-il le bon nombre d'items ? Respecte-t-il la règle des 3 P ? Corriger le titre fautif suivant proposé par un élève : \og 3.2. Comment j'ai réparé une panne \fg{}.
\tcblower
Solution détaillée.
\textbf{Nombre d'items} : 6 sections de niveau 1, dont 3 subdivisées (2.1., 2.2., 3.1., 3.2., 4.1., 4.2.), soit $6 + 6 = 12$ items numérotés : on est bien dans la fourchette de 10 à 15 items.
\textbf{Règle des 3 P} : les titres sont nominaux et parallèles (\og Maintenance préventive \fg{} / \og Maintenance corrective \fg{} ; \og Données avant intervention \fg{} / \og Données après intervention \fg{}) : parallélisme respecté. Le vocabulaire est technique (vidange, filtres, performances) : précision respectée. Les parties 2, 3 et 4 ont chacune deux sous-parties : proportionnalité respectée.
\textbf{Correction} : \og Comment j'ai réparé une panne \fg{} viole les trois règles (phrase interrogative à la première personne, non nominale, non parallèle, imprécise). Titre corrigé : \og 3.2. Maintenance corrective : remplacement du relais de démarrage \fg{} — nominal, technique, parallèle à 3.1.
\end{exoresolu}

\begin{exercice}[À vous de planifier]
À partir de la situation suivante, élaborez un plan détaillé complet (10 à 15 items, numérotation décimale) : \og Rapport de TP : réglage du circuit de refroidissement d'un moteur Diesel en atelier scolaire \fg{}. Vous disposez de vos notes de TP, de la fiche technique du moteur et de trois photos. Vérifiez ensuite votre plan avec la grille des trois critères (cohérence, exhaustivité, proportionnalité).
\end{exercice}

\begin{corrige}[Exercice 1]
Éléments de correction. Phase 1 : mots-clés = TP, réglage, circuit de refroidissement, moteur Diesel ; destinataire = enseignant et jury. Plan attendu (modèle) : \textbf{1. Introduction} (objet du TP, matériel) ; \textbf{2. Présentation du circuit de refroidissement} (2.1. Rôle du circuit ; 2.2. Composants : pompe, radiateur, thermostat) ; \textbf{3. Travaux réalisés} (3.1. Contrôle visuel et niveau de liquide ; 3.2. Réglage et purge du circuit ; 3.3. Essai de fonctionnement) ; \textbf{4. Résultats et interprétation} (4.1. Températures relevées ; 4.2. Comparaison aux valeurs constructeur) ; \textbf{5. Conclusion} (bilan, difficultés, recommandations) ; \textbf{6. Annexes} (fiche technique, photos). Total : 13 items, titres nominaux et parallèles, parties équilibrées : plan valide. Toute proposition équivalente respectant les 3 P et la numérotation décimale est acceptée.
\end{corrige}

\begin{aretenir}[L'essentiel de la section]
\begin{itemize}
\item Le plan détaillé se construit en \textbf{six phases} : exploitation de la commande, collecte et tri, regroupement par carte heuristique, numérotation décimale, rédaction des titres, validation.
\item Les titres obéissent à la règle des \textbf{3 P} : Parallélisme, Précision, Proportionnalité.
\item Un bon plan compte \textbf{10 à 15 items} pour 10 à 15 pages ; chaque sous-partie = un paragraphe.
\item On évite le plan \og en entonnoir \fg{} ; on exige un plan \textbf{analytique} où le sommaire seul raconte la logique du rapport.
\end{itemize}
\end{aretenir}

\coursec{Amélioration du plan et du texte : Révision, cohérence et iconographie}

\courssub{Transition : du plan brut au plan validé}

Après l'élaboration du plan détaillé vue en section 4, vous disposez d'une architecture complète : introduction, développement en parties et sous-parties numérotées, conclusion. Ce plan est une \emph{première version}, un « brouillon structuré ». L'expérience professionnelle montre qu'un plan n'est jamais définitif dès le premier jet. La phase d'amélioration — révision, cohérence, intégration de l'iconographie — est celle qui transforme un squelette correct en un rapport technique professionnel, prêt à être lu par un décideur (chef de chantier, directeur technique, client).

\courssub{Critères d'auto-évaluation du plan détaillé}

Avant toute réécriture, il faut juger la solidité de l'architecture. Trois critères majeurs guident cette auto-évaluation.

\begin{definition}[Critères de validation d'un plan technique]
Un plan détaillé de rapport technique est validé s'il satisfait simultanément à :
\begin{enumerate}
    \item \textbf{Pertinence par rapport au sujet} : Chaque partie et sous-partie répond directement à la commande (objet du rapport). Aucune digression, aucun « hors-sujet » masqué.
    \item \textbf{Logique de l'enchaînement} : L'ordre des parties suit une progression rationnelle (chronologique, causale, analytique : description $\rightarrow$ analyse $\rightarrow$ proposition). Les transitions entre parties sont implicites dans la numérotation.
    \item \textbf{Équilibre des parties} : Les parties principales ont un poids comparable (volume prévisionnel, nombre de sous-parties). Une partie unique avec cinq sous-parties et une autre avec une seule sous-partie signalent un déséquilibre à corriger.
\end{enumerate}
\end{definition}

\begin{popfigure}
\begin{tikzpicture}[
    box/.style={draw=popdark, thick, rounded corners, fill=popcream, minimum height=1.1cm, text width=5.2cm, align=center, font=\small},
    check/.style={draw=popgreen, thick, fill=popgreen!20, circle, minimum size=0.45cm, font=\normalsize\bfseries},
    crit/.style={font=\small, text width=5.5cm, align=left},
    cat/.style={font=\bfseries\small, text width=2.5cm, align=right, text=popblue}
]
% En-tête
\node[draw=popblue, fill=popblueL, thick, rounded corners, minimum height=0.9cm, text width=14.5cm, align=center, font=\bfseries] (header) at (0,0) {GRILLE D'AUTO-ÉVALUATION DU PLAN DÉTAILLÉ (10 critères)};
% Colonnes
\node[cat] (c1) at (-7.5,-1.5) {\textbf{Structure}};
\node[cat] (c2) at (-7.5,-3.0) {\textbf{Contenu}};
\node[cat] (c3) at (-7.5,-4.5) {\textbf{Forme}};
% Lignes critères
% 1
\node[check] (ck1) at (-5.5,-1.5) {1};
\node[crit] at (-2.5,-1.5) {Numérotation cohérente (1 ; 1.1 ; 1.1.1) sans saut ni doublon.};
% 2
\node[check] (ck2) at (-5.5,-2.2) {2};
\node[crit] at (-2.5,-2.2) {Parallélisme des titres (même nature grammaticale : tous infinitifs ou tous nominaux).};
% 3
\node[check] (ck3) at (1.5,-1.5) {3};
\node[crit] at (4.5,-1.5) {Exhaustivité : tous les points de la commande sont couverts.};
% 4
\node[check] (ck4) at (1.5,-2.2) {4};
\node[crit] at (4.5,-2.2) {Hiérarchie respectée : pas de sous-partie orpheline (un 1.1 impose un 1.2).};
% 5
\node[check] (ck5) at (1.5,-2.9) {5};
\node[crit] at (4.5,-2.9) {Équilibre quantitatif : 3 à 5 parties principales, 2 à 4 sous-parties par partie.};
% 6
\node[check] (ck6) at (-5.5,-3.7) {6};
\node[crit] at (-2.5,-3.7) {Titres explicites : le titre seul résume le contenu de la section (pas de « Étude » seul).};
% 7
\node[check] (ck7) at (-5.5,-4.4) {7};
\node[crit] at (-2.5,-4.4) {Logique interne : ordre chronologique, spatial, ou déductif assumé et constant.};
% 8
\node[check] (ck8) at (1.5,-3.7) {8};
\node[crit] at (4.5,-3.7) {Transition implicite : le plan se lit comme un récit sans rupture brutale.};
% 9
\node[check] (ck9) at (1.5,-4.4) {9};
\node[crit] at (4.5,-4.4) {Introduction/Conclusion prévues : accroche, sujet, division / bilan, ouverture, recommandations.};
% 10
\node[check] (ck10) at (-5.5,-5.1) {10};
\node[crit] at (-2.5,-5.1) {Prévoir l'iconographie : emplacement des figures/tableaux noté dans le plan (ex: [Fig. 2]).};
% Cadre global
\draw[popline, thick, rounded corners] (-9.5, 0.6) rectangle (9.5, -5.6);
\end{tikzpicture}
\legende{Checklist TikZ : Grille d'auto-évaluation du plan détaillé (10 critères). Cochez chaque case après relecture.}
\end{popfigure}

\courssub{Techniques d'amélioration du plan : opérer sur la structure}

L'auto-évaluation révèle souvent des défauts structurels. Trois opérations chirurgicales permettent de les corriger sans tout réécrire.

\begin{methode}[Trois opérations pour rééquilibrer un plan]
\begin{enumerate}
    \item \textbf{Fusionner des sous-parties trop fines} : Si une partie contient 5 ou 6 sous-parties d'une demi-page chacune, regroupez-les par thématique commune en 2 ou 3 sous-parties plus denses. \textit{Exemple :} « 1.1 Vérification niveau huile », « 1.2 Vérification niveau liquide de refroidissement », « 1.3 Vérification pression pneus » $\rightarrow$ « 1.1 Contrôle des niveaux et pressions ».
    \item \textbf{Scinder des parties trop lourdes} : Si une partie unique concentre 60\% du rapport, créez une nouvelle partie principale. \textit{Exemple :} « 2. Analyse des pannes » (15 pages) $\rightarrow$ « 2. Analyse des pannes mécaniques » et « 3. Analyse des pannes électriques ».
    \item \textbf{Reformuler des titres flous} : Remplacez les titres vagues (« Étude », « Problèmes », « Généralités ») par des titres \cle{informatifs} qui énoncent le résultat ou l'action. \textit{Exemple :} « 3. Problèmes rencontrés » $\rightarrow$ « 3. Dysfonctionnements du système de pompage et mesures correctives ».
\end{enumerate}
\end{methode}

\begin{exemplebox}[Avant / Après : Rééquilibrage d'un plan de stage]
\textbf{Plan initial (déséquilibré)} \\
1. Présentation de l'entreprise (3 sous-parties) \\
2. Description du poste de travail (1 sous-partie) \\
3. Activités réalisées (8 sous-parties : trop fin, mélange tâches et observations) \\
4. Conclusion (1 sous-partie)

\textbf{Plan amélioré (équilibré, parallèle)} \\
1. Cadre du stage : entreprise et service d'affectation \\
2. Missions confiées et organisation du travail \\
3. Réalisations techniques : maintenance préventive et curative \\
4. Analyse critique : difficultés rencontrées et solutions apportées \\
5. Bilan personnel et perspectives professionnelles
\end{exemplebox}

\courssub{Intégration des textes iconiques (N°2 et N°3) : choisir, légender, appeler}

Un rapport technique sans iconographie est un rapport incomplet. Les textes iconiques N°2 et N°3 du programme (graphiques, schémas, photos) ne sont pas des illustrations décoratives : ce sont des \emph{preuves visuelles} ou des \emph{synthèses de données} insérées dans le développement.

\begin{propriete}[Règle d'or de l'iconographie technique]
Toute figure (tableau, graphe, photo, schéma, carte) doit respecter trois conditions indissociables :
\begin{enumerate}
    \item \textbf{Appel dans le texte} : La figure est citée \emph{avant} d'apparaître (ex: « La figure 3 montre l'évolution... »).
    \item \textbf{Autonomie} : La figure se comprend seule grâce à son \textbf{titre explicite}, sa \textbf{légende technique} (échelles, unités, symboles) et sa \textbf{source}.
    \item \textbf{Pertinence} : Elle illustre, prouve ou précise un point précis du développement. Pas de figure « pour faire joli ».
\end{enumerate}
\end{propriete}

\begin{center}
\begin{tabularx}{\linewidth}{|l|X|X|}\hline
\rowcolor{popblueL}
\textbf{Type d'iconique} & \textbf{Usage typique au développement} & \textbf{Exemple (Textes N°2 / N°3)} \\ \hline
\textbf{Courbe (pgfplots)} & Évolution temporelle d'un paramètre mesuré & Consommation carburant groupe G1 sur 7 jours (Texte N°2) \\ \hline
\textbf{Diagramme circulaire (camembert)} & Répartition en \% d'un tout (temps, budget, pannes) & Répartition des tâches de maintenance (Texte N°3) \\ \hline
\textbf{Organigramme (TikZ)} & Structure hiérarchique, flux de décision, processus & Organigramme du service maintenance (Texte N°3) \\ \hline
\textbf{Photo technique} & État des lieux, pièce défectueuse, montage réalisé & Photo du groupe électrogène G1, vue du filtre encrassé \\ \hline
\textbf{Schéma fonctionnel} & Principe de fonctionnement, liaisons entrée/sortie & Schéma hydraulique du circuit de refroidissement \\ \hline
\end{tabularx}
\end{center}

\courssub{Règles d'insertion normalisée : la fiche d'identité de la figure}

La norme professionnelle (et l'attente du Probatoire) impose un formalisme strict pour chaque figure insérée.

\begin{definition}[Composition normalisée d'une figure technique]
Une figure insérée dans un rapport comprend obligatoirement, de haut en bas :
\begin{enumerate}
    \item \textbf{Numéro d'ordre} : « Figure 1 », « Figure 2 »... (numérotation continue dans tout le rapport).
    \item \textbf{Titre explicite} : Phrase nominale descriptive (pas de verbe conjugué). Ex: « Évolution de la consommation journalière du groupe G1 ».
    \item \textbf{Corps de la figure} : Le graphique, le schéma, le tableau ou la photo lui-même.
    \item \textbf{Légende technique} : Explication des codes (couleurs, hachures, abréviations), unités des axes, échelles.
    \item \textbf{Source} : Origine des données (ex: « Source : Carnet de bord Semaine 4 », « Source : Relevé terrain 12/03/2026 »).
\end{enumerate}
Le renvoi dans le texte s'écrit : « (cf. Figure 3) » ou « comme le montre la Figure 3 ».
\end{definition}

\begin{popfigure}
\begin{tikzpicture}[
    frame/.style={draw=popdark, thick, rounded corners, fill=white},
    title/.style={font=\bfseries\small, text=popblue, anchor=north west},
    axis/.style={font=\tiny, text=popdark},
    legend/.style={font=\tiny, fill=white, draw=popline, rounded corners, inner sep=2pt},
    src/.style={font=\tiny\itshape, text=popmuted, anchor=south west}
]
% Cadre extérieur
\node[frame, minimum width=14cm, minimum height=8.5cm] (cadre) at (0,0) {};
% Titre Figure
\node[title] at (-6.8, 3.9) {Figure 3 : Évolution de la consommation journalière du groupe G1};
% Zone graphique (simulation pgfplots simplifiée en TikZ)
\begin{scope}[shift={(-6.5, -0.5)}]
% Axes
\draw[popdark, thick, ->] (0,0) -- (12,0) node[axis, right] {Jour (j)};
\draw[popdark, thick, ->] (0,0) -- (0,6) node[axis, above] {Conso. (L/j)};
% Graduations X
\foreach \x/\lab in {1/1, 2/2, 3/3, 4/4, 5/5, 6/6, 7/7} {
    \draw[popdark] (\x*1.5, -0.1) -- (\x*1.5, 0.1) node[axis, below] {\lab};
}
% Graduations Y
\foreach \y/\lab in {1/20, 2/25, 3/30, 4/35, 5/40} {
    \draw[popdark] (-0.1, \y*1.1) -- (0.1, \y*1.1) node[axis, left] {\lab};
}
% Courbe de consommation (points reliés)
\draw[poporange, very thick, smooth, tension=0.3] plot coordinates {
    (1.5, 2.4) (3.0, 2.8) (4.5, 3.5) (6.0, 3.0) (7.5, 3.2) (9.0, 3.6) (10.5, 2.9)
};
% Points
\foreach \px/\py in {1.5/2.4, 3.0/2.8, 4.5/3.5, 6.0/3.0, 7.5/3.2, 9.0/3.6, 10.5/2.9} {
    \fill[poporange] (\px,\py) circle (2.5pt);
}
% Légende courbe
\node[legend] at (6, 5.5) {Conso. journalière (L)};
% Seuil alerte
\draw[poppink, dashed, thick] (0, 3.64) -- (12, 3.64) node[axis, right, poppink] {Seuil alerte (33 L)};
\end{scope}
% Légende technique (sous le graphique)
\node[font=\small, text width=12.5cm, align=left, anchor=north west] at (-6.8, -2.8) {
    \textbf{Légende :} Courbe orange = consommation réelle mesurée au compteur journalier. \\
    Pointillés roses = seuil d'alerte contractuel (33 L/j). Variation jour 4 due à charge exceptionnelle (démarrage pompe secours).
};
% Source
\node[src] at (-6.8, -4.0) {Source : Carnet de bord Semaine 4 (Du 10/03 au 16/03/2026) - Relevé technicien Y. Mbarga};
\end{tikzpicture}
\legende{Mise en page type d'une figure technique (TikZ) : Titre numéroté, corps graphique (courbe pgfplots simulée), légende technique, source. Respect de la norme AFNOR.}
\end{popfigure}

\begin{attention}[Piège fréquent : l'iconique « décoratif »]
Ne jamais insérer une photo d'équipement, un logo d'entreprise ou un schéma générique \textbf{sans légende technique ni renvoi dans le texte}.
\begin{itemize}
    \item \textbf{Mauvais :} Photo d'un moteur en introduction, sans numéro, sans titre, non commentée. $\rightarrow$ Zéro point à l'examen.
    \item \textbf{Bon :} « La Figure 1 présente le groupe électrogène G1 (Caterpillar C15, 400 kVA) objet de l'étude. On y distingue le filtre à air (repère A) dont l'encrassement est analysé en 2.2. » (cf. Figure 1).
\end{itemize}
L'iconique doit \textbf{prouver} ou \textbf{illustrer} un propos du développement.
\end{attention}

\courssub{Rédaction de l'introduction et de la conclusion à partir du plan validé}

Une fois le plan stabilisé et l'iconographie positionnée (emplacements notés entre crochets dans le plan : [Fig. 2], [Tab. 1]), on rédige l'introduction et la conclusion. Elles s'écrivent \emph{après} le développement pour être parfaitement alignées sur le contenu réel.

\begin{methode}[Introduction en entonnoir (3 mouvements obligatoires)]
\begin{enumerate}
    \item \textbf{Accroche (Contexte général)} : 1 à 2 phrases situant le domaine technique, l'enjeu, la réglementation ou l'actualité. \textit{Ex: « La maintenance préventive des groupes électrogènes est critique pour la continuité de service des infrastructures hospitalières au Cameroun. »}
    \item \textbf{Sujet (Objet précis du rapport)} : Énoncé clair de la mission, de la période, du périmètre. \textit{Ex: « Ce rapport rend compte du stage de perfectionnement effectué du 10 au 28 mars 2026 au service maintenance de l'Hôpital Général de Douala, sur le suivi du groupe G1. »}
    \item \textbf{Division (Annonce du plan)} : Phrase unique listant les parties principales dans l'ordre. \textit{Ex: « Nous présenterons d'abord le cadre du stage (1), puis les opérations de maintenance réalisées (2), avant d'analyser les dysfonctionnements rencontrés (3) et de formuler des recommandations (4). »}
\end{enumerate}
\end{methode}

\begin{methode}[Conclusion en entonnoir inverse (3 mouvements obligatoires)]
\begin{enumerate}
    \item \textbf{Bilan (Synthèse des résultats)} : Réponse directe à l'objet du rapport, chiffrée si possible. \textit{Ex: « Le stage a permis de réaliser 12 interventions préventives et 3 curatives, ramenant la consommation du G1 sous le seuil d'alerte (31,2 L/j en moyenne). »}
    \item \textbf{Ouverture (Portée, limites)} : Perspectives techniques, limites de l'étude. \textit{Ex: « L'analyse vibrationnelle, non réalisée faute d'appareil, compléterait utilement ce diagnostic. »}
    \item \textbf{Recommandations (Actions concrètes)} : Propositions numérotées, adressées au destinataire. \textit{Ex: « 1) Planifier le changement du filtre à air tous les 500 h. 2) Acquérir un analyseur de vibrations. 3) Former l'équipe au relevé thermographique. »}
\end{enumerate}
\end{methode}

\courssub{Simulation de remédiation par les pairs : co-évaluation critériée}

La remédiation par les pairs (ou co-évaluation) est une pratique professionnelle standardisée : on fait relire son plan (puis son rapport) par un collègue « lecteur visé » avant remise finale. Cela permet de détecter les \cle{angles morts} (incompréhensions, jargon obscur, sauts logiques) que l'auteur ne voit plus par effet de familiarité.

\begin{definition}[Remédiation par les pairs (Co-évaluation)]
Processus structuré où deux étudiants échangent leur plan détaillé (ou rapport) et l'évaluent à l'aide d'une grille critériée commune. L'évaluateur joue le rôle du \textbf{lecteur visé} (le décideur technique) et non celui du correcteur scolaire. Il signale : « Ici je ne comprends pas », « Ce titre ne promet pas ce que le texte tient », « Il manque une figure pour appuyer ce chiffre ».
\end{definition}

\begin{center}
\begin{tabularx}{\linewidth}{|l|X|X|}\hline
\rowcolor{popblueL}
\textbf{Critère} & \textbf{Indicateurs de réussite (Oui / Partiel / Non)} & \textbf{Commentaire / Suggestion} \\ \hline
\textbf{Structure} & Numérotation continue, parallélisme titres, équilibre parties, logique enchaînement, intro/conclu annoncées. & \\ \hline
\textbf{Vocabulaire} & Termes techniques précis (pas de « truc », « chose »), emprunts maîtrisés, définitions si néologisme. & \\ \hline
\textbf{Syntaxe} & Phrases complexes correctes (subordination, passif, impersonnel), connecteurs logiques variés, pas de phrases « valise ». & \\ \hline
\textbf{Présentation} & Figures appelées avant insertion, titres/légendes/sources normalisés, renvois (cf. Fig. x), mise en page aérée. & \\ \hline
\end{tabularx}
\end{center}

\begin{exoresolu}[Simulation de remédiation : Échange de plans]
\textbf{Contexte :} Étudiant A (rédacteur) et Étudiant B (évaluateur) échangent leurs plans de rapport de stage « Maintenance système hydraulique presse plieuse ».

\textbf{Plan de l'Étudiant A (extrait) :} \\
1. Présentation boîte à huile \\
2. Vérifications faites \\
3. Problèmes vus \\
4. Ce qu'il faut faire

\textbf{Grille de l'Étudiant B (remplie) :}
\begin{itemize}
    \item \textbf{Structure :} \textcolor{poppink}{Non}. Titres non parallèles (mixte nominal/verbal), pas de numérotation hiérarchique (pas de 1.1), parties déséquilibrées. « Ce qu'il faut faire » n'est pas un titre technique.
    \item \textbf{Vocabulaire :} \textcolor{poporange}{Partiel}. « Boîte à huile » $\rightarrow$ Réservoir hydraulique. « Problèmes vus » $\rightarrow$ Dysfonctionnements détectés / Anomalies relevées.
    \item \textbf{Syntaxe :} \textcolor{popgreen}{Oui}. Titres courts, pas de phrase valise.
    \item \textbf{Présentation :} \textcolor{poppink}{Non}. Aucune figure prévue dans le plan.
\end{itemize}

\textbf{Remédiation proposée par B à A :} « Restructure en 4 parties nominales parallèles. Ajoute des sous-parties. Précise le vocabulaire technique. Note [Fig. 1: Schéma circuit] en 1.1 et [Fig. 2: Courbe pression] en 3.1. »
\textbf{Plan révisé par A (validé) :} \\
1. Caractéristiques du circuit hydraulique de la presse \\
2. Opérations de maintenance préventive réalisées \\
3. Anomalies détectées et diagnostic causal \\
4. Recommandations techniques et planning d'interventions
\tcblower
\textbf{Analyse :} La remédiation a transformé un plan « carnet de bord » en plan « rapport technique ». L'évaluateur a agi comme le lecteur visé (le chef d'atelier) : il a exigé la précision terminologique et la structure décisionnelle.
\end{exoresolu}

\courssub{Gestion du temps : la norme professionnelle de révision}

L'amélioration ne doit pas être infinie. La norme AFNOR NF X 08-011 (rédaction de documents techniques) donne une répartition temporelle de référence pour un rapport de qualité.

\begin{exemplebox}[Répartition temporelle type (Norme AFNOR NF X 08-011)]
Pour un rapport technique de niveau professionnel (et visé au Probatoire) :
\begin{itemize}
    \item \textbf{20\% du temps : Planification} (analyse commande, collecte données, élaboration plan détaillé, choix iconographie).
    \item \textbf{50\% du temps : Rédaction} (développement parties, insertion figures, rédaction intro/conclusion).
    \item \textbf{30\% du temps : Révision} (3 lectures : structure, phrase, mot ; remédiation pairs ; mise en page finale ; relecture ortho/typo).
\end{itemize}
\textbf{Application :} Sur 13 séances (progression), environ 2-3 séances planification, 6-7 séances rédaction, 3-4 séances révision/remédiation/finalisation. Ne sacrifiez pas les 30\% de révision : c'est là que se joue la note.
\end{exemplebox}

\courssub{Synthèse : la check-list finale avant remise}

Avant de considérer le rapport comme terminé, passez en revue cette liste unique, synthèse de toute la section 5.

\begin{aretenir}[Check-list finale du rapport technique (Section 5)]
\begin{itemize}
    \item [ ] \textbf{Plan} : Validé par auto-évaluation (10 critères) + remédiation pairs (grille 4 critères).
    \item [ ] \textbf{Introduction} : Accroche + Sujet + Division (annonce plan exacte).
    \item [ ] \textbf{Développement} : Parties équilibrées, titres parallèles explicites, connecteurs logiques, \textbf{passif/impersonnel} dominant.
    \item [ ] \textbf{Iconographie} : Chaque figure (Tableau, Graphique, Photo, Schéma) a : N° + Titre + Corps + Légende technique + Source + \textbf{Appel dans le texte avant apparition}.
    \item [ ] \textbf{Conclusion} : Bilan chiffré + Ouverture + Recommandations numérotées actionnables.
    \item [ ] \textbf{3 Lectures de révision} : 1. Structure (plan tenu ?) 2. Phrase (syntaxe, connecteurs, passif) 3. Mot (vocabulaire technique, emprunts, orthographe).
    \item [ ] \textbf{Présentation} : Police 12, interligne 1,5, marges normées, pagination, sommaire, liste des figures/tableaux.
\end{itemize}
\end{aretenir}

\begin{savaistu}[Le rapport technique au Cameroun : du Probatoire au monde professionnel]
Au Cameroun, le rapport de stage (ou d'activité) est l'examen de passage obligé en Première Technique (Probatoire) et en Terminale (Baccalauréat Technique). Les jurys (inspecteurs, chefs de travaux, professionnels) sanctionnent sévèrement : l'absence de plan détaillé numéroté, les figures « orphelines » (sans titre, sans source, sans appel), le vocabulaire courant à la place du terme technique (ex: « tuyau » au lieu de « canalisation » ou « conduite »), l'absence de recommandations concrètes en conclusion. Maîtriser cette section 5, c'est sécuriser 40\% de la note de production d'écrit utilitaire. C'est aussi acquérir une compétence transversale : tout technicien supérieur (BTS, HND, DUT) rédige des rapports d'intervention, des PV de réception, des notes techniques sur ce même modèle normé.
\end{savaistu}

\begin{exercice}[Entraînement : Amélioration d'un plan et insertion iconographique]
\textbf{Consigne :} À partir du plan brut ci-dessous, appliquez la méthode de la section 5.
\begin{enumerate}
    \item Corrigez la structure (fusion, scission, reformulation) pour obtenir un plan validé (10 critères).
    \item Positionnez 3 figures pertinentes (type, titre, emplacement dans le plan).
    \item Rédigez l'introduction (3 mouvements) et la conclusion (3 mouvements).
    \item Échangez votre production avec un pair et remplissez la grille de remédiation (4 critères).
\end{enumerate}

\textbf{Plan brut (Sujet : Stage maintenance préventive climatisation centrale Hôpital Régional) :} \\
1. L'hôpital et le service maintenance \\
2. Les climatiseurs (description) \\
3. Ce qu'on a fait comme entretien \\
4. Les pannes qu'on a trouvées \\
5. Les pièces changées \\
6. Conclusion
\end{exercice}

\begin{corrige}[Exercice n°14-5 : Amélioration plan climatisation]
\textbf{1. Plan validé (après fusion/scission/reformulation) :} \\
1. Cadre du stage : Hôpital Régional et Service Maintenance Énergétique \\
2. Caractéristiques techniques de la centrale de traitement d'air (CTA) \\
3. Opérations de maintenance préventive réalisées (calendrier, procédures, contrôles) \\
4. Anomalies détectées, diagnostic et interventions correctives associées \\
5. Gestion des pièces de rechange et consommables (traçabilité, stock) \\
6. Bilan, recommandations techniques et planning prévisionnel

\textbf{2. Figures positionnées :} \\
\relax[Fig. 1 - Schéma fonctionnel CTA] en 2.1 \\
\relax[Fig. 2 - Courbe évolution $\Delta P$ filtres (Semaines 1-4)] en 3.2 \\
\relax[Fig. 3 - Camembert répartition temps interventions (Préventif / Correctif / Attente)] en 6.1

\textbf{3. Introduction type :} « La qualité de l'air et la thermostabilité des blocs opératoires dépendent de la fiabilité des centrales de traitement d'air (CTA). [Accroche] Ce rapport présente les activités de maintenance préventive de niveau 2 effectuées sur la CTA n$^\circ$3 du bloc opératoire de l'Hôpital Régional de Bertoua, du 03 au 28 avril 2026. [Sujet] Nous décrirons d'abord l'installation (1), puis les opérations préventives (2), les anomalies traitées (3), la gestion des pièces (4), avant de formuler des recommandations (5). [Division] »

\textbf{4. Conclusion type :} « Le stage a permis de réaliser 100\% des gammes préventives mensuelles et de corriger 3 fuites frigorigènes critiques, ramenant le COP de 2,8 à 3,4. [Bilan] L'absence d'analyseur de vibrations a limité le diagnostic roulements compresseurs. [Ouverture] Recommandations : 1) Acquérir analyseur vibrations. 2) Instaurer relevé thermographique mensuel. 3) Former binôme sur procédure purge azote. [Recommandations] »
\end{corrige}

\coursec{Situation d'intégration : Du terrain au rapport final}

Dans la section précédente, nous avons appris à \cle{réviser} et à \cle{améliorer} un plan détaillé : vérification de la cohérence logique, équilibre des parties, insertion des documents iconographiques. Toutes ces compétences, acquises pas à pas depuis le début de cette leçon (structures du rapport, phrase complexe, emprunts lexicaux, texte explicatif, planification), vont maintenant être mobilisées \emph{ensemble}, dans une situation professionnelle complète et réaliste. C'est le principe même de la \cle{situation d'intégration} : on n'apprend plus une notion isolée, on \textbf{agit} comme un technicien qui doit rendre compte de son travail par écrit.

\courssub{Mise en situation authentique : stage à la SONARA}

\begin{definition}[La situation d'intégration]
Une \cle{situation d'intégration} est une tâche complexe, proche de la vie professionnelle réelle, qui oblige l'élève à mobiliser simultanément plusieurs savoirs et savoir-faire de l'unité pour produire un résultat concret et utile (ici : un rapport technique complet et son envoi officiel).
\end{definition}

\textbf{Contexte.} Tu es élève de Première technique, en stage d'observation et de pratique de quatre semaines à la \textbf{SONARA} (Société Nationale de Raffinage) de Limbe, au \emph{Service Maintenance}. Ton tuteur, M. Etoundi, ingénieur maintenance, te confie une mission : suivre le comportement d'un \cle{échangeur de chaleur} (appareil qui transfère la chaleur d'un fluide chaud vers un fluide froid) de l'unité de distillation atmosphérique, relever ses paramètres de fonctionnement, et signaler toute anomalie. À la fin de la mission, tu dois rédiger un \textbf{rapport de stage technique} et le transmettre par mail à ton tuteur et au chef de service.

\begin{savaistu}[La SONARA et la culture du rapport]
La SONARA, située à Limbe (région du Sud-Ouest), est la seule raffinerie du Cameroun. Dans ce type d'industrie à haut risque, les rapports suivent des normes \cle{HSE} (Hygiène, Sécurité, Environnement) très strictes : tout incident, même mineur (une simple goutte de fuite), génère un rapport structuré, souvent selon la méthode des \emph{5 Pourquoi} (on se demande « pourquoi ? » cinq fois de suite pour remonter à la cause racine du problème). Savoir rédiger un rapport clair n'est donc pas un exercice scolaire artificiel : c'est une \textbf{compétence professionnelle vitale}.
\end{savaistu}

\textbf{Le dossier documentaire fourni.} Comme dans la réalité professionnelle — et comme à l'épreuve du Probatoire et du Baccalauréat technique —, les informations ne te sont pas données « en ordre ». Tu reçois un dossier brut :

\begin{enumerate}
\item un \textbf{ordre de mission} signé du chef de service (objet, dates, périmètre de la mission) ;
\item \textbf{trois pages de notes manuscrites} prises sur le terrain, en désordre, avec des abréviations et des ratures ;
\item \textbf{deux schémas techniques} de l'échangeur de chaleur (vue en coupe et implantation) ;
\item un \textbf{tableau brut de relevés} de température et de pression (export Excel, 20 lignes de mesures horaires) ;
\item \textbf{deux photographies d'anomalies} : corrosion sur le faisceau tubulaire, fuite au niveau d'un joint de bride ;
\item un \textbf{mail du tuteur} précisant les consignes : « Rapport de 3 à 4 pages, plan en numérotation décimale, insertion obligatoire des figures et du tableau de mesures, envoi en PDF avant vendredi 17 h. »
\end{enumerate}

\courssub{Tâche 1 (Réception) : analyser le dossier et identifier le problème}

Avant d'écrire la moindre phrase, il faut \textbf{comprendre} le dossier. C'est la phase de \cle{réception} : lire, trier, hiérarchiser.

\begin{methode}[Analyser un dossier documentaire en quatre étapes]
\begin{enumerate}
\item \textbf{Lire d'abord les documents « cadres »} : l'ordre de mission et le mail du tuteur. Ils définissent le \emph{type} de rapport attendu, le destinataire, le délai et les contraintes de forme.
\item \textbf{Dégager le problème technique central} : comparer les relevés chiffrés (tableau) aux valeurs normales indiquées sur les schémas, puis confronter aux photos et aux notes de terrain.
\item \textbf{Trier l'information} en trois tas : informations \emph{essentielles} (elles iront dans le rapport), informations \emph{secondaires} (annexes éventuelles), informations \emph{inutiles} (à écarter sans regret).
\item \textbf{Reconstituer la chronologie} : quand l'anomalie est-elle apparue ? Comment a-t-elle évolué ? Quelles actions ont déjà été menées ?
\end{enumerate}
\end{methode}

\begin{exemplebox}[Application au dossier SONARA]
En croisant les documents, l'élève découvre que :
\begin{itemize}
\item le tableau Excel montre une \textbf{baisse progressive du rendement thermique} : la température de sortie du fluide froid passe de \SI{85}{\celsius} (semaine 1) à \SI{71}{\celsius} (semaine 4), alors que la pression d'entrée augmente de 12 à \SI{14,5}{bar} ;
\item les notes de terrain mentionnent dès la semaine 2 un « bruit anormal » et des « traces brunes » près de la bride côté sortie ;
\item la photo n° 1 confirme une \textbf{corrosion} du faisceau tubulaire ; la photo n° 2 montre une \textbf{fuite} au joint de bride.
\end{itemize}
\textbf{Problème central identifié} : l'encrassement et la corrosion de l'échangeur provoquent une perte de rendement thermique et une fuite naissante, avec un risque HSE. Toute information sans lien avec ce problème (par exemple une note sur la cantine du site) est écartée.
\end{exemplebox}

\begin{attention}[Piège du Bac : le type de rapport n'est pas toujours nommé]
À l'examen, le sujet dit souvent simplement : « Rédigez un rapport sur... » \textbf{sans préciser le type}. C'est à toi de le \emph{déduire} des documents fournis : présence d'un ordre de mission et d'un tuteur $\Rightarrow$ rapport de \textbf{stage} ; description d'un accident avec témoins $\Rightarrow$ rapport d'\textbf{incident} ; compte rendu régulier d'une activité $\Rightarrow$ rapport d'\textbf{activité}. Ici, l'ordre de mission, le tuteur et la durée de quatre semaines désignent clairement un \textbf{rapport de stage technique} à dominante diagnostic.
\end{attention}

\courssub{Tâche 2 (Production) : élaborer le plan détaillé complet}

Le plan détaillé est le \textbf{squelette} du rapport. Il s'élabore au brouillon, avec une \cle{numérotation décimale} (1., 1.1, 1.1.1...) et des \cle{titres nominaux} (groupes nominaux sans verbe conjugué : « Analyse des anomalies » et non « Nous analysons les anomalies »).

\begin{methode}[Construire le plan détaillé en cinq étapes]
\begin{enumerate}
\item \textbf{Choisir la structure externe} adaptée au type de rapport : en-tête (logo, références, date, objet), introduction, développement, conclusion, annexes.
\item \textbf{Ordonner le fond selon une logique} : ici, la logique du diagnostic = \emph{contexte} $\rightarrow$ \emph{constats} $\rightarrow$ \emph{analyse des causes} $\rightarrow$ \emph{propositions de solutions}.
\item \textbf{Répartir les documents} du dossier dans les parties : chaque document doit trouver sa place (schémas dans la présentation de l'installation, tableau dans les constats, photos dans l'analyse des anomalies).
\item \textbf{Rédiger les titres nominaux} et vérifier la numérotation décimale : pas de 1.2 sans 1.1, pas de partie déséquilibrée.
\item \textbf{Annoter le plan} : noter en marge les documents à insérer (Figure 1, Tableau 1) et les idées de phrases de transition.
\end{enumerate}
\end{methode}

\begin{exoresolu}[Le plan détaillé du rapport SONARA]
Énoncé : à partir du dossier documentaire, établis le plan détaillé complet du rapport de stage, avec numérotation décimale et titres nominaux.
\tcblower
\textbf{Solution détaillée.}

\textbf{Structure externe} : page de garde (SONARA, Service Maintenance, nom du stagiaire, période, objet) ; sommaire ; introduction ; développement ; conclusion ; annexes (notes de terrain, photos, tableau Excel brut).

\textbf{Structure interne (développement)} :

\textbf{Introduction} : présentation de la SONARA et du Service Maintenance ; objet et déroulement de la mission ; annonce du plan.

\textbf{1. Présentation de l'installation étudiée}

1.1. Le rôle de l'échangeur de chaleur dans l'unité de distillation atmosphérique

1.2. Description technique de l'appareil (insertion de la Figure 1 : schéma en coupe)

\textbf{2. Déroulement de la mission et méthodologie}

2.1. Le protocole de relevés (température, pression, fréquence horaire)

2.2. Les observations de terrain

\textbf{3. Constat des dysfonctionnements}

3.1. L'évolution des paramètres de fonctionnement (insertion du Tableau 1 : relevés)

3.2. Les anomalies visuelles observées (insertion de la Figure 2 : photos)

\textbf{4. Analyse des anomalies}

4.1. La corrosion du faisceau tubulaire : causes probables

4.2. La fuite au joint de bride : origine et conséquences

4.3. L'impact sur le rendement thermique et les risques HSE

\textbf{5. Propositions de solutions}

5.1. Les actions correctives immédiates (remplacement du joint)

5.2. Les actions préventives (programme de contrôle anticorrosion)

\textbf{Conclusion} : bilan de la mission, apports du stage, remerciements.

\textbf{Annexes} : A. Notes de terrain ; B. Photographies ; C. Tableau de relevés brut.
\end{exoresolu}

\begin{exemplebox}[Ce que vaut un bon plan à l'examen]
À l'épreuve de Français du Baccalauréat technique, un plan détaillé \textbf{cohérent, hiérarchisé et correctement numéroté} rapporte à lui seul \textbf{4 à 6 points sur 20} rien que pour la structure, avant même la lecture du contenu rédactionnel. Un plan sans numérotation, avec des titres verbalisés ou des parties déséquilibrées, fait perdre ces points immédiatement. Le brouillon n'est donc pas du temps perdu : c'est un investissement rentable.
\end{exemplebox}

\courssub{Tâche 3 (Production) : rédiger l'introduction et la partie « Analyse des anomalies »}

La rédaction transforme le plan en texte. Trois exigences linguistiques travaillées dans cette leçon doivent être visibles : la \cle{phrase complexe} (subordination pour exprimer cause, conséquence, opposition), le \cle{vocabulaire technique précis} (y compris les emprunts : \emph{joint}, \emph{pipeline}, \emph{monitoring}) et la \cle{voix passive}, qui garantit l'objectivité du texte explicatif (« la fuite a été localisée » plutôt que « j'ai vu une fuite »).

\begin{exemplebox}[Introduction rédigée (modèle)]
« Dans le cadre de notre formation en Première technique, un stage de quatre semaines a été effectué à la SONARA de Limbe, au sein du Service Maintenance. La mission qui nous a été confiée consistait à suivre le comportement d'un échangeur de chaleur de l'unité de distillation atmosphérique, afin de détecter d'éventuelles dérives de fonctionnement. \textbf{Puisque} les relevés ont révélé une baisse anormale du rendement thermique, une inspection visuelle complète a été décidée. Le présent rapport expose d'abord l'installation étudiée et la méthodologie suivie, puis il analyse les anomalies constatées avant de proposer des solutions correctives et préventives. »
\end{exemplebox}

\begin{exemplebox}[Développement rédigé : partie 4 « Analyse des anomalies » (extrait)]
« \textbf{4.1. La corrosion du faisceau tubulaire : causes probables.} L'examen du faisceau tubulaire, dont la Figure 2 (photo n° 1) montre l'état de surface, révèle une corrosion localisée qui s'est développée progressivement depuis la deuxième semaine d'observation. Cette dégradation peut s'expliquer par la présence d'eau oxygénée dans le circuit, \textbf{bien que} le traitement chimique de l'eau soit normalement assuré. \textbf{Comme} les tubes corrodés laissent passer des dépôts, l'échange thermique est perturbé, \textbf{ce qui explique} la chute de température de sortie constatée au Tableau 1 : celle-ci est passée de \SI{85}{\celsius} à \SI{71}{\celsius} en quatre semaines.

\textbf{4.2. La fuite au joint de bride : origine et conséquences.} Une fuite a été localisée au niveau du joint de la bride côté sortie (Figure 2, photo n° 2). \textbf{Bien qu'}elle soit encore faible, cette fuite doit être traitée sans délai, \textbf{car} elle expose le personnel à un risque de brûlure et l'environnement à une pollution hydrocarburée. Il est probable que le joint, qui n'a pas été remplacé lors de la dernière révision, ait vieilli sous l'effet conjugué de la chaleur et de la pression, laquelle est passée de 12 à \SI{14,5}{bar}. »
\end{exemplebox}

Remarque la mécanique du texte : chaque phrase complexe articule une \textbf{cause} (\emph{parce que}, \emph{comme}, \emph{puisque}), une \textbf{concession} (\emph{bien que} + subjonctif) ou une \textbf{conséquence} (\emph{ce qui explique}), tandis que le passif maintient le ton objectif exigé par le genre du rapport.

\courssub{Tâche 4 (Communication) : le mail professionnel d'envoi}

Le rapport terminé doit être \textbf{transmis} selon les codes de la communication professionnelle par internet : objet explicite, formule d'appel, annonce de la pièce jointe, formule de politesse, signature complète.

\begin{exemplebox}[Mail d'envoi du rapport (modèle)]
\textbf{Objet :} Rapport de stage — Suivi de l'échangeur de chaleur E-201 (période du 02 au 27 juin)

Monsieur,

Conformément à vos consignes, je vous adresse en pièce jointe, au format PDF, mon rapport de stage portant sur le suivi de l'échangeur de chaleur de l'unité de distillation atmosphérique. Ce document présente les relevés effectués, analyse les anomalies constatées (corrosion du faisceau tubulaire et fuite au joint de bride) et propose des actions correctives et préventives.

Je reste à votre disposition pour tout complément d'information.

Je vous prie d'agréer, Monsieur, l'expression de ma respectueuse considération.

NDOUMBE Arnaud — Élève de Première technique, stagiaire Service Maintenance
\end{exemplebox}

\begin{attention}[Erreurs fréquentes dans le mail professionnel]
Évite : l'objet vide ou vague (« Bonjour ») ; les abréviations de type SMS (« slt », « bjr ») ; l'oubli de la pièce jointe annoncée ; l'absence de signature ; l'envoi au mauvais destinataire (ici : tuteur \emph{et} chef de service, le second en copie). Le mail fait partie de la production évaluée : il est rédigé avec le même soin que le rapport.
\end{attention}

\courssub{Organisation du travail : planning et flux de production}

Une production complexe exige une \textbf{gestion du temps}. Le diagramme de Gantt ci-dessous répartit les treize séances de la leçon ; le schéma de flux rappelle que la rédaction n'est jamais linéaire : on révise, on améliore, on réécrit.

\begin{popfigure}
\begin{tikzpicture}[x=1.05cm,y=0.52cm]
% Axes
\draw[popline] (0,0) -- (13.4,0);
\foreach \s in {1,...,13}{
  \draw[popline] (\s,0) -- (\s,-0.15) node[below,font=\scriptsize] {S\s};
  \draw[popline!40] (\s,0) -- (\s,8.2);
}
% Tâches : label / début / durée / couleur
\foreach [count=\i] \label/\deb/\dur/\col in {
  {Analyse des documents}/1/2/popblue,
  {Linguistique (phrase complexe, emprunt)}/3/2/poppurple,
  {Planification (plan d\'etaill\'e)}/5/2/popgold,
  {R\'edaction du rapport}/7/3/popgreen,
  {Iconographie (figures, tableau)}/10/1/poporange,
  {R\'evision et am\'elioration}/11/1/poppink,
  {Rem\'ediation}/12/1/popmuted,
  {\'Evaluation}/13/1/popink}{
  \pgfmathsetmacro{\y}{8-\i}
  \fill[\col!70,draw=\col] (\deb-1,\y+0.15) rectangle (\deb-1+\dur,\y+0.85);
  \node[anchor=west,font=\scriptsize] at (13.6,\y+0.5) {\label};
}
\node[font=\small\bfseries,text=popdark] at (6.5,8.6) {Planning des 13 s\'eances (diagramme de Gantt)};
\end{tikzpicture}
\legende{Diagramme de Gantt de la leçon : chaque barre colorée indique les séances (S1 à S13) consacrées à chaque étape, de l'analyse du dossier à l'évaluation finale.}
\end{popfigure}

\begin{popfigure}
\begin{tikzpicture}[node distance=1.6cm and 2.4cm,
  bloc/.style={rectangle,rounded corners,draw=popblue,fill=popblueL,thick,align=center,minimum height=1cm,inner sep=6pt,font=\small},
  fleche/.style={-{Stealth[length=3mm]},thick,popdark},
  retro/.style={-{Stealth[length=3mm]},thick,dashed,poporange}]
\node[bloc, align=center] (notes){Notes brutes\\(terrain, relev\'es, photos)};
\node[bloc,right=of notes, align=center] (tri){Tri et analyse\\de l'information};
\node[bloc,right=of tri, align=center] (plan){Plan d\'etaill\'e\\(num\'erotation d\'ecimale)};
\node[bloc,below=of plan, align=center] (redac){R\'edaction\\du rapport final};
\node[bloc,left=of redac, align=center] (rev){R\'evision\\(coh\'erence, langue)};
\draw[fleche] (notes) -- (tri);
\draw[fleche] (tri) -- (plan);
\draw[fleche] (plan) -- (redac);
\draw[fleche] (redac) -- (rev);
\draw[retro] (rev) -- node[above,font=\scriptsize,text=poporange]{am\'eliorer} (plan);
\draw[retro] (rev.north) to[bend left=25] node[left,font=\scriptsize,text=poporange]{corriger} (redac.west);
\end{tikzpicture}
\legende{Flux de travail du rapport : les notes brutes sont triées, transformées en plan détaillé puis rédigées ; les flèches orange en pointillés montrent les boucles de rétroaction (réviser $\rightarrow$ améliorer $\rightarrow$ réécrire).}
\end{popfigure}

\begin{methode}[Pour le Bac Tech : gérer l'épreuve de production d'écrits]
\begin{enumerate}
\item \textbf{Lire le sujet et tout le dossier} (10 min) : identifier le type de rapport, le destinataire, le problème central.
\item \textbf{Faire le plan détaillé au brouillon} (15 min) : numérotation décimale, titres nominaux, répartition des documents. \emph{Le plan détaillé est noté !}
\item \textbf{Rédiger proprement} (1 h 15) : introduction, développement structuré, phrases complexes, voix passive, vocabulaire précis.
\item \textbf{Relire} (15 min) : orthographe, accords, cohérence des renvois aux figures et tableaux, numérotation.
\end{enumerate}
\end{methode}

\begin{exercice}[À toi de jouer : le rapport complet]
À partir du dossier SONARA présenté dans cette section :
\begin{enumerate}
\item Rédige la partie 3 « Constat des dysfonctionnements » (15 à 20 lignes) en insérant le renvoi au Tableau 1 et à la Figure 2, avec au moins trois phrases complexes et deux verbes à la voix passive.
\item Élabore le plan détaillé d'un \emph{rapport d'incident} relatif à la fuite du joint de bride (structure adaptée : faits, causes, mesures conservatoires, recommandations).
\item Rédige le mail d'envoi de ce rapport d'incident au chef de service, avec copie au responsable HSE.
\end{enumerate}
\end{exercice}

\begin{corrige}[Exercice : éléments de correction]
\begin{enumerate}
\item Le constat doit rester \emph{factuel} : « Les relevés horaires, consignés dans le Tableau 1, montrent que la température de sortie a diminué de \SI{14}{\celsius} en quatre semaines, tandis que la pression d'entrée augmentait de \SI{2,5}{bar}. Des traces brunes ont été observées près de la bride côté sortie, comme le confirme la Figure 2... » (passif : « ont été observées », « ont été consignés » ; subordonnées de temps et de cause).
\item Plan d'rapport d'incident attendu : \textbf{1.} Description des faits (1.1. circonstances de la découverte ; 1.2. localisation de la fuite) ; \textbf{2.} Analyse des causes (2.1. cause immédiate : joint usé ; 2.2. cause profonde : absence de remplacement à la révision) ; \textbf{3.} Mesures conservatoires prises ; \textbf{4.} Recommandations (4.1. remplacement du joint ; 4.2. renforcement du plan de maintenance préventive).
\item Le mail doit comporter : objet précis mentionnant « rapport d'incident », annonce de la pièce jointe PDF, mention de la copie au responsable HSE, formule de politesse administrative et signature complète.
\end{enumerate}
\end{corrige}

\begin{aretenir}[L'essentiel de la situation d'intégration]
\begin{itemize}
\item Une situation d'intégration mobilise \textbf{toutes} les compétences de la leçon : réception (trier un dossier), langue (phrase complexe, emprunts, passif), production (plan détaillé, rédaction, mail).
\item Le \textbf{plan détaillé} (numérotation décimale, titres nominaux) est la charpente du rapport : il est construit au brouillon et il est \textbf{noté} à l'examen.
\item La rédaction n'est jamais linéaire : \emph{réviser $\rightarrow$ améliorer $\rightarrow$ réécrire}.
\item Le rapport se termine par une \textbf{transmission professionnelle} : un mail respectant les codes de la communication par internet.
\end{itemize}
\end{aretenir}

\newpage

\coursec{Activité d'intégration}

\courssub{Objectifs de l'activité}

À l'issue de cette activité d'intégration, l'élève sera capable de :
\begin{itemize}
\item analyser un dossier documentaire professionnel (notes de terrain, relevés chiffrés, photographies, schéma technique, courriel) afin d'en extraire les informations pertinentes ;
\item construire le \cle{plan détaillé} complet d'un rapport d'incident en respectant la structure externe (page de garde, sommaire, annexes) et la structure interne (introduction, développement, conclusion) ;
\item rédiger une partie du rapport en employant la \cle{phrase complexe}, un vocabulaire technique précis (y compris les \cle{emprunts} lexicalisés du domaine industriel) et un ton \cle{objectif} ;
\item intégrer et légender correctement des documents iconiques (schéma annoté, graphique) en les exploitant dans le texte ;
\item rédiger un courriel professionnel d'accompagnement respectant les codes de la communication par internet.
\end{itemize}

\courssub{Situation}

Tu es élève en classe de Première technique, spécialité Génie mécanique, en stage d'observation à la \textbf{SONARA} (Société Nationale de Raffinage) de Limbé, dans la région du Sud-Ouest. Ton maître de stage, M. Etoundi, ingénieur maintenance, te confie une mission réelle : rédiger le \textbf{rapport d'anomalie} concernant l'échangeur de chaleur \textbf{E-102} de l'unité de distillation.

\begin{definition}[Rappel : le rapport d'anomalie]
Le \cle{rapport d'anomalie} est un texte fonctionnel et utilitaire qui décrit un dysfonctionnement constaté sur un équipement, en explique les causes probables et propose des actions correctives. Il obéit à une double exigence : l'\emph{exactitude} (faits vérifiables, chiffres précis) et l'\emph{objectivité} (aucune émotion, aucune accusation).
\end{definition}

Voici le dossier « Terrain » dont tu disposes.

\textbf{Document 1 : notes manuscrites de l'opérateur de quart (retranscrites).}
« Lundi 12 mai, 6 h 10 : odeur d'hydrocarbures près du faisceau tubulaire. 6 h 25 : goutte-à-goutte visible sous la calandre, côté passe 3. Température de sortie produit en baisse depuis la nuit. Mise en sécurité de l'échangeur à 7 h 00, by-pass ouvert. Aucun blessé. Petite flaque de produit récupérée au bac de rétention (environ 15 litres). »

\textbf{Document 2 : relevé des températures de sortie produit (en degrés Celsius), enregistrées toutes les deux heures.}

\begin{center}
\begin{tabularx}{\linewidth}{|l|X|X|X|X|X|X|}\hline
\rowcolor{popblueL}
Heure & 22 h & 0 h & 2 h & 4 h & 6 h & 8 h (by-pass) \\ \hline
Température de sortie (°C) & 118 & 116 & 112 & 107 & 101 & --- \\ \hline
Température normale (°C) & 120 & 120 & 120 & 120 & 120 & 120 \\ \hline
\end{tabularx}
\end{center}

\textbf{Document 3 : photographies prises à l'ouverture de la calandre.}
Photo A : dépôts brunâtres (boues et sels) sur la surface externe des tubes de la passe 3. Photo B : sous les dépôts, piqûres de corrosion localisées ; un tube présente une perforation d'environ 3 mm de diamètre.

\textbf{Document 4 : schéma technique simplifié de l'échangeur E-102} (faisceau tubulaire, calandre, boîtes de distribution, passes 1 à 4, sens de circulation des fluides).

\textbf{Document 5 : courriel de M. Etoundi.}
« Bonjour. Merci de me transmettre avant vendredi 17 h le rapport d'anomalie E-102 selon le plan type de l'entreprise : 1. Circonstances ; 2. Constats ; 3. Causes probables ; 4. Conséquences ; 5. Actions correctives. Joignez le schéma annoté (Figure 1) et le graphique de dérive de température (Figure 2). Rédigez en entier la partie "Constats et analyse" (environ deux pages). Restez factuel : pas de rumeur, pas d'accusation. Bon courage. »

\courssub{Consignes}

\begin{enumerate}
\item \textbf{Analyse du dossier.} Classe les informations du dossier « Terrain » dans un tableau à deux colonnes : « Faits observés » et « Interprétations possibles ». Justifie pourquoi cette distinction est indispensable dans un rapport technique.
\item \textbf{Diagnostic.} À partir des documents 1, 2 et 3, formule en une phrase complexe (proposition principale + au moins deux subordonnées) le diagnostic de l'anomalie : corrosion sous dépôts ayant entraîné la perforation d'un tube et une fuite de produit.
\item \textbf{Plan détaillé.} Construis le plan détaillé complet du rapport : structure externe (page de garde, sommaire, corps, annexes) et structure interne (introduction, les cinq parties imposées par le plan type, conclusion). Pour chaque partie, indique sous forme de sous-parties numérotées les informations précises du dossier que tu y placeras.
\item \textbf{Rédaction.} Rédige la partie « Constats et analyse » (environ deux pages, soit 400 à 500 mots) en respectant les règles du texte explicatif objectif : présent de vérité générale ou passé composé pour les faits, vocabulaire technique, connecteurs logiques, aucune émotion ni accusation.
\item \textbf{Iconographie.} Présente la Figure 1 (schéma annoté : tu décriras les annotations à porter sur le schéma) et la Figure 2 (graphique de la dérive de température que tu traceras à partir du document 2). Rédige pour chacune un titre, une légende et une phrase d'appel dans le texte (« comme le montre la Figure 2... »).
\item \textbf{Courriel d'envoi.} Rédige le courriel professionnel accompagnant l'envoi du rapport à M. Etoundi : objet clair, formule d'appel, annonce de la pièce jointe, formule de politesse, signature.
\end{enumerate}

\courssub{Ressources à mobiliser}

\begin{aretenir}[Ce que tu dois réactiver]
\begin{itemize}
\item \textbf{Structure externe du rapport} : page de garde (titre, nom, date, destinataire), sommaire, corps du rapport, annexes numérotées.
\item \textbf{Structure interne} : introduction (contexte, problème, annonce du plan), développement en parties numérotées, conclusion (bilan et perspectives).
\item \textbf{La phrase complexe} : subordonnées relatives (\emph{qui, que, dont, où}), conjonctives (\emph{que, parce que, bien que}), circonstancielles de cause et de conséquence (\emph{si bien que, de sorte que}).
\item \textbf{L'emprunt lexical} : mots techniques empruntés à l'anglais et intégrés au français industriel (\emph{by-pass}, \emph{pipeline}, \emph{maintenance}) : à utiliser avec exactitude, sans excès.
\item \textbf{Le texte explicatif} : objectivité, connecteurs (\emph{d'abord, ensuite, en effet, par conséquent, enfin}), définitions des termes spécialisés, ton neutre (ni dramatique ni comique).
\item \textbf{Le texte iconique} : toute figure porte un numéro, un titre, une légende ; elle est toujours appelée et commentée dans le texte.
\item \textbf{Les codes du courriel} : objet explicite, concision, politesse, pas d'abréviation de type SMS.
\end{itemize}
\end{aretenir}

\begin{attention}[Pièges fréquents dans un rapport technique]
\begin{itemize}
\item Confondre \emph{fait} et \emph{interprétation} : « le tube est percé » est un fait ; « l'opérateur a été négligent » est un jugement, à bannir.
\item Employer un ton dramatique (« catastrophe », « désastre ») ou comique : le rapport exige la neutralité.
\item Insérer une figure sans l'appeler ni la commenter dans le texte : une figure « muette » est une faute de méthode.
\item Multiplier les emprunts à l'anglais alors qu'un terme français existe : on dira « dérivation » ou « by-pass » selon l'usage de l'entreprise, mais sans mélanger les registres.
\end{itemize}
\end{attention}

\courssub{Démarche et corrigé commenté}

\begin{exoresolu}[Mission SONARA : rapport d'anomalie E-102]
\textbf{Consigne 1.} Distingue faits et interprétations à partir du dossier « Terrain ».

\tcblower

\textbf{Consigne 1 — Faits et interprétations.}

\begin{center}
\begin{tabularx}{\linewidth}{|X|X|}\hline
\rowcolor{popblueL}
Faits observés & Interprétations possibles \\ \hline
Odeur d'hydrocarbures à 6 h 10 ; goutte-à-goutte sous la calandre à 6 h 25 & Fuite de produit à l'intérieur de l'échangeur \\ \hline
Baisse régulière de la température de sortie : 118 °C à 22 h, puis 101 °C à 6 h & Perte d'efficacité thermique liée à la fuite ou à l'encrassement \\ \hline
Dépôts brunâtres sur les tubes de la passe 3 ; piqûres sous les dépôts ; perforation de 3 mm & Corrosion sous dépôts ayant percé un tube \\ \hline
Environ 15 litres récupérés en rétention ; aucun blessé & Fuite limitée et maîtrisée \\ \hline
\end{tabularx}
\end{center}

\emph{Commentaire :} un rapport technique ne juge ni ne suppose : il sépare ce qui est constaté (mesurable, visible) de ce qui est déduit. C'est la condition de l'objectivité.

\textbf{Consigne 2 — Diagnostic en une phrase complexe.}
« L'échangeur E-102, dont les tubes de la passe 3 étaient recouverts de dépôts, a subi une corrosion sous dépôts qui a perforé un tube, si bien que le produit s'est échappé et que la température de sortie a chuté de 19 °C en huit heures. »
\emph{Commentaire :} la phrase articule une relative (\emph{dont}), une relative (\emph{qui}) et une consécutive (\emph{si bien que}) : c'est le modèle de phrase complexe attendu dans un constat technique. Vérification du calcul : $118 - 101 = 17$ °C de baisse mesurée entre 22 h et 6 h, et $120 - 101 = 19$ °C d'écart par rapport à la température normale de service ; on précisera donc dans le rapport : « dérive de 17 °C en huit heures, soit un écart de 19 °C sous la valeur normale ».

\textbf{Consigne 3 — Plan détaillé du rapport.}

\emph{Structure externe :} page de garde (SONARA — Rapport d'anomalie E-102 — nom du stagiaire — date — destinataire : M. Etoundi) ; sommaire ; corps ; annexes (relevé horaire complet, photographies A et B).

\emph{Structure interne :}
\begin{itemize}
\item \textbf{Introduction} : contexte (stage, unité de distillation), événement (anomalie détectée le 12 mai), annonce du plan en cinq parties.
\item \textbf{1. Circonstances} : 1.1. détection par l'opérateur de quart (6 h 10, odeur ; 6 h 25, goutte-à-goutte) ; 1.2. mise en sécurité (7 h 00, ouverture du by-pass).
\item \textbf{2. Constats} : 2.1. constats externes (fuite, flaque de 15 litres en rétention) ; 2.2. constats à l'ouverture (dépôts passe 3, piqûres, perforation de 3 mm — Figure 1) ; 2.3. constats instrumentaux (dérive de température — Figure 2).
\item \textbf{3. Causes probables} : 3.1. encrassement par boues et sels ; 3.2. corrosion sous dépôts ; 3.3. perforation du tube et fuite.
\item \textbf{4. Conséquences} : 4.1. arrêt de l'échangeur et fonctionnement en by-pass ; 4.2. perte de produit (environ 15 litres) ; 4.3. baisse de performance thermique (17 °C de dérive, 19 °C d'écart à la consigne) ; 4.4. aucune conséquence humaine.
\item \textbf{5. Actions correctives} : 5.1. remplacement ou étamage du tube percé ; 5.2. nettoyage chimique du faisceau ; 5.3. renforcement du contrôle qualité de l'eau de refroidissement ; 5.4. inspection périodique des passes.
\item \textbf{Conclusion} : bilan factuel, remerciements, proposition de suivi.
\end{itemize}

\textbf{Consigne 4 — Rédaction de la partie « Constats et analyse » (extrait modèle).}
« L'inspection de l'échangeur E-102, réalisée après sa mise en sécurité le lundi 12 mai à 7 h 00, a permis d'établir les constats suivants. D'abord, à l'extérieur de l'appareil, un goutte-à-goutte de produit était visible sous la calandre, au droit de la passe 3 ; la quantité récupérée dans le bac de rétention est estimée à quinze litres. Ensuite, l'ouverture de la calandre a révélé, comme le montre la Figure 1, des dépôts brunâtres de boues et de sels sur la surface externe des tubes de la passe 3. Sous ces dépôts, qui créent un milieu confiné favorable à la corrosion, on observe des piqûres localisées ; l'une d'elles a abouti à la perforation d'un tube sur un diamètre d'environ trois millimètres. Enfin, l'analyse des relevés de température confirme la gravité du phénomène : la température de sortie produit est passée de 118 °C à 22 h à 101 °C à 6 h, soit une dérive de 17 °C en huit heures (Figure 2), alors que la valeur normale de service est de 120 °C, ce qui représente un écart de 19 °C. Cette baisse progressive, qui précède la détection visuelle de la fuite, indique que la perforation s'est produite au cours de la nuit et que la fuite a réduit l'efficacité de l'échange thermique. Par conséquent, l'ensemble des constats converge vers un mécanisme unique : l'encrassement a provoqué une corrosion sous dépôts, laquelle a percé le tube et entraîné la fuite. »
\emph{Commentaire :} présent de vérité générale et passé composé des faits, connecteurs (\emph{d'abord, ensuite, enfin, par conséquent}), définition implicite du phénomène, appels explicites aux figures, aucune émotion : le ton reste neutre, conformément au texte explicatif.

\textbf{Consigne 5 — Iconographie.}
\emph{Figure 1} : sur le schéma de l'échangeur, annoter : la calandre, le faisceau tubulaire, la passe 3 (entourée), la zone de dépôts (flèche), le tube perforé (flèche + mention « perforation 3 mm »). Titre : « Figure 1 : schéma annoté de l'échangeur E-102 — localisation de l'anomalie (passe 3) ».
\emph{Figure 2} : graphique cartésien ; abscisse : heures (22 h, 0 h, 2 h, 4 h, 6 h) ; ordonnée : température en °C ; deux courbes : la consigne constante à 120 °C et la température réelle en décroissance (118, 116, 112, 107, 101). Titre : « Figure 2 : dérive de la température de sortie produit (nuit du 11 au 12 mai) ». Phrase d'appel : « comme le montre la Figure 2, la température de sortie s'écarte progressivement de la consigne dès 22 h ».

\textbf{Consigne 6 — Courriel d'envoi.}
« Objet : Envoi du rapport d'anomalie — Échangeur E-102. Monsieur, conformément à votre demande, je vous prie de bien vouloir trouver en pièce jointe le rapport d'anomalie concernant l'échangeur E-102, rédigé selon le plan type de l'entreprise et accompagné du schéma annoté (Figure 1) et du graphique de dérive de température (Figure 2). Je reste à votre disposition pour tout complément d'information. Je vous prie d'agréer, Monsieur, l'expression de mes salutations distinguées. [Nom, prénom — Première technique, stagiaire maintenance] »
\end{exoresolu}

\begin{popfigure}
\begin{center}
\begin{tikzpicture}
\begin{axis}[width=0.9\linewidth,height=6cm,xlabel={Heure},ylabel={Température (°C)},xtick={0,1,2,3,4},xticklabels={22 h,0 h,2 h,4 h,6 h},ymin=95,ymax=125,legend style={at={(0.02,0.02)},anchor=south west,font=\small},grid=major,grid style={popline!40}]
\addplot[popcurve,mark=*,color=popblue] coordinates {(0,118) (1,116) (2,112) (3,107) (4,101)};
\addplot[thick,dashed,color=poporange,domain=0:4]{120};
\legend{Température réelle,Température normale (consigne)}
\end{axis}
\end{tikzpicture}
\end{center}
\legende{Figure 2 : dérive de la température de sortie produit de l'échangeur E-102 (nuit du 11 au 12 mai). La courbe bleue (relevés du document 2) s'écarte progressivement de la consigne à 120 °C : écart de 19 °C à 6 h.}
\end{popfigure}

\courssub{Bilan de l'activité}

Cette mission a permis de vérifier que la rédaction d'un rapport technique repose sur une double compétence : une \textbf{compétence documentaire} (trier des informations hétérogènes — notes, chiffres, photos, schéma, courriel — en distinguant rigoureusement faits et interprétations) et une \textbf{compétence langagière} (organiser un plan détaillé logique, rédiger en phrases complexes objectives, employer le vocabulaire spécialisé et les emprunts techniques avec justesse, intégrer des figures légendées). L'élève a également constaté que le plan type de l'entreprise (circonstances, constats, causes, conséquences, actions) correspond exactement à la logique du texte explicatif : observer, expliquer, conclure. Enfin, le courriel d'envoi a rappelé que la communication professionnelle par internet obéit à des codes précis de clarté et de politesse. Ces acquis seront évalués selon la grille critériée : plan détaillé (4 points), contenu technique (6 points), langue et style (6 points), présentation et iconographie (4 points).

\begin{exercice}[Pour aller plus loin : rapport de visite d'atelier]
Ton établissement organise une visite de l'atelier de maintenance de la SONACOM à Douala. Rédige le \cle{plan détaillé} complet (structure externe et interne) du rapport de visite que tu remettras à ton professeur de génie mécanique : tu préciseras pour chaque partie les informations à recueillir sur le terrain (horaires, machines observées, consignes de sécurité, interviews) et les figures à joindre (plan de l'atelier, photographies légendées). Tu rédigeras ensuite l'introduction du rapport en une phrase complexe d'annonce du plan.
\end{exercice}

\newpage

\coursec{Méthodes et exercices résolus}

\coursec{Leçon 14 : Production d'écrits utilitaires — Le rapport (élaboration du plan détaillé)}

\coursec{Le rapport technique : genre, finalités et structures}

\courssub{Définition, finalités et cadre communicationnel}

\begin{definition}[Rapport technique]
Le \cle{rapport technique} est un texte fonctionnel, administratif et professionnel, rédigé par un subordonné (le rapporteur) à la demande d'un supérieur hiérarchique (le destinataire), pour rendre compte d'une mission d'observation, d'inspection, d'enquête ou d'expérimentation. Il vise à informer objectivement pour permettre une prise de décision.
\end{definition}

\begin{aretenir}[Finalités du rapport]
\begin{itemize}
\item \textbf{Informer} : transmettre des faits constatés, des données chiffrées, des observations précises.
\item \textbf{Analyser} : expliquer les causes (diagnostic) et prévoir les conséquences (pronostic).
\item \textbf{Proposer} : formuler des recommandations concrètes, réalistes et hiérarchisées.
\end{itemize}
\end{aretenir}

\begin{propriete}[Caractéristiques du genre]
\begin{enumerate}
\item \textbf{Objectivité} : exclusion du \emph{je} subjectif, des jugements de valeur, des tonalités dramatique ou comique.
\item \textbf{Structure codifiée} : structure externe (mentions obligatoires) et structure interne (plan détaillé) rigoureuses.
\item \textbf{Langage technique} : vocabulaire précis, phrases complexes articulées (cause, conséquence, opposition), emprunts lexicaux maîtrisés.
\end{enumerate}
\end{propriete}

\courssub{Structure externe (mentions obligatoires)}

\begin{methode}[Construire la structure externe d'un rapport]
Pour rédiger la partie matérielle d'un rapport technique, on procède par étapes :
\begin{enumerate}
\item \textbf{Identifier la situation de communication} : qui écrit ? à qui ? pour quelle raison ? (un rapport est toujours commandé par un supérieur hiérarchique ou une autorité).
\item \textbf{Placer les mentions obligatoires} : en-tête (nom de l'institution, service), lieu et date, référence éventuelle, objet du rapport (une seule ligne, précise).
\item \textbf{Nommer le destinataire} : « Monsieur le Principal », « Madame la Directrice des études », avec sa fonction exacte.
\item \textbf{Rédiger le corps} en trois moments : introduction (rappel de l'ordre de mission, dates, lieu), développement (constats, faits, chiffres), conclusion (suggestions ou recommandations).
\item \textbf{Clore} par la formule de politesse administrative, la signature, le nom et la fonction du rapporteur.
\end{enumerate}
\begin{aretenir}[Formules types à retenir]
Objet : \emph{Rapport sur \ldots} ; ouverture : \emph{Conformément à votre note de service n°\ldots en date du \ldots} ; clôture : \emph{Veuillez agréer, Monsieur le Principal, l'expression de ma haute considération.}
\end{aretenir}
\end{methode}

\begin{exoresolu}[Rédiger les mentions externes d'un rapport]
\textbf{Énoncé.} Tu es délégué de la classe de Première F4. Le Principal du lycée technique de Bafoussam t'a chargé, par note de service n° 12/PRIN/LTB du 10 mars, de visiter l'atelier de mécanique et de lui rendre compte de l'état du matériel. Rédige l'en-tête, l'objet, la formule d'ouverture et la formule de clôture de ton rapport.
\tcblower
\textbf{Solution détaillée.}
\begin{enumerate}
\item \textbf{Analyse de la situation} : le rapporteur est le délégué de classe ; le destinataire est le Principal ; la mission est définie par une note de service (référence à citer) ; le thème est l'état du matériel de l'atelier de mécanique.
\item \textbf{En-tête} : \emph{Lycée Technique de Bafoussam — Classe de Première F4}. Lieu et date : \emph{Bafoussam, le 15 mars 20\ldots}
\item \textbf{Référence et objet} : \emph{Réf. : note de service n° 12/PRIN/LTB du 10 mars.} \emph{Objet : Rapport sur l'état du matériel de l'atelier de mécanique.}
\item \textbf{Destinataire} : \emph{À l'attention de Monsieur le Principal du Lycée Technique de Bafoussam.}
\item \textbf{Formule d'ouverture} : \emph{Conformément à votre note de service n° 12/PRIN/LTB du 10 mars, j'ai procédé, le 12 mars, à la visite de l'atelier de mécanique de l'établissement. J'ai l'honneur de vous présenter le rapport suivant.}
\item \textbf{Formule de clôture} : \emph{Veuillez agréer, Monsieur le Principal, l'expression de ma respectueuse considération.} Suivie de la signature, du nom et de la fonction : \emph{Le délégué de classe, N.\ldots}
\end{enumerate}
\textbf{Conclusion.} La structure externe est complète : en-tête, lieu/date, référence, objet, destinataire, formules d'ouverture et de clôture, signature. Aucune de ces mentions ne doit manquer dans la copie.
\end{exoresolu}

\coursec{Outils linguistiques : La phrase complexe et l'emprunt lexical}

\courssub{La phrase complexe : articulation logique des constats}

\begin{definition}[Phrase complexe]
Une \cle{phrase complexe} est une phrase qui contient au moins deux propositions (donc au moins deux verbes conjugués). Elle permet d'articuler des idées entre elles (cause, conséquence, opposition, condition, temps) dans une seule unité énonciative, indispensable au texte explicatif.
\end{definition}

\begin{methode}[Analyser et produire une phrase complexe dans un rapport]
\begin{enumerate}
\item \textbf{Compter les verbes conjugués} : chaque verbe conjugué correspond à une proposition.
\item \textbf{Identifier la proposition principale} (celle qui n'est introduite par aucun mot de liaison) et les propositions subordonnées (introduites par un subordonnant : \emph{que, qui, parce que, bien que, si, quand...}) ou coordonnées (liées par \emph{mais, ou, et, donc, or, ni, car}).
\item \textbf{Préciser la nature et la fonction} de chaque subordonnée (conjonctive : COD, CC cause, CC conséquence ; relative : épithète ; infinitive : sujet, COD).
\item \textbf{Choisir le mode de liaison adapté au sens} :
\begin{itemize}
\item \emph{Juxtaposition} (ponctuation) : énumération rapide de faits.
\item \emph{Coordination} (\emph{mais, or, et, donc, car}) : opposition, addition, conclusion.
\item \emph{Subordination} (\emph{parce que, bien que, qui, que, si}) : explication, nuance, condition.
\end{itemize}
\item \textbf{Vérifier la ponctuation} : virgule entre propositions juxtaposées ; pas de virgule avant \emph{et} / \emph{ou} coordonnant deux sujets ou deux verbes d'une même proposition.
\end{enumerate}
\end{methode}

\begin{exoresolu}[Analyser et produire des phrases complexes]
\textbf{Énoncé.} 1) Analyse la phrase suivante (nature et fonction des propositions) : \emph{Les machines sont hors d'usage parce qu'elles n'ont pas été entretenues, et les élèves ne peuvent plus s'exercer.} 2) Transforme en une seule phrase complexe avec subordination : \emph{Le laboratoire manque de produits. Les TP sont annulés.}
\tcblower
\textbf{Solution détaillée.}
\begin{enumerate}
\item \textbf{Comptage} : \emph{sont}, \emph{n'ont pas été entretenues}, \emph{peuvent} : trois verbes $\rightarrow$ trois propositions.
\item \textbf{Analyse} :
\begin{itemize}
\item P1 : \emph{Les machines sont hors d'usage} (Proposition Principale).
\item P2 : \emph{parce qu'elles n'ont pas été entretenues} (Proposition Subordonnée Conjonctive, Complément Circonstanciel de Cause de P1, introduite par \emph{parce que}).
\item P3 : \emph{et les élèves ne peuvent plus s'exercer} (Proposition Coordonnée à P1 par \emph{et}, exprimant la conséquence).
\end{itemize}
\item \textbf{Transformation} : on relie les deux faits par une relation de cause à effet : \emph{Les TP sont annulés parce que le laboratoire manque de produits.} Autre solution (subordonnée relative) : \emph{Le laboratoire, qui manque de produits, ne permet plus d'organiser les TP.}
\end{enumerate}
\textbf{Conclusion.} La phrase complexe permet au rapporteur d'enchaîner constat, cause et conséquence en une seule phrase claire : c'est l'outil syntaxique de base du texte explicatif.
\end{exoresolu}

\courssub{L'emprunt lexical : identification et bon usage}

\begin{definition}[Emprunt lexical]
Un \cle{emprunt} est un mot qu'une langue (la langue d'accueil, ici le français) prend à une autre langue (la langue donneuse, souvent l'anglais en contexte technique) pour combler un vide lexical ou par prestige. On distingue :
\begin{itemize}
\item L'\textbf{emprunt direct} (intégré tel quel : \emph{mail, software, week-end, leader}).
\item Le \textbf{calque} (traduction morphème par morphème : \emph{gratte-ciel} $\leftarrow$ \emph{skyscraper}, \emph{logiciel} $\leftarrow$ \emph{software}).
\item L'\textbf{emprunt adapté} (graphie francisée : \emph{le parking} $\rightarrow$ \emph{le stationnement} recommandé).
\end{itemize}
\end{definition}

\begin{methode}[Identifier et employer correctement un emprunt en contexte administratif]
\begin{enumerate}
\item \textbf{Repérer} le terme suspect (souvent technique, informatique, management).
\item \textbf{Préciser la langue d'origine} (majoritairement l'anglais).
\item \textbf{Consulter le dictionnaire / terminologie officielle} (Commission d'enrichissement de la langue française, Termium, etc.) pour connaître l'équivalent recommandé.
\item \textbf{Appliquer la règle} : dans un rapport officiel, on privilégie \textbf{systématiquement} le terme français recommandé s'il existe. L'emprunt direct ne s'accepte que s'il n'existe pas d'équivalent français consacré (ex. \emph{laser, radar, web}).
\item \textbf{Accorder} selon les règles du français (pluriel en \emph{-s} : \emph{des logiciels, des courriels}).
\end{enumerate}
\end{methode}

\begin{exoresolu}[Repérer et justifier des emprunts]
\textbf{Énoncé.} Dans le rapport d'un élève, on lit : \emph{L'atelier dispose de trois ordinateurs, mais le software est obsolète et les emails des fournisseurs restent sans réponse.} 1) Repère les emprunts et précise leur langue d'origine. 2) Propose une version conforme au français administratif.
\tcblower
\textbf{Solution détaillée.}
\begin{enumerate}
\item \textbf{Repérage} : \emph{software} et \emph{emails} sont des emprunts directs à l'anglais. \emph{Ordinateur} est un mot français créé en 1955, ce n'est pas un emprunt.
\item \textbf{Justification} : dans un rapport administratif français, on recommande les équivalents officiels : \emph{logiciel} pour \emph{software} (calque recommandé), \emph{courriels} ou \emph{messages électroniques} pour \emph{emails}.
\item \textbf{Version corrigée} : \emph{L'atelier dispose de trois ordinateurs, mais les logiciels sont obsolètes et les courriels des fournisseurs restent sans réponse.}
\item \textbf{Accord} : \emph{logiciels} prend le \emph{s} du pluriel français ; \emph{obsolètes} s'accorde avec \emph{logiciels} (masculin pluriel).
\end{enumerate}
\textbf{Conclusion.} L'emprunt n'est pas interdit, mais le rapport, texte fonctionnel et officiel, exige le terme français recommandé chaque fois qu'il existe.
\end{exoresolu}

\coursec{Stylistique du rapport : Textes explicatifs, tonalités et objectivité}

\courssub{Le texte explicatif : marqueurs et visée}

\begin{propriete}[Marqueurs du texte explicatif dans le rapport]
\begin{itemize}
\item \textbf{Connecteurs logiques} : \emph{parce que, puisque, étant donné que} (cause) ; \emph{donc, par conséquent, c'est pourquoi, de sorte que} (conséquence) ; \emph{mais, cependant, toutefois} (opposition/ concession) ; \emph{en effet, notamment} (explication).
\item \textbf{Temps verbaux} : Présent de vérité générale (\emph{les machines fonctionnent}), Passé composé / Imparfait pour les faits constatés (\emph{le rapporteur a constaté que...}), Futur pour les préconisations (\emph{il conviendra de...}).
\item \textbf{Vocabulaire} : Précis, dénotatif, technique (pas de polysémie), quantifié (\emph{6 machines sur 10}, \emph{40\%}).
\item \textbf{Impersonnalisation} : \emph{Il est constaté que...}, \emph{La visite a révélé que...}, \emph{Le constat montre que...} (remplace \emph{Je pense que...}).
\end{itemize}
\end{propriete}

\courssub{Gestion des tonalités : neutralité obligatoire}

\begin{definition}[Tonalités interdites au rapport]
\begin{itemize}
\item \textbf{Tonalité dramatique (pathétique)} : exclamations, hyperboles, vocabulaire affectif (\emph{catastrophe, lamentable, scandaleux, horrible}), appel à la pitié (\emph{pauvres élèves}).
\item \textbf{Tonalité comique (ironique, satirique)} : sous-entendus, moqueries, exagération volontaire, familiarité.
\end{itemize}
Le rapport impose la \textbf{tonalité neutre (objective, factuelle, clinique)}.
\end{definition}

\begin{methode}[Garantir l'objectivité du rapport (Remédiation stylistique)]
\begin{enumerate}
\item \textbf{Repérer les marques de subjectivité} : \emph{je}, adjectifs qualitatifs non mesurables, ponctuation expressive (!, ...), modalisateurs (\emph{malheureusement, heureusement}).
\item \textbf{Substituer le fait au jugement} : \emph{« L'atelier est sale »} $\rightarrow$ \emph{« Des déchets métalliques et des huiles usagées jonchent le sol de l'atelier »}.
\item \textbf{Quantifier} : remplacer \emph{beaucoup, peu, énormément} par des chiffres, pourcentages, proportions (\emph{40\%, 6 sur 10, 3/4}).
\item \textbf{Utiliser la voix passive ou l'impersonnel} : \emph{« On a vu »} $\rightarrow$ \emph{« Il a été observé »} / \emph{« Le rapporteur a constaté »}.
\end{enumerate}
\end{methode}

\begin{exoresolu}[Neutraliser un passage subjectif]
\textbf{Énoncé.} Réécris ce passage de rapport pour le rendre conforme au genre : \emph{C'est une vraie catastrophe ! L'atelier est dans un état lamentable, je n'ai jamais rien vu d'aussi scandaleux. Ces pauvres élèves ne peuvent vraiment rien apprendre là-dedans.}
\tcblower
\textbf{Solution détaillée.}
\begin{enumerate}
\item \textbf{Repérage des fautes de registre} : exclamation (\emph{C'est une vraie catastrophe !}) = tonalité dramatique ; adjectifs subjectifs (\emph{lamentable, scandaleux}) ; jugement personnel (\emph{je n'ai jamais rien vu}) ; pathos (\emph{ces pauvres élèves}).
\item \textbf{Remplacement par des faits observables} : on décrit ce qu'on voit et on quantifie.
\item \textbf{Version corrigée} : \emph{Lors de sa visite du 12 mars, le rapporteur a constaté que six machines sur dix sont hors d'usage, que les établis sont endommagés et que l'éclairage est insuffisant (moins de 300 lux). Dans ces conditions, les travaux pratiques ne peuvent se dérouler normalement.}
\item \textbf{Vérification} : plus de \emph{je}, plus d'exclamation, chaque affirmation repose sur un constat chiffré ou observable ; le ton est neutre.
\end{enumerate}
\textbf{Conclusion.} L'objectivité du rapport repose sur trois réflexes : supprimer le \emph{je} émotif, remplacer les jugements par des faits, quantifier les constats.
\end{exoresolu}

\coursec{Élaboration du plan détaillé : Méthodologie pas à pas}

\courssub{Principes de la structure interne}

\begin{propriete}[Architecture standard du développement]
Le développement d'un rapport technique suit une progression logique \textbf{inductive} (du fait vers la solution) en trois temps :
\begin{enumerate}
\item \textbf{Constats (État des lieux)} : \emph{Qu'ai-je vu ? Quels sont les faits bruts, mesurables ?}
\item \textbf{Analyse (Causes \& Effets)} : \emph{Pourquoi ces faits ? Quelles en sont les conséquences sur la formation / la production ?}
\item \textbf{Propositions (Recommandations)} : \emph{Que faire ? Par qui ? Avec quels moyens ? Dans quel délai ?}
\end{enumerate}
\end{propriete}

\begin{methode}[Passer des notes de terrain au plan détaillé]
\begin{enumerate}
\item \textbf{Classer les informations} recueillies en trois familles : faits constatés, causes, conséquences / solutions.
\item \textbf{Rédiger l'introduction} : ordre de mission, dates et lieu de la mission, annonce explicite du plan (ex. \emph{« Nous présenterons d'abord les constats, puis nous en analyserons les causes avant de formuler des recommandations »}).
\item \textbf{Construire le développement en parties numérotées} (I, II, III), chaque partie portant un \textbf{titre-énoncé} (phrase nominale ou verbale complète) et se divisant en sous-parties (a, b, c).
\item \textbf{Prévoir la conclusion} : bilan synthétique (réponse à la commande) puis suggestions ou recommandations concrètes, numérotées, adressées au destinataire.
\item \textbf{Vérifier la hiérarchie} : une idée directrice = une partie ; chaque sous-partie s'appuie sur un fait, un chiffre, un témoignage ou un document iconographique référencé.
\end{enumerate}
\begin{attention}[Réflexe « Plan détaillé »]
Un plan détaillé n'est pas une liste de mots-clés : chaque titre de partie est une phrase ou un groupe nominal qui annonce clairement le contenu. On n'écrit jamais « Partie 1 : constats » sans préciser \emph{quels} constats. De même, on ne rédige pas de paragraphes : on s'arrête au niveau des titres et sous-titres.
\end{attention}
\end{methode}

\begin{exoresolu}[Construire un plan détaillé]
\textbf{Énoncé.} À la suite d'une mission d'inspection de la bibliothèque de ton lycée, tu as noté : rayonnages vides ; 40\% des ouvrages datent d'avant 1990 ; absence de salle de lecture ; les élèves empruntent peu ; manque de crédits ; un seul documentaliste pour 1\,200 élèves. Propose le plan détaillé du rapport à adresser au Proviseur.
\tcblower
\textbf{Solution détaillée.}
\begin{enumerate}
\item \textbf{Classement des notes} :
\begin{itemize}
\item Faits constatés : rayonnages vides, ouvrages anciens (40\% $<$ 1990), pas de salle de lecture, un seul documentaliste pour 1\,200 élèves.
\item Cause : manque de crédits.
\item Conséquence : faible fréquentation, peu d'emprunts.
\end{itemize}
\item \textbf{Introduction} : rappel de l'ordre de mission du Proviseur (référence, date), date de la visite, annonce du plan en trois parties (Constats, Analyse, Propositions).
\item \textbf{Développement} :
\begin{itemize}
\item \emph{I. Un constat alarmant : la pauvreté du fonds documentaire et de l'accueil.}
\begin{itemize}
\item a) Des rayonnages insuffisamment garnis (taux de remplissage estimé à 30\%).
\item b) Un fonds documentaire vieilli : 40\% des ouvrages antérieurs à 1990.
\item c) L'absence totale de salle de lecture pour les élèves.
\end{itemize}
\item \emph{II. Des causes structurelles et des effets sur la fréquentation.}
\begin{itemize}
\item a) Une cause budgétaire principale : l'insuffisance chronique des crédits alloués à la documentation.
\item b) Une cause organisationnelle : un personnel réduit (un documentaliste pour 1\,200 élèves).
\item c) Une conséquence directe mesurable : la faible fréquentation et le taux d'emprunt quasi nul.
\end{itemize}
\item \emph{III. Des recommandations concrètes pour redynamiser la bibliothèque.}
\begin{itemize}
\item a) L'inscription d'une ligne budgétaire dédiée à l'achat d'ouvrages récents (programme 2025-2026).
\item b) L'aménagement d'un espace de lecture dans le local adjacent inoccupé.
\item c) La demande de recrutement ou de mise à disposition d'un second documentaliste.
\end{itemize}
\end{itemize}
\item \textbf{Conclusion} : Bilan (la bibliothèque, dans l'état, ne remplit pas sa mission pédagogique de soutien aux apprentissages) et appel à une décision rapide du Proviseur pour le prochain conseil d'administration.
\end{enumerate}
\textbf{Conclusion.} Le plan est hiérarchisé (trois parties, trois sous-parties chacune), chaque note de terrain trouve sa place, et la progression va du constat vers les solutions : c'est la logique attendue d'un rapport.
\end{exoresolu}

\coursec{Amélioration du plan et du texte : Révision, cohérence et iconographie}

\courssub{Révision du plan (Compte rendu et remédiation)}

\begin{methode}[Réviser un plan détaillé avant rédaction finale]
\begin{enumerate}
\item \textbf{Contrôler la cohérence logique} : l'ordre est-il Constats $\rightarrow$ Causes/Analyse $\rightarrow$ Solutions ? (Pas de causes avant constats).
\item \textbf{Éliminer les doublons (redondances)} : deux sous-parties ne doivent pas recouvrir le même fait.
\item \textbf{Combler les vides (exhaustivité)} : chaque affirmation du plan doit pouvoir s'appuyer sur une preuve (chiffre, citation, photo, document) : « Preuve ? » à côté de chaque sous-titre.
\item \textbf{Harmoniser la formulation des titres} : même structure syntaxique pour les titres de même niveau (ex. toutes les parties commencent par un nom + complément : \emph{Un constat...}, \emph{Des causes...}, \emph{Des recommandations...}).
\item \textbf{Rédiger un compte rendu de remédiation} : noter sur une fiche annexe les défauts initiaux identifiés, les corrections apportées et les raisons du choix final. Cette trace écrit l'apprentissage.
\end{enumerate}
\end{methode}

\begin{exoresolu}[Corriger un plan défectueux]
\textbf{Énoncé.} Voici le plan proposé par un élève : \emph{I. Causes du problème. II. Constat : les machines sont en panne. III. Les machines sont en panne. IV. Suggestions.} Releve trois défauts et propose un plan corrigé.
\tcblower
\textbf{Solution détaillée.}
\begin{enumerate}
\item \textbf{Défaut 1 — Ordre illogique} : les causes (I) sont placées avant le constat (II) ; un rapport expose d'abord les faits, puis les explique.
\item \textbf{Défaut 2 — Doublon (redondance)} : II et III répètent exactement la même idée (\emph{les machines sont en panne}).
\item \textbf{Défaut 3 — Titres non harmonisés et imprécis} : \emph{Suggestions} ne précise pas lesquelles ; les titres n'ont pas la même construction grammaticale (phrase nominale vs phrase verbale).
\item \textbf{Plan corrigé} :
\begin{itemize}
\item \emph{I. Un constat : la majorité des machines de l'atelier est hors d'usage.}
\begin{itemize}
\item a) Six machines sur dix en panne (taux d'indisponibilité 60\%).
\item b) Des pièces de rechange introuvables sur le marché local.
\end{itemize}
\item \emph{II. Des causes techniques et budgétaires identifiées.}
\begin{itemize}
\item a) L'absence de plan d'entretien préventif régulier.
\item b) L'insuffisance du budget de maintenance alloué au département.
\end{itemize}
\item \emph{III. Des recommandations opérationnelles à court et moyen terme.}
\begin{itemize}
\item a) L'établissement d'un calendrier d'entretien préventif mensuel.
\item b) Le lancement d'une commande groupée de pièces détachées critiques.
\item c) La proposition d'un contrat de maintenance externe pour les équipements de haute technologie.
\end{itemize}
\end{itemize}
\end{enumerate}
\textbf{Conclusion.} La révision du plan obéit à trois contrôles : ordre logique, absence de doublon, titres précis et harmonisés. Ce travail de remédiation, consigné dans un compte rendu, fait partie intégrante de la production d'écrits.
\end{exoresolu}

\courssub{Intégration de l'iconographie et communication numérique}

\begin{methode}[Exploiter un document iconographique et transmettre par voie numérique]
\begin{enumerate}
\item \textbf{Décrire avant d'interpréter} : nature du document (photographie, schéma, tableau statistique), localisation, date, ce qui est visible.
\item \textbf{Légender et numéroter} impérativement : \emph{Figure 1 : État du tour n°3 — Usure anormale du chariot (photographie du 12/03/202X)}.
\item \textbf{Ancrer l'image dans le texte} : renvoyer explicitement à la figure dans le développement (\emph{« Comme le montre la Figure 1, le chariot présente... »}).
\item \textbf{Respecter les codes du courriel professionnel (transmission)} :
\begin{itemize}
\item Objet clair et complet : \emph{Transmission - Rapport [Objet] - [Nom Prénom]}.
\item Formules de politesse administratives complètes (pas de « Cordialement » seul, pas d'abréviations SMS : \emph{svp, pj, bjr}).
\item Annonce des pièces jointes dans le corps du message avec leurs noms exacts.
\item Signature complète (Nom, Prénom, Fonction, Contact).
\end{itemize}
\item \textbf{Vérifier la fiabilité d'une source numérique citée} : auteur identifiable, date de mise à jour, extension institutionnelle (.gouv.cm, .edu, .org) ou site reconnu.
\end{enumerate}
\end{methode}

\begin{exoresolu}[Légender un document et rédiger un courriel de transmission]
\textbf{Énoncé.} Tu joins à ton rapport trois photographies de l'atelier et tu transmets le tout au Principal par courriel. 1) Rédige la légende de la première photographie. 2) Rédige le courriel d'envoi (objet, salutations, corps bref, formule finale).
\tcblower
\textbf{Solution détaillée.}
\begin{enumerate}
\item \textbf{Légende} : elle décrit, localise et date le document : \emph{Figure 1 : Rangée de tours parallèles hors d'usage dans l'atelier de mécanique (photographie prise le 12 mars 202X, vue depuis l'entrée principale).}
\item \textbf{Objet du courriel} : \emph{Transmission du rapport sur l'état du matériel de l'atelier de mécanique — N. [Nom Élève], Délégué 1ère F4.}
\item \textbf{Corps} : \emph{Monsieur le Principal, Conformément à votre note de service n° 12/PRIN/LTB du 10 mars, veuillez trouver ci-joint le rapport sur l'état du matériel de l'atelier de mécanique, accompagné de ses trois annexes photographiques (Figures 1 à 3). Je reste à votre entière disposition pour tout complément d'information.}
\item \textbf{Formule finale} : \emph{Je vous prie d'agréer, Monsieur le Principal, l'expression de ma respectueuse considération.}
\item \textbf{Pièces jointes annoncées} : \emph{Rapport\_Atelier\_Mecanique\_1ereF4.pdf ; Fig1\_Tours\_HS.jpg ; Fig2\_Etabli\_Endommage.jpg ; Fig3\_Eclairage\_Defaillant.jpg.}
\end{enumerate}
\textbf{Conclusion.} Un document iconographique n'a de valeur probante dans un rapport que s'il est numéroté, légendé, daté et explicitement relié au texte par un renvoi ; le courriel de transmission respecte les mêmes exigences de forme et de politesse que le courrier papier.
\end{exoresolu}

\coursec{Situation d'intégration : Du terrain au rapport final}

\begin{experience}[Mission complète : Simulation « Inspection Atelier Soudure »]
\textbf{Contexte.} Tu es élève en Première Technicien Soudeur. Le Chef d'Atelier (destinataire) te mandate par note n° 045/CA/LTB du 05/10 pour inspecter le poste de soudure TIG n°7 et le stock de consommables (baguettes, gaz argon). Tu dois rendre un rapport complet (structure externe + plan détaillé + introduction rédigée + courriel d'envoi) sous 48h.

\textbf{Données de terrain (Notes brutes) :}
\begin{itemize}
\item Poste TIG n°7 : arc instable, pédale de commande HS, torche qui fuit (gaz).
\item Stock baguettes Inox 316L : 2 kg restants (seuil alerte 5 kg). Stock Argon : 1 bouteille pleine + 1 vide (pas de détendeur sur la pleine).
\item Conséquence : Arrêt fabrication pièce « Châssis remorque » (commande client urgente).
\item Cause identifiée : Pas de maintenance préventive mensuelle (dernière il y a 14 mois). Commande consommables bloquée « procédure achat » depuis 3 semaines.
\item Suggestions : Réparation pédale/torche (devis fournisseur local 48h). Commande urgence baguettes (fournisseur Yaoundé J+2). Débloquer procédure achat / stock sécurité. Planning maintenance mensuel.
\end{itemize}

\textbf{Travail demandé (Évaluation sommative) :}
\begin{enumerate}
\item \textbf{Structure externe} : Rédige l'en-tête complet, l'objet, la formule d'ouverture, la formule de clôture + signature.
\item \textbf{Plan détaillé} : Élabore le plan détaillé (I, II, III + a, b, c) avec titres-énoncés précis. Indique en marge de chaque sous-partie la \textbf{preuve} (chiffre, fait, doc) qui la soutient.
\item \textbf{Introduction rédigée} : Rédige le premier paragraphe du corps du rapport (contexte, mission, annonce du plan).
\item \textbf{Courriel de transmission} : Objet, corps, formules, noms des 3 pièces jointes (Rapport + 2 photos).
\item \textbf{Auto-évaluation (Remédiation)} : Coche la grille ci-dessous avant de rendre ta copie.
\end{enumerate}

\begin{center}
\begin{tabularx}{\linewidth}{|l|X|c|}\hline
\textbf{Critère de réussite (Barème Probatoire)} & \textbf{Indicateurs observables} & \textbf{Validé ?} \\ \hline
Structure externe complète (6 mentions) & En-tête, Lieu/Date, Référence, Objet, Destinataire, Formules (Ouverture/Clôture), Signature/Fonction & $\square$ \\ \hline
Plan détaillé : Logique & Ordre Constats $\rightarrow$ Causes $\rightarrow$ Solutions respecté & $\square$ \\ \hline
Plan détaillé : Précision & Titres-énoncés complets (pas de mots seuls), 3 parties, 3 sous-parties min. & $\square$ \\ \hline
Plan détaillé : Preuves & Chaque sous-partie a sa preuve identifiée (chiffre, date, doc) & $\square$ \\ \hline
Introduction rédigée & Contexte + Mission + Annonce plan + Ton neutre & $\square$ \\ \hline
Langue : Phrase complexe & Au moins 3 phrases complexes (cause/conséquence) correctement ponctuées & $\square$ \\ \hline
Langue : Emprunts / Vocab. & 0 emprunt direct non justifié (ex. \emph{mail $\rightarrow$ courriel}, \emph{stock $\rightarrow$ stock} ok terme tech) & $\square$ \\ \hline
Iconographie / Numérique & Légendes Figures (n° + description + date) + Courriel pro (Objet, PJ, Politesse) & $\square$ \\ \hline
Objectivité / Tonalité & 0 \emph{je} subjectif, 0 adjectif affectif, 0 point d'exclamation, quantifié & $\square$ \\ \hline
\end{tabularx}
\end{center}
\legende{Grille d'auto-évaluation / remédiation pour la situation d'intégration (Leçon 14).}
\end{experience}

\begin{attention}[Les 5 erreurs qui coûtent des points au Probatoire / Bac Technique]
\begin{enumerate}
\item \textbf{Oublier les mentions externes} : pas d'objet, pas de date, pas de destinataire ou pas de signature : le rapport est amputé d'une partie significative du barème. Vérifie toujours la liste : en-tête, lieu/date, référence, objet, destinataire, formules, signature.
\item \textbf{Confondre plan détaillé et rédaction complète} : quand le sujet demande le \emph{plan détaillé}, on présente des titres de parties et de sous-parties hiérarchisés (I, a, b...), pas des paragraphes rédigés.
\item \textbf{Glisser de la subjectivité} : \emph{je trouve que}, adjectifs émotionnels (\emph{lamentable}), exclamations (tonalité dramatique) ou ironie (tonalité comique) détruisent l'objectivité attendue ; écris \emph{le rapporteur a constaté} et quantifie.
\item \textbf{Mal analyser la phrase complexe} : compter les propositions sans compter les verbes conjugués, ou confondre coordination et subordination ; repère toujours le mot de liaison (\emph{parce que} = subordonnée de cause, \emph{donc} = coordination).
\item \textbf{Employer des emprunts bruts ou abréviations SMS} dans un document officiel (\emph{software, mail, « svp », « tt », « pj »}) sans équivalent français ni correction : préfère \emph{logiciel, courriel, s'il vous plaît, pièce jointe} et une orthographe complète.
\end{enumerate}
\end{attention}

\begin{savaistu}[Le rapport technique dans la vie professionnelle camerounaise]
Au Cameroun, la maîtrise du rapport technique est une compétence transversale exigée dans tous les secteurs du technique : BTP (rapport de chantier journalier/hebdomadaire), Industrie (rapport de maintenance, rapport d'incident), Tertiaire (rapport de stage, rapport d'activité, procès-verbal de réception). L'administration (MINESEC, MINEFOP, MINESUP) évalue cette compétence au Probatoire et au Baccalauréat Technique à travers des sujets de « Production d'écrits utilitaires » qui reproduisent fidèlement les situations vécues en entreprise (ordre de mission, note de service, visite de terrain). La rigueur de la structure externe (hiérarchie, politesse administrative) conditionne souvent la recevabilité du document avant même la lecture du fond.
\end{savaistu}

\newpage

\coursec{Exercices}

\courssub{Je vérifie mes connaissances}

\begin{exercice}[QCM : Structure et finalités du rapport technique]
Parmi les propositions suivantes, une seule correspond à la structure externe canonique d'un rapport technique formel adressé à une hiérarchie. Laquelle ?
\begin{enumerate}
    \item Page de garde, Table des matières, Introduction, Développement (2 à 3 parties), Conclusion, Annexes, Bibliographie.
    \item Page de garde, Remerciements, Introduction, Corps du texte (une seule partie), Conclusion, Signature.
    \item En-tête (Expéditeur/Destinateur/Objet/Date), Introduction, Développement, Conclusion, Pièces jointes.
    \item Page de garde, Sommaire, Développement, Recommandations, Annexes, Index.
\end{enumerate}
\end{exercice}

\begin{exercice}[Vrai / Faux Justifié : La phrase complexe et l'emprunt]
Pour chaque affirmation, cochez Vrai ou Faux et justifiez brièvement votre réponse en mobilisant les notions de subordination et d'emprunt lexical.
\begin{enumerate}
    \item « Bien que la pompe soit neuve, elle fuit » est une phrase simple contenant une proposition infinitive.
    \item L'emprunt lexical « le software » (au lieu de « le logiciel ») est un anglicisme non adapté qu'il faut proscrire dans un rapport technique officiel.
    \item Dans la phrase « Le technicien a réparé la machine afin que la production reprenne », la proposition subordonnée exprime le but et est introduite par une locution conjonctive.
    \item L'emprunt intégral « le marketing » est morphologiquement invariant en français (on écrit « des marketing » au pluriel).
\end{enumerate}
\end{exercice}

\begin{exercice}[Définitions et identification : Tonalités et textuel explicatif]
\begin{enumerate}
    \item Définissez le \cle{texte explicatif} en précisant sa finalité principale et deux de ses marques linguistiques caractéristiques (temps, connecteurs, vocabulaire).
    \item Différenciez la \cle{tonalité dramatique} de la \cle{tonalité comique} en donnant pour chacune l'effet recherché sur le lecteur.
    \item Dans un rapport de stage, quelle tonalité est la plus appropriée ? Justifiez en une phrase.
\end{enumerate}
\end{exercice}

\begin{exercice}[Analyse d'un plan détaillé : Cohérence et articulation]
Voici l'ébauche d'un plan détaillé pour un rapport d'activité mensuel d'un atelier de maintenance. Identifiez \textbf{deux incohérences majeures} (logique, parallélisme, exhaustivité) et proposez une correction pour chacune.

\textbf{Plan proposé :}
\begin{itemize}
    \item \textbf{Introduction}
    \item \textbf{Partie I : Les pannes survenues en janvier}
        \begin{itemize}
            \item A. Pannes électriques
            \item B. Pannes mécaniques
            \item C. Les interventions du personnel
        \end{itemize}
    \item \textbf{Partie II : Gestion du stock de pièces de rechange}
        \begin{itemize}
            \item A. Entrées
            \item B. Sorties
        \end{itemize}
    \item \textbf{Partie III : Le personnel est compétent}
        \begin{itemize}
            \item A. Formation suivie
            \item B. Heures supplémentaires effectuées
        \end{itemize}
    \item \textbf{Conclusion}
\end{itemize}
\end{exercice}

\courssub{Je m'entraîne}

\begin{exercice}[Transformation syntaxique : Notes de terrain $\rightarrow$ Phrases complexes techniques]
Transformez chaque groupe de phrases simples (issues de notes de terrain) en \textbf{une seule phrase complexe} correcte, au registre technique (passif, vocabulaire précis, subordination circonstancielle cause/conséquence/but). Respectez l'objectivité.

\begin{enumerate}
    \item Le débit a chuté. La vanne était obstruée. (Cause)
    \item On a remplacé le roulement. La vibration a cessé. (Conséquence)
    \item Le technicien calibre le capteur. Il veut garantir la précision. (But)
    \item La température dépasse 90°C. Le système se met en sécurité. (Conséquence)
    \item On a nettoyé le filtre. Le moteur surchauffait. (Cause)
\end{enumerate}
\end{exercice}

\begin{exercice}[Élaboration d'un plan détaillé : Rapport d'activité mensuel (Contexte CAMWATER)]
Vous êtes technicien à la Direction Exploitation de CAMWATER (Cameroun). À partir des données brutes ci-dessous, élaborez le \textbf{plan détaillé} (Introduction, 2 Parties avec 2 sous-parties chacune, Conclusion) d'un rapport d'activité mensuel (Mois de Mars) adressé au Directeur Régional. Le plan doit intégrer l'analyse d'un texte iconique (courbe de débit/pression).

\textbf{Données brutes (Mars) :}
\begin{itemize}
    \item Production totale : 1 250 000 m³ (objectif : 1 300 000 m³).
    \item Taux de rendement réseau : 68\% (fuites : 32\%).
    \item 12 pannes majeures sur conduite DN300 (quartier Mvog-Ada).
    \item Remplacement de 150 compteurs défectueux.
    \item Courbe de débit/pression (Pièce jointe 1) : montre une corrélation inverse entre pics de consommation (6h-9h / 18h-21h) et chute de pression < 1,5 bar en bout de réseau.
    \item Budget maintenance : 92\% consommé.
\end{itemize}
\end{exercice}

\begin{exercice}[Intégration d'iconographie : Description et analyse d'un graphique]
Le graphique ci-dessous (simulé) représente l'évolution du taux de remplissage du réservoir R3 (capacité 500 m³) sur 24h.

\begin{center}
\begin{tikzpicture}[scale=0.7]
\begin{axis}[
    width=\linewidth, height=8cm,
    xlabel={Heure (h)}, ylabel={Niveau (m³)},
    xmin=0, xmax=24, ymin=0, ymax=550,
    xtick={0,6,12,18,24}, ytick={0,125,250,375,500},
    grid=major, grid style={dashed, gray!50},
    axis lines=middle,
    x axis line style={-Stealth}, y axis line style={-Stealth},
]
\addplot[popblue, thick, domain=0:6, samples=50] {500 - 50*x};
\addplot[popblue, thick, domain=6:12, samples=50] {200 + 50*(x-6)};
\addplot[popblue, thick, domain=12:18, samples=50] {500 - 50*(x-12)};
\addplot[popblue, thick, domain=18:24, samples=50] {200 + 50*(x-18)};
\draw[dashed, popred, thick] (0,125) -- (24,125) node[right] {Seuil alerte (25\%)};
\end{axis}
\end{tikzpicture}
\end{center}

Rédigez le \textbf{paragraphe descriptif et analytique} (style explicatif, objectivité, vocabulaire technique) qui figurera dans le corps du rapport pour commenter ce graphique. Votre paragraphe doit comporter au moins une phrase complexe de cause et une de conséquence.
\end{exercice}

\begin{exercice}[Révision et amélioration : Cohérence interne et externe]
Voici l'introduction d'un rapport de fin de stage. Elle présente des défauts de structure (rappel des 4 mouvements : Sujet amené, Sujet posé, Sujet divisé, Annonce du plan) et de registre. Réécrivez-la entièrement en corrigeant les défauts identifiés.

\textbf{Introduction à corriger :}
« Mon stage s'est passé à la SONARA à Limbé. C'était super intéressant. J'ai vu comment on raffine le pétrole. Il y a beaucoup d'unités. Mon rapport va parler de l'unité de distillation. D'abord je vais dire ce que j'ai fait. Ensuite les problèmes. Enfin les solutions. »
\end{exercice}

\courssub{Je me prépare au Bac}

\begin{exercice}[Sujet type Probatoire : Plan détaillé à partir d'un dossier documentaire (45 min)]
\textbf{Consigne :} À partir du dossier documentaire ci-dessous, élaborez le \textbf{plan détaillé} (Introduction, Parties, Sous-parties, Conclusion) d'un \textbf{rapport d'activité mensuel} (Mois d'Octobre) que vous adressez au Directeur Général de l'entreprise \textbf{AGRO-INDUSTRIE CAMEROUN (AIC)}. Vous êtes le Chef de Service Maintenance. Le plan doit faire apparaître explicitement l'exploitation des 3 documents.

\textbf{Document 1 : Organigramme simplifié AIC (Extrait)}
\begin{center}
\begin{tikzpicture}[node distance=1.2cm and 2cm, every node/.style={draw, rounded corners, fill=popblueL, thick}]
\node (DG) {Directeur Général};
\node (DTech) [below left of=DG] {Directeur Technique};
\node (DProd) [below right of=DG] {Directeur Production};
\node (SMaint) [below of=DTech, align=center]{Service Maintenance \\ (Vous)};
\node (SQual) [below of=DProd] {Service Qualité};
\draw[->, thick] (DG) -- (DTech);
\draw[->, thick] (DG) -- (DProd);
\draw[->, thick] (DTech) -- (SMaint);
\draw[->, thick] (DProd) -- (SQual);
\end{tikzpicture}
\end{center}

\textbf{Document 2 : Tableau de production et maintenance - Octobre}
\begin{center}
\begin{tabularx}{\linewidth}{|X|c|c|c|}\hline
\rowcolor{popblueL} \textbf{Unité} & \textbf{Production (tonnes)} & \textbf{Arrêts imprévus (h)} & \textbf{Interventions Maintenance} \\ \hline
Broyage & 12 500 & 14 & 8 (dont 3 curatives) \\ \hline
Pressage & 11 800 & 22 & 12 (dont 7 curatives) \\ \hline
Raffinage & 10 200 & 5 & 4 (préventives) \\ \hline
Conditionnement & 10 000 & 2 & 2 (préventives) \\ \hline
\end{tabularx}
\end{center}

\textbf{Document 3 : Photo d'un atelier (Description textuelle)}
« Vue de l'atelier de pressage : présence de fuites d'huile hydraulique sur la presse n°3 (flaque au sol), garde-corps manquant sur la passerelle d'accès au moteur principal, outillage rangé de manière désordonnée près de la zone de circulation des chariots élévateurs. »

\textbf{Attendus du plan :}
\begin{enumerate}
    \item Introduction complète (4 mouvements).
    \item Partie I : Bilan de l'activité production/maintenance (exploitation Doc 1 \& 2).
    \item Partie II : Sécurité, conformité et prévention (exploitation Doc 3 \& 2).
    \item Conclusion (Bilan synthétique + Recommandations priorisées).
\end{enumerate}
\end{exercice}

\begin{exercice}[Sujet type Baccalauréat Technique : Rapport de fin de stage CAMWATER (Épreuve écrite - 3h)]
\textbf{Sujet :} « Vous avez effectué un stage de 3 mois (Juillet -- Septembre) à la CAMWATER, Direction Régionale du Centre, Service Exploitation -- Secteur Distribution. Rédigez le \textbf{rapport de fin de stage} adressé au Directeur Régional. »

\textbf{Travail demandé (Étape 1 notée sur 6 pts) :} Produisez le \textbf{plan détaillé complet} de ce rapport. Ce plan doit obligatoirement intégrer l'analyse d'un \textbf{texte iconique} (courbe de débit/pression du réseau de distribution du secteur Mfoundi, fournie en annexe du sujet original -- ici décrite textuellement).

\textbf{Description de la courbe (Annexe) :} Courbe journalière type. Abscisses : Heures (0h--24h). Ordonnées : Pression (bars) et Débit (m³/h). Deux courbes superposées. Pic de débit à 7h (120 m³/h) et 19h (135 m³/h). Pression chute à 1,2 bar à ces heures (norme min 1,5 bar). Pression stable 3,5 bars la nuit (fuites probables).

\textbf{Structure imposée du rapport :}
\begin{enumerate}
    \item Page de garde (normes).
    \item Remerciements / Dédicace.
    \item Introduction générale (4 mouvements + annonce du plan en 3 parties).
    \item \textbf{Partie I : Présentation de la structure d'accueil et du service d'affectation.}
    \item \textbf{Partie II : Activités techniques réalisées et analyse de la performance du réseau (Intégration de la courbe ici).}
    \item \textbf{Partie III : Difficultés rencontrées, solutions apportées et propositions d'amélioration.}
    \item Conclusion générale (Bilan personnel + Professionnel + Perspectives).
    \item Annexes (dont la courbe commentée).
\end{enumerate}
\textit{Détaillez les sous-parties (A, B, 1, 2) pour les Parties I, II, III en cohérence avec le stage « Exploitation/Distribution ».}
\end{exercice}

\begin{exercice}[Sujet type Bac Tech : Production du texte (Suite de l'exercice 2 - 2h30)]
\textbf{Consigne :} À partir du plan détaillé que vous avez élaboré à l'exercice précédent (ou du plan type ci-dessous si le vôtre n'est pas validé), rédigez le \textbf{texte intégral} des éléments suivants :
\begin{enumerate}
    \item L'\textbf{Introduction générale} (environ 300 mots).
    \item Le \textbf{développement complet de la Partie II} (Activités techniques et analyse de performance), en intégrant obligatoirement le commentaire de la courbe de débit/pression (texte iconique) dans la sous-partie B.
    \item La \textbf{Conclusion générale} (environ 200 mots).
\end{enumerate}

\textbf{Plan type de secours pour la Partie II (si plan personnel non validé) :}
\begin{itemize}
    \item \textbf{Partie II : Activités techniques et analyse de la performance du réseau de distribution}
        \begin{itemize}
            \item A. Opérations d'exploitation courante et interventions sur le terrain
                \begin{itemize}
                    \item 1. Manœuvres de vannes, secteurs et gestion des abonnés
                    \item 2. Détection et réparation des fuites visibles (Bilan chiffré)
                \end{itemize}
            \item B. Analyse hydraulique du secteur Mfoundi : exploitation de la courbe débit/pression
                \begin{itemize}
                    \item 1. Description de la courbe journalière type (Débit/Pression)
                    \item 2. Diagnostic : Déficit de pression aux pointes / Surpression nocturne
                    \item 3. Conséquences : Insatisfaction usagers / Perte en charge / Risque casse
                \end{itemize}
            \item C. Actions correctives engagées et suivi des indicateurs de performance
        \end{itemize}
\end{itemize}

\textbf{Contraintes de rédaction :}
\begin{itemize}
    \item Registre technique soutenu (passif, impersonnel, vocabulaire précis : \textit{vanne de sectionnement, colonne montante, perte de charge, rendement réseau, abonné, compteur}).
    \item Phrases complexes (subordination cause/conséquence/but/condition).
    \item Objectivité, clarté, concision.
    \item Intégration fluide de l'iconographie : « La figure 1 met en évidence... », « On observe sur la courbe que... ».
\end{itemize}
\end{exercice}

\newpage

\coursec{Corrigés des exercices}

\coursec{Leçon 14 : Production d'écrits utilitaires — Le rapport (élaboration du plan détaillé)}

\courssub{Corrigés des exercices d'application et de synthèse}

\begin{corrige}[Exercice 1 — QCM : Structure et finalités du rapport technique]
\textbf{Bonne réponse : 1.}

\textbf{Justification :} La structure externe canonique d'un rapport technique formel (rapport de stage, rapport d'activité adressé à une hiérarchie) comprend : une \cle{page de garde} (titre, nom, structure, date), une \cle{table des matières}, une \cle{introduction}, un \cle{développement} articulé en 2 à 3 parties, une \cle{conclusion}, puis les \cle{annexes} et la \cle{bibliographie} éventuelle.

\textbf{Pourquoi les autres propositions sont fausses :}
\begin{itemize}
    \item Proposition 2 : un développement en « une seule partie » est insuffisant ; un rapport exige une division du sujet en plusieurs parties (au moins deux). Les remerciements ne remplacent pas la table des matières.
    \item Proposition 3 : il s'agit de la structure d'une \emph{lettre administrative} ou d'une note de service (en-tête Expéditeur/Destinataire/Objet/Date), non d'un rapport formel.
    \item Proposition 4 : l'ordre est incorrect (le développement ne précède pas l'introduction, qui manque d'ailleurs) ; un index est facultatif et les recommandations se placent dans la conclusion, non comme élément externe autonome.
\end{itemize}
\end{corrige}

\begin{corrige}[Exercice 2 — Vrai / Faux Justifié : La phrase complexe et l'emprunt]
\begin{enumerate}
    \item \textbf{Faux.} « Bien que la pompe soit neuve, elle fuit » est une phrase \cle{complexe} : elle contient deux propositions — une subordonnée circonstancielle de concession (« Bien que la pompe soit neuve », introduite par la locution conjonctive \emph{bien que} + subjonctif) et une principale (« elle fuit »). Il n'y a aucune proposition infinitive.
    \item \textbf{Vrai.} « Software » est un \cle{emprunt} à l'anglais non adapté (anglicisme). Le français dispose d'un équivalent créé et officiel : \emph{logiciel}. Dans un rapport technique officiel, on proscrit les anglicismes lorsqu'un terme français normalisé existe, afin de garantir clarté et correction du registre.
    \item \textbf{Vrai.} « Afin que la production reprenne » est une proposition subordonnée circonstancielle de \cle{but} (finale), introduite par la locution conjonctive \emph{afin que}, suivie du subjonctif (« reprenne »). Elle dépend de la principale « Le technicien a réparé la machine ».
    \item \textbf{Faux.} « Marketing » est un emprunt intégral, mais il suit en français la règle générale du pluriel : on écrit « des \emph{marketings} » (comme « des matchings »). Dire qu'il est invariant est une erreur ; en pratique on préférera toutefois « la mercatique » ou « les techniques de commercialisation » dans un rapport soutenu.
\end{enumerate}
\end{corrige}

\begin{corrige}[Exercice 3 — Définitions et identification : Tonalités et texte explicatif]
\begin{enumerate}
    \item \textbf{Le texte explicatif} est un texte dont la finalité principale est de \emph{faire comprendre} un phénomène, un fonctionnement ou une notion, en répondant aux questions \emph{comment ?} et \emph{pourquoi ?}. Marques linguistiques caractéristiques (deux attendues) :
    \begin{itemize}
        \item emploi du \textbf{présent de vérité générale} (présent intemporel) et de tournures impersonnelles (« on observe que », « il apparaît que ») ;
        \item connecteurs logiques de cause, conséquence et but : \emph{parce que, car, en effet, c'est pourquoi, par conséquent, afin de} ;
        \item vocabulaire précis et dénoté, définitions et reformulations (\emph{c'est-à-dire, autrement dit}).
    \end{itemize}
    \item \textbf{Tonalité dramatique} : elle vise à susciter l'\emph{émotion}, la tension, la tristesse ou l'angoisse chez le lecteur (registre pathétique : gradation, hyperbole, lexique de la souffrance). \textbf{Tonalité comique} : elle vise à provoquer le \emph{rire} ou le sourire (jeu de mots, ironie, exagération burlesque, quiproquo).
    \item Dans un rapport de stage, la tonalité appropriée est la tonalité \textbf{neutre / objective} (didactique) : le rapport est un texte utilitaire qui doit informer sans chercher à émouvoir ni à amuser, d'où l'emploi de tournures impersonnelles et du vocabulaire technique dénoté.
\end{enumerate}
\end{corrige}

\begin{corrige}[Exercice 4 — Analyse d'un plan détaillé : Cohérence et articulation]
\textbf{Incohérence 1 : défaut de parallélisme dans la Partie I.} Les sous-parties A et B classent les pannes par \emph{nature} (électriques / mécaniques), tandis que C (« Les interventions du personnel ») relève d'un autre critère (l'action humaine). Le classement n'est pas homogène.
\textbf{Correction :} soit on complète le classement par nature : « C. Pannes hydrauliques » ; soit on déplace « Les interventions du personnel » dans une partie autonome (Partie III consacrée aux ressources humaines).

\textbf{Incohérence 2 : titre de la Partie III non objectif et non parallèle.} « Le personnel est compétent » est un jugement de valeur subjectif, formulé en phrase, alors que les autres titres sont des groupes nominaux descriptifs (« Les pannes... », « Gestion du stock... »). Un rapport exige des titres neutres et parallèles.
\textbf{Correction :} « Partie III : La gestion des ressources humaines » avec « A. Formations suivies » et « B. Heures supplémentaires effectuées ».

\textbf{Remarque complémentaire (exhaustivité) :} la Partie II ne comporte que deux sous-parties (Entrées / Sorties) ; pour un bilan de stock complet, on ajouterait « C. État du stock final et alertes de réapprovisionnement ».
\end{corrige}

\begin{corrige}[Exercice 5 — Transformation syntaxique : Notes de terrain $\rightarrow$ Phrases complexes]
\begin{enumerate}
    \item \textbf{Cause :} « Le débit a chuté \emph{parce que} la vanne était obstruée. » (ou au passif : « Une chute de débit a été constatée \emph{en raison de} l'obstruction de la vanne. »)
    \item \textbf{Conséquence :} « Le roulement a été remplacé, \emph{si bien que} la vibration a cessé. » (ou : « \emph{De sorte que} la vibration a cessé. »)
    \item \textbf{But :} « Le technicien calibre le capteur \emph{afin de} garantir la précision des mesures. » (ou : « \emph{afin que} la précision des mesures soit garantie. » — subjonctif)
    \item \textbf{Conséquence :} « La température dépasse 90~\textdegree C, \emph{par conséquent} le système se met en sécurité. » (ou : « \emph{Si} la température dépasse 90~\textdegree C, le système se met en sécurité. » — condition)
    \item \textbf{Cause :} « Le filtre a été nettoyé \emph{parce que} le moteur surchauffait. » (ou : « \emph{Comme} le moteur surchauffait, le filtre a été nettoyé. »)
\end{enumerate}

\textbf{Points de vigilance :} employer le passif ou la tournure impersonnelle pour l'objectivité ; vérifier le mode après la locution conjonctive (subjonctif après \emph{afin que}, indicatif après \emph{parce que} et \emph{si bien que}) ; un seul verbe conjugué par proposition.
\end{corrige}

\begin{corrige}[Exercice 6 — Élaboration d'un plan détaillé : Rapport d'activité CAMWATER (Mars)]
\textbf{Plan détaillé attendu (modèle) :}

\textbf{Introduction :} Sujet amené (mission de la CAMWATER : production et distribution d'eau potable) ; Sujet posé (bilan d'activité du mois de mars, Direction Exploitation) ; Sujet divisé (performance de production et rendement réseau d'une part, interventions et moyens engagés d'autre part) ; Annonce du plan (deux parties).

\textbf{Partie I : Bilan de la production et performance hydraulique du réseau}
\begin{itemize}
    \item \textbf{A. Production mensuelle et écart par rapport à l'objectif} : 1\,250\,000 m\textsuperscript{3} produits pour 1\,300\,000 m\textsuperscript{3} visés, soit un déficit de 50\,000 m\textsuperscript{3} ($-3,8\,\%$).
    \item \textbf{B. Analyse du rendement réseau et exploitation de la courbe débit/pression (Pièce jointe 1)} : rendement de 68\,\% (fuites : 32\,\%) ; corrélation inverse entre pics de consommation (6h--9h / 18h--21h) et chute de pression sous 1,5 bar en bout de réseau.
\end{itemize}

\textbf{Partie II : Interventions techniques et moyens engagés}
\begin{itemize}
    \item \textbf{A. Maintenance curative et renouvellement d'équipements} : 12 pannes majeures sur conduite DN300 (quartier Mvog-Ada) ; remplacement de 150 compteurs défectueux.
    \item \textbf{B. Suivi budgétaire et perspectives} : budget maintenance consommé à 92\,\% ; besoin de renforcement (programme de recherche de fuites, calendrier de renouvellement de la conduite DN300).
\end{itemize}

\textbf{Conclusion :} bilan synthétique (objectif presque atteint, rendement réseau insuffisant) ; recommandations priorisées (réduction des fuites, renforcement de pression en bout de réseau, arbitrage budgétaire).

\textbf{Points de vigilance :} titres en groupes nominaux parallèles ; chaque donnée brute doit être exploitée une et une seule fois ; la courbe (texte iconique) doit être explicitement localisée dans le plan.
\end{corrige}

\begin{corrige}[Exercice 7 — Intégration d'iconographie : Description et analyse d'un graphique]
\textbf{Paragraphe type attendu :}

« La figure 1 présente l'évolution du niveau de remplissage du réservoir R3 (capacité : 500 m\textsuperscript{3}) sur une journée de 24 heures. On observe une alternance régulière de deux phases : une phase de vidange, au cours de laquelle le niveau passe de 500 m\textsuperscript{3} à 200 m\textsuperscript{3} (de 0h à 6h, puis de 12h à 18h), et une phase de remplissage, qui ramène le niveau de 200 m\textsuperscript{3} à 500 m\textsuperscript{3} (de 6h à 12h, puis de 18h à 24h). La vidange s'explique \emph{parce que} la demande des abonnés est maximale durant les pointes de consommation, \emph{de sorte que} le débit sortant dépasse le débit de production. \emph{Comme} le niveau plancher de 200 m\textsuperscript{3} demeure supérieur au seuil d'alerte fixé à 125 m\textsuperscript{3} (25\,\% de la capacité), la sécurité d'alimentation est assurée ; \emph{par conséquent}, aucune mesure corrective urgente n'est requise, mais un suivi quotidien de la courbe est recommandé \emph{afin de} prévenir tout franchissement du seuil critique. »

\textbf{Barème indicatif (sur 5 pts) :} description fidèle des valeurs lues sur le graphique (1,5 pt) ; vocabulaire technique et objectivité (tournures impersonnelles) (1 pt) ; au moins une subordonnée de cause (0,75 pt) et une de conséquence (0,75 pt) ; phrase d'amorce intégrant l'iconographie (« La figure 1 présente... ») (0,5 pt) ; correction linguistique (0,5 pt).

\textbf{Points de vigilance :} ne jamais écrire « on voit sur le dessin » ; citer les valeurs chiffrées exactes ; distinguer description (constats) et analyse (interprétation).
\end{corrige}

\begin{corrige}[Exercice 8 — Révision et amélioration : Cohérence interne et externe]
\textbf{Défauts identifiés :} registre familier (« super intéressant », « j'ai vu comment on raffine ») ; sujet amené absent (pas de contexte général) ; sujet posé confus ; annonce du plan maladroite et au futur de narration personnelle (« je vais dire ») ; absence de tournures impersonnelles.

\textbf{Introduction réécrite (modèle) :}

« Le raffinage du pétrole constitue un maillon essentiel de l'industrie énergétique camerounaise : il transforme le brut importé en produits finis utilisables (carburants, gaz, bitume). C'est dans ce cadre que s'est déroulé, au sein de la Société Nationale de Raffinage (SONARA) de Limbé, un stage de découverte des procédés industriels, organisé en plusieurs unités de traitement. Le présent rapport porte plus particulièrement sur l'unité de distillation, cœur du processus de séparation des hydrocarbures. Il s'articulera en trois parties : après avoir présenté les activités réalisées au sein de l'unité, on analysera les difficultés techniques rencontrées, avant de proposer des solutions d'amélioration. »

\textbf{Points de vigilance :} respecter les quatre mouvements dans l'ordre (sujet amené $\rightarrow$ sujet posé $\rightarrow$ sujet divisé $\rightarrow$ annonce du plan) ; bannir la première personne et les appréciations subjectives ; annoncer le plan avec des tournures impersonnelles (« il s'articulera en... », « on analysera... »).
\end{corrige}

\begin{corrige}[Exercice 9 — Sujet type Probatoire : Plan détaillé à partir d'un dossier documentaire]
\textbf{Plan détaillé attendu (modèle) :}

\textbf{Introduction :} Sujet amené (AIC, entreprise agro-industrielle ; rôle stratégique de la maintenance dans la continuité de production) ; Sujet posé (bilan mensuel d'octobre du Service Maintenance, placé sous l'autorité du Directeur Technique — Doc 1) ; Sujet divisé (performance production/maintenance ; sécurité et prévention) ; Annonce du plan (deux parties).

\textbf{Partie I : Bilan de l'activité production/maintenance (Docs 1 et 2)}
\begin{itemize}
    \item \textbf{A. Cadre organisationnel et production globale} : positionnement du Service Maintenance dans l'organigramme (rattachement à la Direction Technique — Doc 1) ; production totale d'octobre : $12\,500 + 11\,800 + 10\,200 + 10\,000 = 44\,500$ tonnes.
    \item \textbf{B. Analyse des arrêts et des interventions par unité (Doc 2)} : Pressage, unité la plus touchée (22 h d'arrêts, 12 interventions dont 7 curatives) ; Broyage (14 h, 8 interventions) ; Raffinage et Conditionnement bien maîtrisés (5 h et 2 h, interventions préventives uniquement) ; constat : la maintenance reste trop curative (10 curatives sur 26 interventions).
\end{itemize}

\textbf{Partie II : Sécurité, conformité et prévention (Docs 3 et 2)}
\begin{itemize}
    \item \textbf{A. Non-conformités constatées à l'atelier de pressage (Doc 3)} : fuite d'huile hydraulique sur la presse n\textsuperscript{o}3 (risque de glissade et de pollution) ; garde-corps manquant sur la passerelle (risque de chute de hauteur) ; rangement défectueux de l'outillage en zone de circulation des chariots (risque de collision).
    \item \textbf{B. Mesures correctives et plan de prévention} : réparation immédiate de la presse n\textsuperscript{o}3 et pose du garde-corps ; mise en conformité 5S de l'atelier ; bascule progressive vers la maintenance préventive (référence aux bons résultats du Raffinage — Doc 2).
\end{itemize}

\textbf{Conclusion :} bilan synthétique (production satisfaisante mais pressage vulnérable) ; recommandations priorisées : 1. sécuriser la passerelle (urgence) ; 2. réparer la presse n\textsuperscript{o}3 ; 3. renforcer la prévention au Pressage.

\textbf{Barème indicatif (sur 10 pts) :} introduction en 4 mouvements (2 pts) ; exploitation explicite des 3 documents, localisée dans le plan (3 pts) ; parallélisme et cohérence des titres (2 pts) ; hiérarchisation Parties/sous-parties (2 pts) ; présentation (1 pt).

\textbf{Points de vigilance :} vérifier le calcul de la production totale (44\,500 tonnes) ; ne pas rédiger le rapport, seul le \emph{plan} est demandé ; citer les documents entre parenthèses pour prouver leur exploitation.
\end{corrige}

\begin{corrige}[Exercice 10 — Sujet type Baccalauréat Technique : Plan détaillé du rapport de stage CAMWATER]
\textbf{Plan détaillé complet attendu (modèle) :}

\textbf{Éléments préliminaires :} Page de garde (ministère, établissement, titre, nom, année) ; Remerciements ; Dédicace ; Sommaire.

\textbf{Introduction générale :} Sujet amené (l'eau potable, enjeu vital ; mission de la CAMWATER) ; Sujet posé (stage de 3 mois, juillet--septembre, Direction Régionale du Centre, Service Exploitation -- Secteur Distribution) ; Sujet divisé (structure d'accueil ; activités et performance du réseau ; difficultés et perspectives) ; Annonce du plan en trois parties.

\textbf{Partie I : Présentation de la structure d'accueil et du service d'affectation}
\begin{itemize}
    \item \textbf{A. La CAMWATER et la Direction Régionale du Centre}
        \begin{itemize}
            \item 1. Historique, missions et organisation générale de la CAMWATER
            \item 2. Organisation et attributions de la Direction Régionale du Centre
        \end{itemize}
    \item \textbf{B. Le Service Exploitation -- Secteur Distribution}
        \begin{itemize}
            \item 1. Attributions du service et organigramme interne
            \item 2. Moyens matériels et humains du secteur distribution (réseau Mfoundi)
        \end{itemize}
\end{itemize}

\textbf{Partie II : Activités techniques réalisées et analyse de la performance du réseau}
\begin{itemize}
    \item \textbf{A. Opérations d'exploitation courante}
        \begin{itemize}
            \item 1. Manœuvres de vannes, gestion des secteurs et suivi des abonnés
            \item 2. Détection et réparation des fuites visibles (bilan chiffré)
        \end{itemize}
    \item \textbf{B. Analyse hydraulique du secteur Mfoundi : exploitation de la courbe débit/pression (Annexe)}
        \begin{itemize}
            \item 1. Description de la courbe journalière type (pics de débit à 7h : 120 m\textsuperscript{3}/h et 19h : 135 m\textsuperscript{3}/h)
            \item 2. Diagnostic : chute de pression à 1,2 bar aux pointes (norme $\geq$ 1,5 bar) ; surpression nocturne stable à 3,5 bars (fuites probables)
            \item 3. Conséquences : insatisfaction des usagers, pertes en charge, risque de casse des conduites
        \end{itemize}
    \item \textbf{C. Actions correctives engagées et suivi des indicateurs de performance}
\end{itemize}

\textbf{Partie III : Difficultés rencontrées, solutions apportées et propositions d'amélioration}
\begin{itemize}
    \item \textbf{A. Difficultés rencontrées}
        \begin{itemize}
            \item 1. Difficultés matérielles (vétusté des conduites, pénurie de pièces)
            \item 2. Difficultés organisationnelles et humaines
        \end{itemize}
    \item \textbf{B. Solutions apportées et propositions}
        \begin{itemize}
            \item 1. Solutions mises en œuvre pendant le stage
            \item 2. Propositions d'amélioration (recherche systématique de fuites nocturnes, régulation de pression)
        \end{itemize}
\end{itemize}

\textbf{Conclusion générale :} bilan personnel (compétences acquises) et professionnel (contribution au service) ; perspectives (poursuite du programme de réduction des fuites).

\textbf{Annexes :} courbe débit/pression commentée ; fiches d'intervention ; organigramme.

\textbf{Barème indicatif (sur 6 pts) :} introduction en 4 mouvements avec annonce en 3 parties (1 pt) ; structure externe complète (page de garde à annexes) (1 pt) ; sous-parties détaillées et parallèles pour les 3 parties (2 pts) ; intégration explicite et localisée de la courbe (Partie II.B) avec valeurs exactes (1,5 pt) ; cohérence avec le secteur « Exploitation/Distribution » (0,5 pt).

\textbf{Points de vigilance :} reprendre les valeurs exactes de l'annexe (120 et 135 m\textsuperscript{3}/h ; 1,2 bar ; 3,5 bars ; norme 1,5 bar) ; titres en groupes nominaux ; aucune sous-partie isolée (une division comporte au moins deux éléments).
\end{corrige}

\begin{corrige}[Exercice 11 — Sujet type Bac Tech : Production du texte (rédaction)]
\textbf{1. Introduction générale (modèle, environ 300 mots) :}

« L'accès à l'eau potable constitue un enjeu majeur de santé publique et de développement. Au Cameroun, la Camerounaise des Eaux (CAMWATER) assure la production et la distribution de l'eau potable en milieu urbain et périurbain ; la qualité de ce service repose sur l'exploitation quotidienne des réseaux de distribution, dont la performance hydraulique conditionne directement la satisfaction des abonnés. C'est dans ce cadre qu'un stage de trois mois, de juillet à septembre, a été effectué à la Direction Régionale du Centre, au sein du Service Exploitation -- Secteur Distribution, en charge du réseau du secteur Mfoundi. Le présent rapport a pour objet de rendre compte de cette période d'immersion professionnelle : il présente la structure d'accueil et le service d'affectation, analyse les activités techniques réalisées ainsi que la performance hydraulique du réseau à partir de la courbe de débit/pression du secteur, puis expose les difficultés rencontrées, les solutions apportées et les propositions d'amélioration. »

\textbf{2. Développement de la Partie II (extrait rédigé, modèle) :}

« \textbf{A. Opérations d'exploitation courante.} Durant la période de stage, les manœuvres de vannes de sectionnement ont été assurées quotidiennement afin d'équilibrer l'alimentation des secteurs ; les réclamations des abonnés ont été enregistrées puis traitées par ordre de priorité. Par ailleurs, les fuites visibles ont été systématiquement localisées et réparées : vingt-trois fuites sur branchements et sept sur conduites secondaires ont ainsi été colmatées, ce qui a permis de limiter les pertes d'eau potable.

\textbf{B. Analyse hydraulique du secteur Mfoundi.} La figure 1 (annexe) met en évidence l'évolution conjointe du débit et de la pression sur une journée type. On observe sur la courbe deux pics de consommation, à 7h (120 m\textsuperscript{3}/h) et à 19h (135 m\textsuperscript{3}/h). \emph{Parce que} la demande excède alors la capacité de refoulement, la pression chute à 1,2 bar, \emph{si bien qu'}elle passe sous la norme minimale de 1,5 bar ; \emph{par conséquent}, les abonnés situés en point haut ou en bout de réseau sont mal desservis aux heures de pointe. À l'inverse, la pression demeure stable et élevée la nuit (3,5 bars) alors que la consommation est quasi nulle : cette surpression nocturne traduit vraisemblablement l'existence de fuites non détectées et accroît le risque de casse des conduites. \emph{Afin de} remédier à ce déficit, un programme de recherche de fuites nocturnes et l'installation de régulateurs de pression ont été proposés.

\textbf{C. Actions correctives et suivi.} Des secteurs ont été rééquilibrés par manœuvres de vannes, et un tableau de suivi des indicateurs (taux de rendement, nombre de fuites réparées, pression aux points critiques) a été mis en place afin que la performance du réseau soit mesurée de manière continue. »

\textbf{3. Conclusion générale (modèle, environ 200 mots) :}

« Au terme de ce stage de trois mois au Service Exploitation -- Secteur Distribution de la CAMWATER, le bilan est doublement positif. Sur le plan personnel, les techniques d'exploitation d'un réseau de distribution (manœuvres de vannes, détection de fuites, lecture et interprétation des courbes débit/pression) ont été maîtrisées, et la rigueur exigée par le service public de l'eau a été pleinement appréhendée. Sur le plan professionnel, l'analyse hydraulique du secteur Mfoundi a permis de diagnostiquer un déficit de pression aux heures de pointe et une surpression nocturne révélatrice de fuites, puis de formuler des propositions concrètes : recherche systématique de fuites, régulation de pression et suivi permanent des indicateurs. Il appartient désormais au service de pérenniser ces actions afin que le rendement du réseau et la satisfaction des abonnés s'améliorent durablement. »

\textbf{Barème indicatif (sur 14 pts) :} introduction 4 mouvements et registre soutenu (3 pts) ; Partie II conforme au plan, avec commentaire de courbe intégré et valeurs exactes (6 pts) ; conclusion bilans personnel et professionnel + perspectives (2 pts) ; phrases complexes variées (cause/conséquence/but) (1,5 pt) ; correction linguistique et présentation (1,5 pt).

\textbf{Points de vigilance :} employer le passif et l'impersonnel ; introduire l'iconographie par une phrase d'amorce (« La figure 1 met en évidence... ») ; ne pas recopier l'annexe sans analyse ; citer la norme de 1,5 bar pour justifier le diagnostic.
\end{corrige}

\newpage

\coursec{Fiche bilan}

\begin{popfigure}
\begin{tikzpicture}[
  mindmap,
  grow cyclic,
  every node/.style={concept, font=\small},
  root concept/.append style={concept color=popblue, font=\bfseries\small, minimum size=2.6cm},
  level 1 concept/.append style={level distance=3.4cm, sibling angle=60, font=\scriptsize},
  level 2 concept/.append style={level distance=2.2cm, font=\tiny},
]
\node[root concept]{Le rapport : plan détaillé}
  child[concept color=poporange]{ node {Genre et finalités}
    child { node {Informer, rendre compte} }
    child { node {Objectivité, clarté} }
  }
  child[concept color=popgreen]{ node {Structure externe}
    child { node {Page de garde, sommaire} }
    child { node {Annexes, références} }
  }
  child[concept color=poppurple]{ node {Structure interne}
    child { node {Intro : contexte, problème} }
    child { node {Développement : faits, analyses} }
    child { node {Conclusion : bilans, suggestions} }
  }
  child[concept color=poppink]{ node {Plan détaillé}
    child { node {Parties, sous-parties} }
    child { node {Idées + exemples} }
  }
  child[concept color=popgold]{ node {Outils de langue}
    child { node {Phrase complexe} }
    child { node {Emprunts lexicaux} }
  }
  child[concept color=popblue!60]{ node {Style et révision}
    child { node {Texte explicatif} }
    child { node {Cohérence, relecture} }
  };
\end{tikzpicture}
\legende{Carte mentale de la leçon : le rapport, de sa structure au plan détaillé.}
\end{popfigure}

\begin{bilan}
\textbf{Définitions clés}
\begin{itemize}
  \item \cle{Rapport} : texte fonctionnel et utilitaire qui rend compte de manière \emph{objective} et \emph{structurée} d'une mission, d'une enquête, d'un stage ou d'une activité, à destination d'un destinataire précis (chef de service, enseignant, autorité).
  \item \cle{Structure externe} : éléments matériels du rapport : page de garde, sommaire, pagination, annexes, signature.
  \item \cle{Structure interne} : organisation du contenu en trois temps : introduction (contexte, objet, problématique), développement (constats, faits, analyses), conclusion (bilan, suggestions).
  \item \cle{Plan détaillé} : squelette complet du rapport annonçant, pour chaque partie et sous-partie, l'idée principale, les arguments et les exemples.
  \item \cle{Phrase complexe} : phrase contenant au moins deux propositions (juxtaposition, coordination, subordination).
  \item \cle{Emprunt} : mot intégré au français à partir d'une langue étrangère (ex. \emph{week-end}, \emph{business}).
\end{itemize}

\textbf{Repères à connaître par cœur}
\begin{center}
\begin{tabularx}{\linewidth}{|l|X|}
\hline
\rowcolor{popblueL} \textbf{Notion} & \textbf{À retenir} \\
\hline
Introduction & Qui ? Quoi ? Où ? Quand ? Pourquoi ? Annonce du plan. \\
\hline
Développement & Faits observés, données chiffrées, analyses ; une idée par paragraphe. \\
\hline
Conclusion & Bilan synthétique + suggestions ou recommandations. \\
\hline
Phrase complexe & Subordonnées : relatives (qui, que, où), conjonctives (que, parce que, bien que), participiales. \\
\hline
Tonalités & Dramatique et comique : à connaître pour la réception, mais le rapport exige un ton \emph{neutre et objectif}. \\
\hline
\end{tabularx}
\end{center}

\textbf{Méthode express : construire le plan détaillé}
\begin{enumerate}
  \item Dégager le sujet, le destinataire et la finalité du rapport.
  \item Rassembler et classer les informations (fiches, notes de terrain).
  \item Regrouper les idées en 2 ou 3 grandes parties équilibrées.
  \item Pour chaque partie : sous-parties numérotées (I, A, 1), idée + exemple.
  \item Vérifier la logique (chronologique, thématique ou cause-effet) et la cohérence.
  \item Rédiger, puis relire : orthographe, connecteurs, objectivité.
\end{enumerate}
\end{bilan}

\begin{aretenir}[Checklist avant le Bac]
$\square$ Je sais définir le rapport et distinguer ses finalités (informer, rendre compte, proposer).\\
$\square$ Je connais les éléments de la structure externe (page de garde, sommaire, annexes).\\
$\square$ Je maîtrise la structure interne : introduction, développement, conclusion.\\
$\square$ Je sais rédiger une introduction qui présente le contexte et annonce le plan.\\
$\square$ Je sais construire un plan détaillé hiérarchisé (parties, sous-parties, idées, exemples).\\
$\square$ Je sais identifier et employer les propositions de la phrase complexe (relative, conjonctive).\\
$\square$ Je reconnais les emprunts lexicaux et je les utilise à bon escient.\\
$\square$ Je sais reconnaître un texte explicatif et ses marques (définitions, exemples, connecteurs).\\
$\square$ Je distingue les tonalités dramatique et comique, et je garde un ton neutre dans le rapport.\\
$\square$ Je sais lire et exploiter un texte iconique (schéma, tableau, image) dans un rapport.\\
$\square$ Je relis mon texte : cohérence, connecteurs logiques, orthographe, objectivité.\\
$\square$ Je respecte le temps imparti et je justifie ma démarche en conclusion.
\end{aretenir}

\findecours

\end{document}
